\chapter{Attention Models}\label{ch:attention}

Attention mechanisms constitute the natural complement to the geometric convolutions developed in the previous chapter: where convolutions impose rigid structure through shared kernels with content-independent weights, attention introduces adaptive flexibility, dynamically weighting the contributions of each element according to the information it carries. This chapter develops the mathematical foundations of attention mechanisms, from their classical algebraic formulation to the geometric extensions that allow operating over non-Euclidean domains while preserving symmetries.

We start from traditional attention mechanisms, formalized as integral operators over Hilbert spaces, and then progressively incorporate geometric structure: first equivariance with respect to symmetry groups, then operation over graphs via attention networks and spectral analysis, the extension to Riemannian manifolds with parallel transport in hyperbolic and spherical spaces, the formulation in Clifford algebras that unifies scalars, vectors, and multivectors within a coherent algebraic framework, and finally attention over simplicial complexes and cohomological structures. Each extension incorporates richer geometric structure while preserving the dynamic adaptation capability that distinguishes attention from convolution.

This progression from the algebraic to the geometric connects directly with the preceding chapters: \autoref{ch:cnn} established that attention generalizes convolution as a limiting case when the Gaussian kernel relaxes its local support, and \autoref{ch:gdl} identified \glspl{Transformer} as instances of the \gls{MPNN} framework over adaptive complete graphs (\autoref{prop:attention_mpnn}). This chapter develops both connections in depth.

\section{Traditional Attention Mechanisms}\label{sec:atencion_tradicional}

\subsection{Introduction and History}\label{subsec:historia_atencion}

Attention mechanisms emerge historically from the need to overcome the inherent limitations of sequential models in natural language processing. The early works of \citet{bahdanau2014neural} and \citet{luong2015effective} introduced attention mechanisms in machine translation systems based on recurrent neural networks, allowing the model to ``attend'' to different parts of the input sequence during generation.

The turning point came with the seminal work \textit{"Attention is All You Need"} by \citet{vaswani2017attentionneed}, which demonstrated that attention mechanisms alone, without recurrence or convolution, could achieve state-of-the-art performance on machine translation tasks. This \gls{Transformer} architecture revolutionized not only natural language processing, but became the foundation of computer vision models, graph processing, and, eventually, geometric deep learning.

The historical evolution can be summarized in three main phases:

\begin{enumerate}
    \item \textbf{Attention as a complement} (2014-2016): Attention mechanisms added to \glspl{RNN} to improve the handling of long-range dependencies
    \item \textbf{Attention as the main architecture} (2017-2019): Pure \glspl{Transformer} replacing recurrent architectures
    \item \textbf{Geometric attention} (2020-present): Extension to non-Euclidean data with preservation of symmetries
\end{enumerate}

\subsection{Algebraic Foundations}\label{subsec:fundamentos_algebraicos_atencion}

Attention mechanisms are grounded in concepts from linear algebra and the geometry of vector spaces. At their core, attention operates over three interrelated vector spaces that encode different aspects of the information.

\subsubsection{Feature Spaces}\label{subsubsec:espacios_caracteristicas}

Consider a feature space $\mathcal{H} = \mathbb{R}^d$ where the input representations reside. For a set of $n$ input elements, we have:

\begin{equation}
    \mathbf{X} = [\mathbf{x}_1, \mathbf{x}_2, \ldots, \mathbf{x}_n]^T \in \mathbb{R}^{n \times d}
    \label{eq:input_matrix}
\end{equation}

where each $\mathbf{x}_i \in \mathcal{H}$ represents an element of the input sequence.

\subsubsection{Fundamental Linear Transformations}

The attention mechanism defines three fundamental linear transformations:

\begin{align}
    \mathbf{Q} &= \mathbf{X}\mathbf{W}_Q \in \mathbb{R}^{n \times d_k} \quad \text{(Queries)} \\
    \mathbf{K} &= \mathbf{X}\mathbf{W}_K \in \mathbb{R}^{n \times d_k} \quad \text{(Keys)} \\
    \mathbf{V} &= \mathbf{X}\mathbf{W}_V \in \mathbb{R}^{n \times d_v} \quad \text{(Values)}
    \label{eq:linear_projections}
\end{align}

where the weight matrices $\mathbf{W}_Q, \mathbf{W}_K \in \mathbb{R}^{d \times d_k}$ and $\mathbf{W}_V \in \mathbb{R}^{d \times d_v}$ are trainable parameters that project the input space into specialized subspaces.

The dimensions are related as follows: if $\mathbf{X} \in \mathbb{R}^{n \times d}$ represents $n$ elements with $d$ features each, then the projections result in queries and keys of dimension $d_k$ and values of dimension $d_v$. In the standard multi-head implementation, $d_k = d_v = d/h$ is used, where $h$ is the number of attention heads.

\begin{definition}[Dot Product as a Similarity Measure]\label{def:dot_product_similarity}
    The fundamental compatibility function in attention is defined via the dot product:
    \begin{equation}
        \text{sim}(\mathbf{q}_i, \mathbf{k}_j) = \mathbf{q}_i^T \mathbf{k}_j = \|\mathbf{q}_i\| \|\mathbf{k}_j\| \cos(\theta_{ij})
        \label{eq:dot_product_similarity}
    \end{equation}
    where $\theta_{ij}$ is the angle between the vectors $\mathbf{q}_i$ and $\mathbf{k}_j$.
\end{definition}

This geometric interpretation reveals that attention measures the normalized cosine similarity between representations in the feature space.

\subsubsection{Geometry of the Q, K, V Spaces}

The underlying geometric structure of the query, key, and value spaces provides deep insights into how attention works.

\begin{definition}[Similarity Matrix in the Attention Space]\label{def:gram_metric_attention}
    The \textit{similarity matrix} (or cross inner-product matrix) associated with the queries and keys is defined as:
    \begin{equation}
        \mathbf{S} = \mathbf{Q}\mathbf{K}^T, \quad S_{ik} = \langle \mathbf{q}_i, \mathbf{k}_k \rangle \in \mathbb{R}^{n \times n}
        \label{eq:gram_matrix_attention}
    \end{equation}
    This matrix encodes all the geometric relations between queries and keys. When $\mathbf{Q} = \mathbf{K}$ (pure self-attention), $\mathbf{S}$ coincides with the Gram matrix of the keys.
\end{definition}

\begin{proposition}[Geometric Properties of the Similarity Matrix]\label{prop:geometric_properties_gram}
    The matrix $\mathbf{S} = \mathbf{Q}\mathbf{K}^T$ satisfies:
    \begin{enumerate}
        \item \textbf{Singular value decomposition}: $\mathbf{S} = \mathbf{U}\mathbf{\Sigma}\mathbf{V}^T$ (SVD), where $\mathbf{U}, \mathbf{V}$ capture the principal directions of queries and keys respectively and $\mathbf{\Sigma}$ contains the singular values.
        \item \textbf{Frobenius norm}: $\|\mathbf{S}\|_F^2 = \text{tr}(\mathbf{Q}\mathbf{K}^T\mathbf{K}\mathbf{Q}^T) = \|\mathbf{K}\mathbf{Q}^T\|_F^2$
        \item \textbf{Angular interpretation}: $S_{ij} = \|\mathbf{q}_i\|\|\mathbf{k}_j\|\cos(\theta_{ij})$, where $\theta_{ij}$ is the angle between $\mathbf{q}_i$ and $\mathbf{k}_j$.
    \end{enumerate}
\end{proposition}

\begin{proof}
    Property (1) is the standard singular value decomposition, applicable to every real matrix. Property (2) follows from $\|\mathbf{S}\|_F^2 = \text{tr}(\mathbf{S}\mathbf{S}^T) = \text{tr}(\mathbf{Q}\mathbf{K}^T(\mathbf{Q}\mathbf{K}^T)^T) = \text{tr}(\mathbf{Q}\mathbf{K}^T\mathbf{K}\mathbf{Q}^T)$, and by the cyclic property of the trace, $\text{tr}(\mathbf{Q}\mathbf{K}^T\mathbf{K}\mathbf{Q}^T) = \text{tr}(\mathbf{K}\mathbf{Q}^T\mathbf{Q}\mathbf{K}^T) = \|\mathbf{K}\mathbf{Q}^T\|_F^2$. Property (3) is an immediate consequence of the definition of the inner product in Euclidean spaces.
\end{proof}

\begin{theorem}[Geometric Projection Theorem in Attention]\label{theorem:geometric_projection_attention}
    The attention mechanism can be interpreted as a geometric projection of the value space $\mathcal{V}$ onto the subspace generated by the queries $\text{span}(\mathbf{Q})$, mediated by the metric induced by the keys.

    Formally, for each query $\mathbf{q}_i$:
    \begin{equation}
        \mathbf{y}_i = \sum_{j=1}^n \alpha_{ij} \mathbf{v}_j = \mathbf{V}^T \boldsymbol{\alpha}_i
        \label{eq:geometric_projection_formula}
    \end{equation}
    where $\boldsymbol{\alpha}_i = \text{softmax}(\mathbf{G}_{i,:}/\sqrt{d_k})$ defines a probability distribution over the value index.
\end{theorem}

\begin{proof}
    The proof proceeds by recognizing that the attention weights $\alpha_{ij}$ form a convex (probabilistic) combination over the values. This combination is determined by the geometric proximity between $\mathbf{q}_i$ and each key $\mathbf{k}_j$, measured by the scaled inner product and normalized by softmax.

    The resulting projection preserves the structure of the value space while emphasizing directions aligned with the query according to the metric defined by the keys.
\end{proof}

\begin{figure}[ht]
    \centering
    \includegraphics[width=0.95\textwidth]{figures/qkv_geometric_projection.pdf}
    \caption{Geometric interpretation of the attention mechanism. \textbf{Left}: in the space of queries and keys, the similarity between $\mathbf{q}_i$ and each $\mathbf{k}_j$ is measured by the angle $\theta_{ij}$ (scaled inner product). \textbf{Right}: the values $\mathbf{v}_j$ are weighted by the attention weights $\alpha_{ij}$, and the output $\mathbf{y}_i$ is their convex combination.}
    \label{fig:qkv_geometric_projection}
\end{figure}

\subsubsection{Multiresolution Analysis of Attention}

The analysis of attention from a harmonic analysis perspective reveals hierarchical structure at multiple scales.

\begin{definition}[Attention Spectrum]\label{def:attention_spectrum}
    For an attention matrix $\mathbf{A} = \text{softmax}(\mathbf{Q}\mathbf{K}^T/\sqrt{d_k})$, we define its \textit{spectrum} as the set of eigenvalues:
    \begin{equation}
        \text{Spec}(\mathbf{A}) = \{\lambda_1, \lambda_2, \ldots, \lambda_n : \mathbf{A}\mathbf{v}_i = \lambda_i \mathbf{v}_i\}
        \label{eq:attention_spectrum}
    \end{equation}
    The eigenvectors $\{\mathbf{v}_i\}$ reveal the principal attention patterns.
\end{definition}

\begin{example}[Explicit Computation: $3 \times 3$ Attention]\label{ex:explicit_3x3_attention}
    Consider a concrete example with $n=3$ elements and $d_k=2$:

    \begin{equation}
        \mathbf{Q} = \begin{bmatrix} 1 & 0 \\ 0 & 1 \\ 1 & 1 \end{bmatrix}, \quad
        \mathbf{K} = \begin{bmatrix} 1 & 0 \\ 1 & 1 \\ 0 & 1 \end{bmatrix}, \quad
        \mathbf{V} = \begin{bmatrix} 2 & 0 \\ 0 & 2 \\ 1 & 1 \end{bmatrix}
    \end{equation}

    \textbf{Step 1: Compatibility matrix}
    \begin{equation}
        \mathbf{S} = \mathbf{Q}\mathbf{K}^T = \begin{bmatrix} 1 & 1 & 0 \\ 0 & 1 & 1 \\ 1 & 2 & 1 \end{bmatrix}
    \end{equation}

    \textbf{Step 2: Attention matrix (without scaling, for simplicity)}
    \begin{equation}
        \mathbf{A} = \text{softmax}(\mathbf{S}) \approx \begin{bmatrix} 0.42 & 0.42 & 0.16 \\ 0.16 & 0.42 & 0.42 \\ 0.21 & 0.58 & 0.21 \end{bmatrix}
    \end{equation}

    \textbf{Step 3: Attention output}
    \begin{equation}
        \mathbf{Y} = \mathbf{A}\mathbf{V} \approx \begin{bmatrix} 1.00 & 1.00 \\ 0.74 & 1.26 \\ 0.63 & 1.37 \end{bmatrix}
    \end{equation}

    This example illustrates how different query-key patterns produce different convex combinations of values.
\end{example}

\subsection{Basic Attention Mechanism}\label{subsec:atencion_basico}

\subsubsection{Fundamental Attention Equation}

The scaled dot-product attention mechanism is formally defined as:

\begin{equation}
    \text{Attention}(\mathbf{Q}, \mathbf{K}, \mathbf{V}) = \text{softmax}\left(\frac{\mathbf{Q}\mathbf{K}^T}{\sqrt{d_k}}\right)\mathbf{V}
    \label{eq:scaled_dot_product_attention}
\end{equation}

The attention score matrix $\mathbf{S} = \mathbf{Q}\mathbf{K}^T \in \mathbb{R}^{n \times n}$ encodes the compatibility between each query-key pair, while the scaling factor $\frac{1}{\sqrt{d_k}}$ prevents the dot-product values from growing excessively with the dimension.

\subsubsection{Geometric Interpretation}

The attention operation can be interpreted as a convex interpolation in the value space:

\begin{equation}
    \mathbf{y}_i = \sum_{j=1}^{n} \alpha_{ij} \mathbf{v}_j, \quad \sum_{j=1}^{n} \alpha_{ij} = 1, \quad \alpha_{ij} \geq 0
    \label{eq:convex_combination}
\end{equation}

where the attention weights $\alpha_{ij}$ form a probability distribution over the values.

\begin{definition}[Generalized Attention Function]\label{def:generalized_attention}
    A \textit{generalized attention function} is a function $f: \mathcal{X}^n \times \mathcal{Y}^n \times \mathcal{Z}^n \to \mathcal{W}^n$ where:
    \begin{align}
        f(\mathbf{Q}, \mathbf{K}, \mathbf{V}) &= \sum_{j=1}^{n} \alpha_{ij} \mathbf{v}_j \\
        \alpha_{ij} &= \frac{\exp(e_{ij})}{\sum_{k=1}^{n} \exp(e_{ik})} \\
        e_{ij} &= a(\mathbf{q}_i, \mathbf{k}_j)
    \end{align}
    where $a: \mathcal{X} \times \mathcal{Y} \to \mathbb{R}$ is a compatibility function.
\end{definition}

\subsubsection{Self-Attention and Global Dependencies}

Self-attention constitutes the special case where the queries, keys, and values come from the same input:

\begin{equation}
    \mathbf{Q} = \mathbf{K} = \mathbf{V} = \mathbf{X}
    \label{eq:self_attention_case}
\end{equation}

This formulation allows each element of the sequence to attend to all the other elements, capturing global dependencies in a single operation.

\begin{theorem}[Universality of Self-Attention]\label{theorem:self_attention_universality}
    Let $K \subseteq \mathbb{R}^{n \times d}$ be a compact set. For any continuous function
    $f: K \to \mathbb{R}^{n \times d}$ that is permutation-equivariant (i.e., $f(P\mathbf{X}) = Pf(\mathbf{X})$
    for any permutation matrix $P$), there exists a self-attention network that can approximate
    $f$ uniformly on $K$ with arbitrary precision.

    This result extends the universal approximation theorems for symmetric functions
    \citep{zaheer2017deepsets} to the context of attention mechanisms. Full universality requires the complete Transformer architecture (self-attention combined with feed-forward layers and residual connections), not solely isolated self-attention layers; see the rigorous proof in \citet{yun2020transformers}.
\end{theorem}

\begin{proof}
    The proof proceeds in three main steps:

    \textbf{Step 1: Characterization of permutation-equivariant functions.}
    Let $f: K \to \mathbb{R}^{n \times d}$ be a continuous, permutation-equivariant function, where $K \subseteq \mathbb{R}^{n \times d}$ is compact. By the representation theorem for symmetric functions \citep{zaheer2017deepsets}, any permutation-equivariant function can be expressed as:

    \begin{equation}
        f(\mathbf{X})_i = \phi\left(\mathbf{x}_i, \sum_{j=1}^n \rho(\mathbf{x}_j)\right)
        \label{eq:symmetric_representation}
    \end{equation}

    where $\phi: \mathbb{R}^d \times \mathbb{R}^d \to \mathbb{R}^d$ and $\rho: \mathbb{R}^d \to \mathbb{R}^d$ are appropriate continuous functions.

    \textbf{Step 2: Approximation via pairwise interactions.}
    By the universal approximation theorem for neural networks, for any $\epsilon > 0$, there exist functions $g_1, g_2, g_3: \mathbb{R}^d \to \mathbb{R}^{d'}$ and $h: \mathbb{R}^{d'} \to \mathbb{R}^d$ such that:

    \begin{equation}
        \left\|f(\mathbf{X})_i - \sum_{j=1}^n \text{softmax}(g_1(\mathbf{x}_i)^T g_2(\mathbf{x}_j)) \cdot h(g_3(\mathbf{x}_j))\right\| < \epsilon
        \label{eq:pairwise_approximation}
    \end{equation}

    uniformly for $\mathbf{X} \in K$.

    \textbf{Step 3: Exact implementation via self-attention.}
    We define the projection matrices as:
    \begin{align}
        \mathbf{W}_Q &= \begin{bmatrix} \mathbf{W}_{g_1} \\ \mathbf{0} \end{bmatrix} \in \mathbb{R}^{d \times d_k} \\
        \mathbf{W}_K &= \begin{bmatrix} \mathbf{W}_{g_2} \\ \mathbf{0} \end{bmatrix} \in \mathbb{R}^{d \times d_k} \\
        \mathbf{W}_V &= \begin{bmatrix} \mathbf{W}_{g_3} \\ \mathbf{0} \end{bmatrix} \in \mathbb{R}^{d \times d_v}
    \end{align}

    where $\mathbf{W}_{g_1}, \mathbf{W}_{g_2}, \mathbf{W}_{g_3}$ implement the functions $g_1, g_2, g_3$ respectively via linear layers.

    With this construction, the self-attention mechanism produces exactly:
    \begin{equation}
        \text{SelfAttention}(\mathbf{X})_i = \sum_{j=1}^n \text{softmax}\left(\frac{g_1(\mathbf{x}_i)^T g_2(\mathbf{x}_j)}{\sqrt{d_k}}\right) h(g_3(\mathbf{x}_j))
        \label{eq:self_attention_implementation}
    \end{equation}

    The scaling factor $\sqrt{d_k}$ can be absorbed into the functions $g_1$ and $g_2$ without loss of generality.

    By composition of universally approximable functions and the density of continuous functions on compact spaces, we conclude that self-attention networks can uniformly approximate any continuous permutation-equivariant function on compact sets.
\end{proof}

\subsection{Multi-Head Attention}

\subsubsection{Decomposition into Parallel Subspaces}

The Multi-Head Attention mechanism extends basic attention through parallel processing in multiple representation subspaces:

\begin{equation}
    \text{MultiHead}(\mathbf{Q}, \mathbf{K}, \mathbf{V}) = \text{Concat}(\text{head}_1, \ldots, \text{head}_h)\mathbf{W}_O
    \label{eq:multihead_attention}
\end{equation}

where each attention head is defined as:

\begin{equation}
    \text{head}_i = \text{Attention}(\mathbf{Q}\mathbf{W}_Q^i, \mathbf{K}\mathbf{W}_K^i, \mathbf{V}\mathbf{W}_V^i)
    \label{eq:attention_head}
\end{equation}

with projection matrices $\mathbf{W}_Q^i, \mathbf{W}_K^i \in \mathbb{R}^{d \times d_k}$, $\mathbf{W}_V^i \in \mathbb{R}^{d \times d_v}$ for $i = 1, \ldots, h$.

\subsubsection{Advantages of the Decomposition}

The multi-head architecture provides several fundamental benefits:

\begin{enumerate}
    \item \textbf{Specialization of representations}: Each head can specialize in capturing different types of relations
    \item \textbf{Increased capacity}: The model can simultaneously attend to multiple positions and subspaces
    \item \textbf{Robustness}: The partial redundancy among heads improves training stability
\end{enumerate}

\begin{figure}[ht]
\centering
\includegraphics[width=0.85\textwidth]{figures/multi_head_attention.pdf}
\caption{Diagram of Multi-Head Attention. The input is projected into $h$ parallel subspaces, each with its own matrices $\mathbf{W}^Q_i$, $\mathbf{W}^K_i$, $\mathbf{W}^V_i$. Each head computes attention independently, and the results are concatenated and projected via $\mathbf{W}^O$.}
\label{fig:multi_head_attention}
\end{figure}

\begin{proposition}[Multi-Head Representational Capacity for Pre-Softmax Scores]\label{theorem:multihead_representational_capacity}
    The Multi-Head Attention mechanism with $h$ heads can represent any pre-softmax score pattern that admits a decomposition as a sum of $h$ low-rank functions.

    Formally, if $\mathbf{S}^* = \sum_{i=1}^h \mathbf{U}_i \mathbf{\Sigma}_i \mathbf{V}_i^T$ where each $\mathbf{U}_i \mathbf{\Sigma}_i \mathbf{V}_i^T$ has rank $r_i \leq d_k$, then there exists a Multi-Head configuration with $h$ heads whose pre-softmax score matrix is exactly $\mathbf{S}^*$.
\end{proposition}

\begin{proof}
    The proof proceeds by direct construction. For each term $\mathbf{U}_i \mathbf{\Sigma}_i \mathbf{V}_i^T$ in the decomposition, we define the $i$-th head as:
    \begin{align}
        \mathbf{W}_Q^i &= \mathbf{U}_i \mathbf{\Sigma}_i^{1/2} \\
        \mathbf{W}_K^i &= \mathbf{V}_i \mathbf{\Sigma}_i^{1/2} \\
        \mathbf{W}_V^i &= \mathbf{I}
    \end{align}

    With this choice, the pre-softmax score matrix of head $i$ is $\mathbf{Q}_i \mathbf{K}_i^T = \mathbf{X}\mathbf{U}_i \mathbf{\Sigma}_i \mathbf{V}_i^T \mathbf{X}^T$, and the sum over heads reproduces the desired structure $\mathbf{S}^*$.
\end{proof}

\begin{remark}
The above construction operates at the level of pre-softmax scores. In practice, each head applies softmax to its scores before aggregation with the values, which introduces the constraint that the effective attention weights are non-negative and sum to 1 per row. This nonlinearity prevents the multi-head mechanism from exactly representing arbitrary attention functions at the output level.
\end{remark}

\subsubsection{Orthogonality Analysis Between Heads}

\begin{definition}[Functional Orthogonality Between Heads]\label{def:functional_orthogonality_heads}
    Two attention heads $\text{head}_i$ and $\text{head}_j$ are \textit{functionally orthogonal} if their projection matrices satisfy:
    \begin{equation}
        \langle \mathbf{W}_Q^i, \mathbf{W}_Q^j \rangle_F = \text{tr}((\mathbf{W}_Q^i)^T \mathbf{W}_Q^j) = 0
        \label{eq:head_orthogonality}
    \end{equation}
    where $\langle \cdot, \cdot \rangle_F$ denotes the Frobenius inner product.
\end{definition}

\begin{proposition}[Benefits of Orthogonality]\label{prop:orthogonality_benefits}
    If the attention heads are mutually orthogonal in the sense of \autoref{def:functional_orthogonality_heads}, then:
    \begin{enumerate}
        \item The information captured by each head is complementary (non-redundant)
        \item The Frobenius norm of the Multi-Head is the sum of the individual norms: $\|\text{MultiHead}\|_F^2 = \sum_{i=1}^h \|\text{head}_i\|_F^2$
        \item The gradient during backpropagation decomposes cleanly across heads
    \end{enumerate}
\end{proposition}

\begin{proof}
\leavevmode
\begin{enumerate}
    \item \textbf{Complementarity.} If $\langle \mathbf{W}_Q^i, \mathbf{W}_Q^j \rangle_F = 0$ for $i \neq j$, the query subspaces are orthogonal. Hence each head projects onto a distinct subspace of $\mathbb{R}^d$, capturing linearly independent information.

    \item \textbf{Norm decomposition.} The Multi-Head output is $\text{MultiHead} = \sum_{i=1}^h \text{head}_i \mathbf{W}_O^i$. By the orthogonality of the heads and Parseval's identity for orthogonal subspaces:
    \[
    \|\text{MultiHead}\|_F^2 = \left\|\sum_{i=1}^h \text{head}_i \mathbf{W}_O^i\right\|_F^2 = \sum_{i=1}^h \|\text{head}_i \mathbf{W}_O^i\|_F^2 = \sum_{i=1}^h \|\text{head}_i\|_F^2
    \]
    where the last equality assumes $\mathbf{W}_O^i$ orthonormal.

    \item \textbf{Gradients.} By orthogonality, $\frac{\partial \mathcal{L}}{\partial \mathbf{W}_Q^i}$ has no components in the direction of the subspaces of $\text{head}_j$ ($j \neq i$), avoiding interference between the updates of different heads. \qedhere
\end{enumerate}
\end{proof}

\subsubsection{Complexity Analysis}

The computational complexity of Multi-Head Attention is:

\begin{align}
    \text{Time} &: O(n^2 d + ndhd_k) \\
    \text{Space} &: O(n^2 + ndh)
\end{align}

where the $n^2$ term dominates for long sequences, motivating subsequent optimizations such as sparse and linear attention.

\subsection{Formalization as an Integral Operator}\label{subsec:attention_integral_operator}

The self-attention mechanism can be rigorously formalized as an integral operator acting on functions defined on manifolds, providing a deep mathematical perspective that connects with the classical theory of operators.

\subsubsection{Continuous Formulation}

Let $\mathcal{M}$ be a compact Riemannian manifold equipped with a measure $\mu$. For a function $f: \mathcal{M} \to \mathbb{R}^d$, the attention operator is defined as:

\begin{equation}
    \mathcal{T}f(x) = \int_{\mathcal{M}} K(x, y) f(y) \, d\mu(y)
    \label{eq:attention_integral_operator}
\end{equation}

where the \textbf{attention kernel} $K: \mathcal{M} \times \mathcal{M} \to \mathbb{R}^{d \times d}$ satisfies:

\begin{equation}
    K(x, y) = \sigma\left(\frac{\langle Q(x), K(y) \rangle}{\sqrt{d_k}}\right) V(y)^T
    \label{eq:attention_kernel}
\end{equation}

with $Q, K, V: \mathcal{M} \to \mathbb{R}^{d_k}, \mathbb{R}^{d_k}, \mathbb{R}^{d_v}$ respectively, and $\sigma$ the generalized softmax function.

\subsubsection{Properties of the Operator}

\begin{proposition}[Compactness of the Attention Operator]\label{prop:attention_compactness}
    If the kernel $K(x,y)$ is continuous on $\mathcal{M} \times \mathcal{M}$ and $\mathcal{M}$ is compact, then the operator $\mathcal{T}: L^2(\mathcal{M}) \to L^2(\mathcal{M})$ is compact.
\end{proposition}

\begin{proof}
    We will show that $\mathcal{T}: L^2(\mathcal{M}) \to L^2(\mathcal{M})$ is compact by applying the Arzelà-Ascoli criterion. To this end, we will show that the image of the unit ball $B = \{f \in L^2(\mathcal{M}) : \|f\|_{L^2} \leq 1\}$ under $\mathcal{T}$ is relatively compact in $C(\mathcal{M})$ (and therefore in $L^2(\mathcal{M})$).

    \textbf{Step 1: Uniform boundedness.}
    For $f \in B$ and $x \in \mathcal{M}$, we have:
    \begin{align}
        \|\mathcal{T}f(x)\| &= \left\|\int_{\mathcal{M}} K(x, y) f(y) \, d\mu(y)\right\| \\
        &\leq \int_{\mathcal{M}} \|K(x, y)\| \|f(y)\| \, d\mu(y) \\
        &\leq \sup_{x,y \in \mathcal{M}} \|K(x, y)\| \cdot \|f\|_{L^2} \cdot \sqrt{\mu(\mathcal{M})}
    \end{align}

    Since $K$ is continuous on the compact set $\mathcal{M} \times \mathcal{M}$, the function $\|K(x,y)\|$ attains its maximum. Therefore, $\{\mathcal{T}f : f \in B\}$ is uniformly bounded.

    \textbf{Step 2: Equicontinuity.}
    Let $x_1, x_2 \in \mathcal{M}$ and $f \in B$. Then:
    \begin{align}
        \|\mathcal{T}f(x_1) - \mathcal{T}f(x_2)\| &= \left\|\int_{\mathcal{M}} [K(x_1, y) - K(x_2, y)] f(y) \, d\mu(y)\right\| \\
        &\leq \int_{\mathcal{M}} \|K(x_1, y) - K(x_2, y)\| \|f(y)\| \, d\mu(y) \\
        &\leq \|f\|_{L^2} \sqrt{\mu(\mathcal{M})} \sup_{y \in \mathcal{M}} \|K(x_1, y) - K(x_2, y)\|
    \end{align}

    Since $K$ is continuous on the compact $\mathcal{M} \times \mathcal{M}$, it is uniformly continuous. Therefore, for any $\epsilon > 0$, there exists $\delta > 0$ such that if $d_{\mathcal{M}}(x_1, x_2) < \delta$, then:
    \[\sup_{y \in \mathcal{M}} \|K(x_1, y) - K(x_2, y)\| < \frac{\epsilon}{\sqrt{\mu(\mathcal{M})}}\]

    This implies $\|\mathcal{T}f(x_1) - \mathcal{T}f(x_2)\| < \epsilon$ for $f \in B$, establishing equicontinuity.

    \textbf{Step 3: Application of the Arzelà-Ascoli theorem.}
    By the previous steps, $\{\mathcal{T}f : f \in B\}$ is an equicontinuous and uniformly bounded family in $C(\mathcal{M})$. Since $\mathcal{M}$ is compact, the Arzelà-Ascoli theorem guarantees that this family is relatively compact in $C(\mathcal{M})$.

    Since $C(\mathcal{M})$ embeds densely in $L^2(\mathcal{M})$ and uniform convergence implies convergence in $L^2$, we conclude that $\mathcal{T}$ is a compact operator.
\end{proof}

\subsubsection{Connection with Spectral Analysis}

The integral formulation allows analyzing attention through spectral theory, following the analytical framework developed by \citet{sander2023attention}. The \textbf{eigenfunctions} of the operator $\mathcal{T}$ capture the fundamental attention patterns:

\begin{equation}
    \mathcal{T}\phi_i = \lambda_i \phi_i, \quad \lambda_1 \geq \lambda_2 \geq \cdots \geq 0
    \label{eq:attention_eigendecomposition}
\end{equation}

The eigenvalues $\lambda_i$ indicate the relative importance of each attention pattern, while the eigenfunctions $\phi_i$ define the principal directions in the feature space.

\begin{figure}[ht]
    \centering
    \includegraphics[width=0.95\textwidth]{figures/attention_spectrum.pdf}
    \caption{Spectral decomposition of the attention operator. \textbf{Left}: the eigenvalues $\lambda_i$ of the attention kernel exhibit rapid decay ($\lambda_i \sim i^{-\alpha}$), which indicates a finite effective rank $r$ of the operator. The shaded region corresponds to the dominant modes. \textbf{Right}: the associated eigenfunctions $\phi_i$ show increasing frequency, analogously to the Fourier basis; the purple bars indicate the magnitude of the corresponding eigenvalue.}
    \label{fig:attention_spectrum}
\end{figure}

\subsubsection{Mercer Kernels and Advanced Spectral Theory}

The attention mechanism $\text{softmax}(QK^T/\sqrt{d})V$ can be interpreted as an integral operator with a data-dependent kernel. When this kernel satisfies the Mercer conditions (symmetry and positive definiteness), the powerful machinery of functional analysis (in particular the spectral theorem for compact operators) provides guarantees on the properties of the attention operator, such as its convergence and approximation capability.

\begin{definition}[Gaussian Attention Kernel]\label{def:gaussian_attention_kernel}
    A common variant of the attention kernel is the \textit{Gaussian kernel}:
    \begin{equation}
        K_{\text{gauss}}(x, y) = \exp\left(-\frac{\|Q(x) - K(y)\|^2}{2\sigma^2}\right)
        \label{eq:gaussian_attention_kernel}
    \end{equation}
    where $\sigma > 0$ is a scale parameter that controls the locality of the attention.
\end{definition}

\begin{theorem}[Spectral Decomposition of the Attention Kernel]\label{theorem:spectral_decomposition_attention}
    For an attention kernel $K: \mathcal{X} \times \mathcal{X} \to \mathbb{R}$ that is a Mercer kernel, there exists a spectral decomposition:
    \begin{equation}
        K(x, y) = \sum_{i=1}^{\infty} \lambda_i \phi_i(x) \phi_i(y)
        \label{eq:mercer_decomposition_attention}
    \end{equation}
    where $\{\lambda_i\}$ are non-negative eigenvalues in decreasing order and $\{\phi_i\}$ are orthonormal eigenfunctions.

    The corresponding attention operator admits the representation:
    \begin{equation}
        \mathcal{T}f = \sum_{i=1}^{\infty} \lambda_i \langle f, \phi_i \rangle_{L^2} \phi_i
        \label{eq:spectral_representation_attention}
    \end{equation}
\end{theorem}

\begin{proof}
    The proof follows from the spectral theorem for compact self-adjoint operators. Since $K$ is a Mercer kernel, the associated integral operator $\mathcal{T}$ is compact and positive. The spectral theorem guarantees the existence of the spectral decomposition, where the eigenvalues are real and non-negative.

    The convergence of the series is in the $L^2$ norm, and the spectral representation provides a complete characterization of the behavior of the operator.
\end{proof}

\begin{corollary}[Automatic Spectral Filtering]\label{cor:spectral_filtering}
    The attention operator acts as a natural low-pass filter, where:
    \begin{enumerate}
        \item High-frequency components (small eigenvalues) are attenuated
        \item The eigenvalue decay rate determines the smoothness of the attention function
        \item The effective rank of the operator is determined by the number of significant eigenvalues
    \end{enumerate}
\end{corollary}

\begin{proof}
Direct consequence of \autoref{theorem:spectral_decomposition_attention}. By the spectral representation $\mathcal{T}f = \sum_{i=1}^\infty \lambda_i \langle f, \phi_i \rangle \phi_i$ with eigenvalues in decreasing order $\lambda_1 \geq \lambda_2 \geq \cdots \geq 0$: (1)~the high-frequency components (large index $i$) are multiplied by small eigenvalues $\lambda_i$, attenuating them naturally; (2)~the decay rate $\lambda_i \to 0$ controls the smoothness of $\mathcal{T}f$ (rapid decay implies greater regularity); (3)~the effective rank is $r_\epsilon = |\{i : \lambda_i > \epsilon\}|$, which determines the intrinsic dimensionality of the representation.
\end{proof}

\subsubsection{Spectral Convergence Analysis}

\begin{theorem}[Convergence of Finite-Rank Approximations]\label{theorem:finite_rank_convergence}
    Let $\mathcal{T}_r$ be the rank-$r$ approximation of the attention operator:
    \begin{equation}
        \mathcal{T}_r f = \sum_{i=1}^{r} \lambda_i \langle f, \phi_i \rangle \phi_i
        \label{eq:finite_rank_approximation}
    \end{equation}

    Then, for any $f \in L^2(\mathcal{M})$:
    \begin{equation}
        \|\mathcal{T}f - \mathcal{T}_r f\|_{L^2} \leq \lambda_{r+1} \|f\|_{L^2}
        \label{eq:approximation_error_bound}
    \end{equation}
\end{theorem}

\begin{proof}
    By the spectral representation:
    \begin{align}
        \mathcal{T}f - \mathcal{T}_r f &= \sum_{i=r+1}^{\infty} \lambda_i \langle f, \phi_i \rangle \phi_i \\
        \|\mathcal{T}f - \mathcal{T}_r f\|_{L^2}^2 &= \sum_{i=r+1}^{\infty} \lambda_i^2 |\langle f, \phi_i \rangle|^2 \\
        &\leq \lambda_{r+1}^2 \sum_{i=r+1}^{\infty} |\langle f, \phi_i \rangle|^2 \\
        &\leq \lambda_{r+1}^2 \|f\|_{L^2}^2
    \end{align}
    where the last inequality uses Parseval's identity.
\end{proof}

\subsubsection{Physical and Geometric Interpretation}

The spectral analysis reveals that attention mechanisms naturally implement principles of harmonic analysis:

\begin{enumerate}
    \item \textbf{Decomposition into modes}: The eigenfunctions $\phi_i$ represent fundamental "attention modes"
    \item \textbf{Hierarchy of importance}: The eigenvalues $\lambda_i$ order these modes by importance
    \item \textbf{Implicit regularization}: The decay of eigenvalues provides automatic regularization
    \item \textbf{Frequency interpretation}: Modes with large eigenvalues correspond to "low-frequency" patterns in the feature space
\end{enumerate}

\textbf{Connection with Fourier Analysis.}\label{rem:fourier_connection} When the attention kernel is translation-invariant, the spectral decomposition reduces to classical Fourier analysis, where the eigenfunctions are complex exponentials and the eigenvalues are related to the Fourier transform of the kernel.

\subsubsection{Computational Applications}

The spectral analysis has important practical implications:

\begin{itemize}
    \item \textbf{Model compression}: Truncating the spectral expansion reduces the computational cost
    \item \textbf{Capacity analysis}: The effective rank indicates the representational capacity
    \item \textbf{Initialization}: The spectral modes can guide the initialization of weights
    \item \textbf{Regularization}: Penalizing large eigenvalues prevents overfitting
\end{itemize}

\begin{theorem}[Approximation by Finite Attention]\label{theorem:finite_attention_approximation}
    Let $\{x_i\}_{i=1}^n$ be an $\epsilon$-dense discretization of $\mathcal{M}$. The discrete attention operator:
    \begin{equation}
        \mathcal{T}_n f(x_i) = \sum_{j=1}^n \alpha_{ij} f(x_j), \quad \alpha_{ij} = \frac{\exp(K(x_i, x_j))}{\sum_{k=1}^n \exp(K(x_i, x_k))}
    \end{equation}
    converges to the continuous operator $\mathcal{T}$ in operator norm as $n \to \infty$ and $\epsilon \to 0$.
\end{theorem}

\begin{proof}[Sketch]
Let $f \in L^2(\mathcal{M})$ and $\mathcal{T}f(x) = \int_{\mathcal{M}} \alpha(x,y) f(y) \, d\mu(y)$ the continuous operator. The discrete approximation is $\mathcal{T}_n f(x_i) = \sum_{j=1}^n \alpha_{ij} f(x_j)$, which can be seen as a quadrature rule. For an $\epsilon$-dense discretization, the density of the points ensures that $|\mu(B_\epsilon(x_i)) - 1/n| \to 0$. The error decomposes into: (1) quadrature error $\|\mathcal{T}f - \mathcal{T}_n f\| = O(\epsilon^s)$ for $C^s$ kernels, by convergence of Riemann sums over compact manifolds; and (2) normalization error $|\sum_k \exp(K(x_i,x_k)) - \int \exp(K(x_i,y)) d\mu(y)| \to 0$ by the same argument. Under kernel regularity conditions (Lipschitz in both arguments), both errors converge to zero uniformly.
\end{proof}

This formulation reveals that standard \glspl{Transformer} can be seen as discrete approximations of integral operators over manifolds, providing a solid theoretical basis for geometric extensions.

\subsection{Functional Analysis of Attention}\label{subsec:functional_analysis_attention}

The theory of reproducing kernel Hilbert spaces (RKHS) provides a rigorous framework for understanding the functional properties of attention mechanisms.

\subsubsection{Attention as an Operator in RKHS}

\begin{definition}[Attention Kernel as a Mercer Kernel]\label{def:attention_mercer_kernel}
    The attention kernel $K_{\text{att}}: \mathcal{X} \times \mathcal{X} \to \mathbb{R}$ defined by:
    \begin{equation}
        K_{\text{att}}(x, y) = \exp\left(\frac{\langle Q(x), K(y) \rangle}{\sqrt{d_k}}\right)
        \label{eq:attention_mercer_kernel}
    \end{equation}
    is a \textit{Mercer kernel} if it is symmetric and positive definite.
\end{definition}

\begin{theorem}[RKHS Induced by Attention]\label{theorem:attention_rkhs}
    If $K_{\text{att}}$ is a Mercer kernel, then there exists a Hilbert space $\mathcal{H}_{\text{att}}$ with reproducing kernel such that:
    \begin{enumerate}
        \item $K_{\text{att}}(x, \cdot) \in \mathcal{H}_{\text{att}}$ for all $x \in \mathcal{X}$
        \item $\langle f, K_{\text{att}}(x, \cdot) \rangle_{\mathcal{H}_{\text{att}}} = f(x)$ for all $f \in \mathcal{H}_{\text{att}}$ (reproducing property)
        \item Attention can be expressed as: $\text{Attention}(x) = \sum_{i} \alpha_i K_{\text{att}}(x, x_i) V_i$
    \end{enumerate}
\end{theorem}

\begin{proof}
    The proof follows the standard Moore-Aronszajn construction:

    \textbf{Step 1}: We define the pre-Hilbert $\mathcal{H}_0$ as the linear span of $\{K_{\text{att}}(x, \cdot) : x \in \mathcal{X}\}$.

    \textbf{Step 2}: The inner product is defined by bilinearity:
    \[\left\langle \sum_i \alpha_i K_{\text{att}}(x_i, \cdot), \sum_j \beta_j K_{\text{att}}(y_j, \cdot) \right\rangle = \sum_{i,j} \alpha_i \beta_j K_{\text{att}}(x_i, y_j)\]

    \textbf{Step 3}: $\mathcal{H}_{\text{att}}$ is the completion of $\mathcal{H}_0$ under this norm.

    The reproducing property is verified directly by construction.
\end{proof}

\subsubsection{Regularization and Stability}

\begin{proposition}[Lipschitz Stability for Attention]\label{prop:lipschitz_stability_attention}
    If the functions $Q, K, V$ are $L$-Lipschitz and the data satisfy $\|x\|, \|y\| \leq R$ (essential boundedness hypothesis, since standard self-attention is not Lipschitz on unbounded domains), then the attention operator satisfies:
    \begin{equation}
        \|\text{Attention}(x) - \text{Attention}(y)\| \leq C_L \|x - y\|
        \label{eq:lipschitz_bound_attention}
    \end{equation}
    where $C_L$ depends on $L$, $R$, $d_k$ and the number of input elements $n$.
\end{proposition}

\begin{proof}
Let $\alpha_j(x) = \frac{\exp(\langle Q(x), K(x_j)\rangle / \sqrt{d_k})}{\sum_l \exp(\langle Q(x), K(x_l)\rangle / \sqrt{d_k})}$ be the attention weights. The softmax function is Lipschitz with constant 1 in the $\ell^\infty$ norm. Since $Q$ and $K$ are $L$-Lipschitz and $\|x\| \leq R$, the function $x \mapsto \langle Q(x), K(x_j)\rangle / \sqrt{d_k}$ is Lipschitz with constant $O(LR/\sqrt{d_k})$. Composing with softmax and the linear combination with $V(x_j)$ (also $L$-Lipschitz and bounded), one obtains $\|\text{Attention}(x) - \text{Attention}(y)\| \leq C_L \|x - y\|$ where $C_L = O(nL^2R/\sqrt{d_k})$. For a detailed derivation, see \citep{kim2021lipschitz}.
\end{proof}

\begin{theorem}[Automatic Spectral Regularization]\label{theorem:spectral_regularization}
    The attention operator acts as a natural spectral regularizer, filtering high-frequency components according to:
    \begin{equation}
        \|\mathcal{T}f\|_{\dot{H}^s} \leq C_s \|f\|_{\dot{H}^{s-\sigma}}
        \label{eq:spectral_regularization}
    \end{equation}
    where $\dot{H}^s$ denotes the homogeneous Sobolev space and $\sigma > 0$ is the order of regularization determined by the smoothness of the attention kernel.
\end{theorem}

\begin{proof}
By the spectral representation (\autoref{theorem:spectral_decomposition_attention}), $\mathcal{T}f = \sum_i \lambda_i \langle f, \phi_i\rangle \phi_i$. In the Sobolev norm $\dot{H}^s$, the $i$-th component has weight $\mu_i^s$ where $\mu_i$ are the eigenvalues of the Laplacian of $\mathcal{M}$. Then:
\[
\|\mathcal{T}f\|_{\dot{H}^s}^2 = \sum_i \lambda_i^2 \mu_i^{2s} |\langle f, \phi_i\rangle|^2
\]
If the attention kernel $K$ has Sobolev regularity $\sigma$ (that is, $\lambda_i = O(\mu_i^{-\sigma})$ as $i \to \infty$), then $\lambda_i^2 \mu_i^{2s} \leq C_s^2 \mu_i^{2(s-\sigma)}$, from which $\|\mathcal{T}f\|_{\dot{H}^s} \leq C_s \|f\|_{\dot{H}^{s-\sigma}}$.
\end{proof}

\begin{remark}
The above proof implicitly assumes that the eigenfunctions $\phi_i$ of the attention operator coincide (or are closely related) with the eigenfunctions of the Laplacian of $\mathcal{M}$. This assumption holds for isotropic kernels that commute with the Laplacian (for example, Matérn kernels on compact manifolds), but it is not valid for general attention kernels.
\end{remark}

\subsubsection{Convergence and Approximation}

\begin{theorem}[Monte Carlo Convergence for Attention]\label{theorem:monte_carlo_convergence}
    Let $\mathcal{T}_n$ be the discrete attention operator with $n$ points sampled i.i.d. according to a measure $\mu$. Then:
    \begin{equation}
        \mathbb{E}[\|\mathcal{T}_n - \mathcal{T}\|_{L^2 \to L^2}^2] \leq \frac{C}{n}
        \label{eq:monte_carlo_rate}
    \end{equation}
    where $C$ is a constant that depends on the properties of the kernel and the measure.
\end{theorem}

\begin{proof}
We consider the version with $1/n$ normalization (without softmax) to facilitate the analysis. The continuous operator is $\mathcal{T}f(x) = \int_\mathcal{M} K(x,y) f(y)\, d\mu(y)$ and its Monte Carlo approximation is $\mathcal{T}_n f(x) = \frac{1}{n}\sum_{j=1}^n K(x, y_j) f(y_j)$ with $y_j \sim \mu$ i.i.d. For fixed $f$, $\mathcal{T}_n f(x)$ is the empirical mean of $n$ i.i.d. random variables $Z_j = K(x, y_j)f(y_j)$ with expectation $\mathcal{T}f(x)$. For the rate:
\[
\mathbb{E}[\|\mathcal{T}_n f - \mathcal{T}f\|_{L^2}^2] = \frac{1}{n}\int_\mathcal{M} \text{Var}(K(x,\cdot)f(\cdot))\, d\mu(x) \leq \frac{C}{n}
\]
where $C = \|K\|_\infty^2 \|f\|_{L^2}^2 \mu(\mathcal{M})$. The extension to operator norm and to the case with softmax normalization (as in the statement) requires additional covering arguments and control of the convergence of the denominator, respectively.
\end{proof}

\section{Group Theory and Equivariance}

The incorporation of geometric symmetries into attention mechanisms constitutes one of the most significant developments in geometric deep learning, allowing models to preserve fundamental structures of the data.

\subsection{Foundations of Group Theory}

The foundations of group theory (groups, group actions, equivariance) needed for equivariant attention mechanisms are developed in full in \autoref{ch:gdl}, Sections~\ref{subsec:teoria_grupos_representaciones} and~\ref{subsec:equivarianza_invarianza}.

Attention mechanisms extend these principles through adaptive operations that preserve the symmetries of arbitrary groups, overcoming the fixed-equivariance limitations of \glspl{CNN}.

\subsubsection{Relevant Groups for Geometric Attention}

In the context of attention mechanisms, the most important groups include:

\begin{enumerate}
    \item \textbf{$SO(3)$}: 3D rotations, essential for Point \glspl{Transformer} and molecular data
    \item \textbf{$SE(3)$}: Special Euclidean group, fundamental in SE(3)-\glspl{Transformer}
    \item \textbf{$E(3)$}: Full Euclidean group, including reflections
    \item \textbf{General Lie groups}: Unified framework for Lie\glspl{Transformer}
    \item \textbf{Permutation groups}: For invariance to reorderings in sets
\end{enumerate}

\autoref{theorem:equivariant_composition} of \autoref{ch:gdl} ensures that equivariant attention layers can be composed while preserving the desired symmetry properties.

\subsection{Equivariant Attention}\label{subsec:equivariant_attention}\label{sec:equivariant_attention}

\subsubsection{Definition and Construction}

Applying the definition of equivariance from \autoref{ch:gdl} (\autoref{def:equivariance_gdl}) to attention mechanisms, an attention model is $G$-equivariant if it preserves the group transformations in queries, keys, and values.

\begin{theorem}[Construction of Equivariant Attention]
    For a group $G$ acting on the feature space, an attention mechanism is $G$-equivariant if and only if:
    \begin{enumerate}
        \item The compatibility functions are $G$-invariant
        \item The query, key, and value transformations commute with the group action
    \end{enumerate}
\end{theorem}

\begin{proof}[Proof sketch]
\textbf{Forward direction} ($\Leftarrow$): If the compatibility functions $a(\mathbf{q}_i, \mathbf{k}_j)$ are $G$-invariant, then $a(\rho(g)\mathbf{q}_i, \rho(g)\mathbf{k}_j) = a(\mathbf{q}_i, \mathbf{k}_j)$ for all $g \in G$, so the attention weights $\alpha_{ij}$ are invariant. If, moreover, the value transformations commute with $G$ ($W_V \rho(g) = \rho(g) W_V$), the output $\sum_j \alpha_{ij} W_V \rho(g)\mathbf{x}_j = \rho(g) \sum_j \alpha_{ij} W_V \mathbf{x}_j$, which gives equivariance.

\textbf{Converse direction} ($\Rightarrow$): If the attention is $G$-equivariant, the weights $\alpha_{ij}$ must be invariant under the simultaneous action of $G$ on Q and K (since they determine the linear combination), which requires invariance of the compatibility functions. The equivariance of the output with fixed weights implies that $W_V$ commutes with $\rho(g)$.
\end{proof}

\begin{lemma}[Sufficient Condition for Equivariance in Self-Attention]\label{lemma:self_attention_equivariance}
Let $G$ be a group acting on the input space $\mathbb{R}^{n \times d}$ via permutations.
If the projection matrices satisfy $\mathbf{W}_Q = \mathbf{W}_K = \mathbf{W}_V = \mathbf{W}$
and $\mathbf{W}$ commutes with the group action, then the self-attention mechanism
is $G$-equivariant.
\end{lemma}

\begin{proof}
For any permutation $P \in G$ and input $\mathbf{X} \in \mathbb{R}^{n \times d}$, we have:
\begin{align}
\text{SelfAttention}(P\mathbf{X}) &= \text{softmax}\left(\frac{(P\mathbf{X}\mathbf{W})(P\mathbf{X}\mathbf{W})^T}{\sqrt{d_k}}\right)(P\mathbf{X}\mathbf{W}) \\
&= \text{softmax}\left(\frac{P(\mathbf{X}\mathbf{W})(\mathbf{X}\mathbf{W})^T P^T}{\sqrt{d_k}}\right)P(\mathbf{X}\mathbf{W}) \\
&= P \cdot \text{SelfAttention}(\mathbf{X})
\end{align}
where the second step uses commutativity and the third the equivariance of softmax to permutations.
\end{proof}

\subsubsection{Irreducible Representations}

The theory of irreducible representations (\autoref{def:representacion_irreducible}, \autoref{ch:gdl}) provides the fundamental mathematical framework for designing equivariant attention mechanisms under compact Lie groups. This theory, developed initially by \citet{Peter1927} and extended to the context of deep learning by \citet{kondor2018generalizationconvolutionalneural}, establishes that any representation of a compact group decomposes into fundamental irreducible blocks.

By the Peter-Weyl theorem (\autoref{thm:peter_weyl}, \autoref{ch:gdl}), every representation decomposes into a direct sum of irreducible representations:
\begin{equation}
V = \bigoplus_{\lambda \in \widehat{G}} n_\lambda \mathcal{H}_\lambda
\label{eq:peter_weyl_decomposition_attention}
\end{equation}
where $\widehat{G}$ indexes the inequivalent irreducible representations, $\mathcal{H}_\lambda$ are the representation spaces, and $n_\lambda$ are the multiplicities.

\paragraph{Irreducible Representations of SO(3)}

For the group of three-dimensional rotations $SO(3)$, the irreducible representations are classified by a non-negative integer $\ell \in \mathbb{N}_0$ called the \textit{degree} or \textit{angular momentum}. Each irreducible representation has dimension $2\ell + 1$.

\begin{theorem}[Classification of Irreps of SO(3)]\label{thm:so3_irreps}
The irreducible representations of $SO(3)$ are indexed by $\ell \in \mathbb{N}_0$ and have dimension $d_\ell = 2\ell + 1$. The representation space $\mathcal{H}_\ell$ is generated by the spherical harmonics $Y_\ell^m$ with $m \in \{-\ell, -\ell+1, \ldots, \ell-1, \ell\}$.
\end{theorem}

\begin{proof}[Sketch]
The classification is established through the relation between $SO(3)$ and $SU(2)$. The Lie algebra $\mathfrak{so}(3) \cong \mathfrak{su}(2)$ has generators $\{J_x, J_y, J_z\}$ with commutation relations $[J_i, J_j] = i\epsilon_{ijk}J_k$. The Casimir operator $\mathbf{J}^2 = J_x^2 + J_y^2 + J_z^2$ commutes with the entire algebra, and by Schur's lemma it acts as a scalar on each irreducible. The simultaneous diagonalization of $\mathbf{J}^2$ and $J_z$ produces eigenvalues $\ell(\ell+1)$ and $m$ respectively, with $-\ell \leq m \leq \ell$, giving dimension $2\ell+1$. The spherical harmonics $Y_\ell^m$ realize these representations as functions on $S^2$. See \citep{brocker1985representations} for a complete proof.
\end{proof}

\begin{definition}[Spherical Harmonics]\label{def:spherical_harmonics_attention}
The \textit{spherical harmonics} $Y_\ell^m: S^2 \to \mathbb{C}$ of degree $\ell$ and order $m$ are eigenfunctions of the Laplacian operator on the sphere:
\begin{equation}
Y_\ell^m(\theta, \phi) = \sqrt{\frac{2\ell+1}{4\pi}\frac{(\ell-m)!}{(\ell+m)!}} P_\ell^m(\cos\theta) e^{im\phi}
\label{eq:spherical_harmonics_definition}
\end{equation}
where $P_\ell^m$ are the associated Legendre polynomials. Under a rotation $R \in SO(3)$, the spherical harmonics transform according to:
\begin{equation}
(R \cdot Y_\ell^m)(\mathbf{x}) = Y_\ell^m(R^{-1}\mathbf{x}) = \sum_{m'=-\ell}^{\ell} D_\ell^{m'm}(R) Y_\ell^{m'}(\mathbf{x})
\label{eq:spherical_harmonics_transformation}
\end{equation}
where $D_\ell^{m'm}(R)$ are the Wigner matrices.
\end{definition}

\begin{definition}[Wigner Matrices]\label{def:wigner_matrices}
The \textit{Wigner matrices} $D_\ell: SO(3) \to GL(2\ell+1, \mathbb{C})$ provide the explicit matrix representation of rotations in the space of spherical harmonics of degree $\ell$. For rotations parametrized by Euler angles $(\alpha, \beta, \gamma)$, they have the form:
\begin{equation}
D_\ell^{m'm}(\alpha, \beta, \gamma) = e^{-im'\alpha} d_\ell^{m'm}(\beta) e^{-im\gamma}
\label{eq:wigner_d_matrices}
\end{equation}
where $d_\ell^{m'm}(\beta)$ are the reduced Wigner matrices, which depend only on the rotation angle $\beta$.
\end{definition}

\begin{figure}[ht]
    \centering
    \includegraphics[width=0.85\textwidth]{figures/spherical_harmonics.pdf}
    \caption{Spherical harmonics $Y_\ell^m(\theta, \varphi)$ for degrees $\ell = 0, 1, 2$. The blue and red regions indicate positive and negative values, respectively. The degree $\ell$ controls the angular frequency (number of sign changes) and the order $m$ determines the orientation of the nodal pattern. These functions form the basis of the irreducible representations of $SO(3)$ of dimension $2\ell+1$ used in the decomposition of equivariant features~\eqref{eq:irrep_decomposition}.}
    \label{fig:spherical_harmonics_attention}
\end{figure}

In the context of equivariant attention mechanisms, the features in each layer decompose according to representation types:

\begin{equation}
    \mathbf{f} = \bigoplus_{\ell=0}^{L_{\max}} \mathbf{f}^{(\ell)}, \quad \mathbf{f}^{(\ell)} \in \mathbb{R}^{n_\ell \times (2\ell+1)}
    \label{eq:irrep_decomposition}
\end{equation}

where $\mathbf{f}^{(\ell)}$ contains $n_\ell$ channels of type $\ell$, each transforming as a spherical harmonic of degree $\ell$ under rotations.

\paragraph{Clebsch-Gordan Coefficients and Tensor Products}

The composition of features of different types requires the tensor product of irreducible representations. The Clebsch-Gordan coefficients describe how this product decomposes.

\begin{theorem}[Tensor Product Decomposition]\label{thm:tensor_product_so3}
The tensor product of two irreducible representations of $SO(3)$ of degrees $\ell_1$ and $\ell_2$ decomposes as:
\begin{equation}
D^{(\ell_1)} \otimes D^{(\ell_2)} = \bigoplus_{\ell=|\ell_1-\ell_2|}^{\ell_1+\ell_2} D^{(\ell)}
\label{eq:tensor_product_decomposition_attention}
\end{equation}
This decomposition is fundamental for constructing equivariant operations that combine features of different types. The proof is obtained by applying the rule of addition of angular momenta, see \citep{brocker1985representations}.
\end{theorem}

\begin{definition}[Clebsch-Gordan Coefficients]\label{def:clebsch_gordan_attention}
The \textit{Clebsch-Gordan coefficients} $C_{l_1 m_1, l_2 m_2}^{l m}$ (see notation in the file \texttt{notation.tex}) provide the expansion coefficients of the product of spherical harmonics:
\begin{equation}
Y_{\ell_1}^{m_1} \otimes Y_{\ell_2}^{m_2} = \sum_{\ell=|\ell_1-\ell_2|}^{\ell_1+\ell_2} \sum_{m=-\ell}^{\ell} C_{l_1 m_1, l_2 m_2}^{l m} Y_{\ell}^{m}
\label{eq:clebsch_gordan_expansion}
\end{equation}
These coefficients satisfy orthogonality relations and can be computed via recursive formulas or precomputed tables.
\end{definition}

\begin{example}[Low-Degree Cases]\label{ex:low_degree_clebsch_gordan}
To illustrate concretely, consider some low-degree cases:

\textbf{Product $\ell_1 = 0, \ell_2 = 1$} (scalar $\times$ vector):
\begin{equation}
D^{(0)} \otimes D^{(1)} = D^{(1)}
\end{equation}
This product preserves the vector type, simply scaling the vector.

\textbf{Product $\ell_1 = 1, \ell_2 = 1$} (vector $\times$ vector):
\begin{equation}
D^{(1)} \otimes D^{(1)} = D^{(0)} \oplus D^{(1)} \oplus D^{(2)}
\end{equation}
This product generates a scalar ($\ell=0$, the dot product), a pseudovector ($\ell=1$, the cross product), and a symmetric traceless tensor ($\ell=2$).
\end{example}

\textbf{Application in Equivariant Attention.} In equivariant attention mechanisms such as SE(3)-Transformers \citep{fuchs2020se3transformers} and Equiformer \citep{liao2022equiformer}, the tensor products via Clebsch-Gordan coefficients allow:
\begin{itemize}
    \item \textbf{Combining queries and keys}: $\mathbf{q}_i^{(\ell_1)} \otimes \mathbf{k}_j^{(\ell_2)}$ to compute equivariant attention scores
    \item \textbf{Mixing geometric information}: Incorporating relative directions $\hat{\mathbf{r}}_{ij}$ (type $\ell=1$) with node features
    \item \textbf{Maintaining exact equivariance}: Each operation respects the group structure of $SO(3)$ by construction
\end{itemize}

\begin{example}[Feature Space Decomposition]\label{ex:feature_space_decomposition}
Consider a feature space with $L_{\max} = 2$:
\begin{align}
\mathbf{f} &= \mathbf{f}^{(0)} \oplus \mathbf{f}^{(1)} \oplus \mathbf{f}^{(2)} \\
&= \underbrace{\mathbb{R}^{n_0 \times 1}}_{\text{scalars}} \oplus \underbrace{\mathbb{R}^{n_1 \times 3}}_{\text{vectors}} \oplus \underbrace{\mathbb{R}^{n_2 \times 5}}_{\text{degree-2 tensors}}
\label{eq:feature_space_example}
\end{align}

An equivariant attention operation computes:
\begin{equation}
e_{ij} = \sum_{\ell_q, \ell_k, \ell_{\text{out}}} \text{MLP}_{\ell_q, \ell_k, \ell_{\text{out}}} \left( \mathbf{q}_i^{(\ell_q)} \otimes \mathbf{k}_j^{(\ell_k)} \right)^{(\ell_{\text{out}})}
\label{eq:equivariant_attention_score_irreps}
\end{equation}

where the tensor product decomposes according to Equation~\eqref{eq:tensor_product_decomposition_attention} and the MLPs operate independently on each representation type.
\end{example}

\subsection{Equivariant Architectures}\label{sec:equivariant_transformers}

\subsubsection[SE(3)-Transformers]{SE(3)-\glspl{Transformer}}

The SE(3)-\glspl{Transformer} of \citet{fuchs2020se3transformers} extend attention to be equivariant under rotations and translations in 3D.

For scalar features $f_0$, vector features $\mathbf{f}_1$, and higher-order tensor features, the attention is constructed as:

\begin{align}
    f_0^{\text{out}} &= \sum_j \alpha_{ij} f_{0,j} \\
    \mathbf{f}_1^{\text{out}} &= \sum_j \alpha_{ij} \mathbf{f}_{1,j} \\
    \alpha_{ij} &= \text{softmax}(\text{\gls{MLP}}(f_{0,i}, f_{0,j}, \|\mathbf{r}_{ij}\|, \mathbf{f}_{1,i} \cdot \mathbf{f}_{1,j}))
    \label{eq:se3_transformer_attention}
\end{align}

where $\mathbf{r}_{ij} = \mathbf{r}_i - \mathbf{r}_j$ is the relative position vector.

\begin{theorem}[SE(3)-Equivariance of Attention]
    The attention mechanism defined in equation~\eqref{eq:se3_transformer_attention} is equivariant under the group $SE(3)$ of Euclidean rotations and translations \citep{fuchs2020se3transformers}.
\end{theorem}

\begin{proof}
Let $(R, \mathbf{t}) \in SE(3)$ be a roto-translation. The attention weights $\alpha_{ij}$ depend on: (i) scalars $f_{0,i}, f_{0,j}$ (invariant by definition), (ii) the norm $\|\mathbf{r}_{ij}\| = \|\mathbf{r}_i - \mathbf{r}_j\|$ which is invariant under $SE(3)$: $\|R\mathbf{r}_i + \mathbf{t} - R\mathbf{r}_j - \mathbf{t}\| = \|R(\mathbf{r}_i - \mathbf{r}_j)\| = \|\mathbf{r}_i - \mathbf{r}_j\|$, and (iii) $\mathbf{f}_{1,i} \cdot \mathbf{f}_{1,j}$ which is invariant since $\mathbf{f}_1$ transforms as a type-1 vector under $SO(3)$: $(R\mathbf{f}_{1,i}) \cdot (R\mathbf{f}_{1,j}) = \mathbf{f}_{1,i} \cdot \mathbf{f}_{1,j}$. Therefore, the scores $\alpha_{ij}$ are \emph{invariant} under $SE(3)$.

For the vector output: $\mathbf{f}_1^{\text{out}} = \sum_j \alpha_{ij} \mathbf{f}_{1,j}$. Under $(R, \mathbf{t})$: $\sum_j \alpha_{ij} R\mathbf{f}_{1,j} = R \sum_j \alpha_{ij} \mathbf{f}_{1,j} = R \mathbf{f}_1^{\text{out}}$. Thus, the type-1 output is \emph{equivariant}. The scalar (type-0) outputs are trivially invariant.
\end{proof}

\subsubsection{Equiformer: Equivariant \gls{MLP} Attention}

The Equiformer of \citet{liao2022equiformer} surpasses the traditional dot product through equivariant \gls{MLP}-based attention.

Instead of the dot product $\mathbf{q}_i^T \mathbf{k}_j$, it uses:

\begin{equation}
    e_{ij} = \text{\gls{MLP}}_{\text{attn}}(\mathbf{q}_i \otimes \mathbf{k}_j, \|\mathbf{r}_{ij}\|)
    \label{eq:equiformer_attention_scores}
\end{equation}

where $\otimes$ denotes the tensor product of irreducible features.

The tensor product of irreducible representations decomposes as:

\begin{equation}
    D^{(\ell_1)} \otimes D^{(\ell_2)} = \bigoplus_{\ell=|\ell_1-\ell_2|}^{\ell_1+\ell_2} D^{(\ell)}
    \label{eq:tensor_product_decomposition}
\end{equation}

This construction preserves SO(3) equivariance while allowing richer interactions than the simple dot product.

\subsubsection{Lie\gls{Transformer}: General Framework}

The Lie\gls{Transformer} of \citet{hutchinson2020lietransformer} provides a unified framework for attention equivariant to arbitrary Lie groups.

The approach is based on ``lifting'' the representations to the group:

\begin{enumerate}
    \item \textbf{Lift}: $\mathbf{x} \mapsto \{(g, \mathbf{x})\}_{g \in G}$
    \item \textbf{Process}: Apply attention in the group space
    \item \textbf{Project}: Reduce back to the original space
\end{enumerate}

\begin{equation}
    \text{LieAttn}(\mathbf{x}) = \text{Project}(\text{Attention}(\text{Lift}(\mathbf{x})))
    \label{eq:lie_attention}
\end{equation}

For a group $G$, the left regular representation is defined as:

\begin{equation}
    (L_g f)(h) = f(g^{-1}h), \quad g, h \in G, f \in L^2(G)
    \label{eq:left_regular_representation}
\end{equation}

This construction guarantees equivariance for any Lie group.

\subsubsection{EPT: Equivariant Foundation Models}

The Equivariant Pretrained Transformer (EPT) of \citet{jiao2024ept} represents the first foundation model that unifies the geometric learning of small molecules and proteins under an equivariant framework.

\begin{definition}[Multi-Domain Equivariant Pretraining]
    A foundation model $F_{\text{EPT}}$ is multi-domain equivariant if for domains $\mathcal{D}_1, \mathcal{D}_2$ (e.g., molecules, proteins) and group $G$:
    \begin{equation}
        F_{\text{EPT}}(g \cdot \mathbf{x}) = g \cdot F_{\text{EPT}}(\mathbf{x}), \quad \forall \mathbf{x} \in \mathcal{D}_1 \cup \mathcal{D}_2, g \in G
        \label{eq:multi_domain_equivariance}
    \end{equation}
\end{definition}

The EPT architecture incorporates enhanced representation blocks:

\begin{equation}
    \mathbf{h}_i^{\text{EPT}} = \text{E(3)-TransformerBlock}(\mathbf{h}_i, \{\mathbf{h}_j\}_{j \in \mathcal{N}_{\text{ext}}(i)})
    \label{eq:ept_enhanced_block}
\end{equation}

EPT uses multiple pretraining objectives that preserve geometry:

\begin{align}
    \mathcal{L}_{\text{EPT}} &= \mathcal{L}_{\text{masked}} + \mathcal{L}_{\text{geom}} + \mathcal{L}_{\text{cross}} \\
    \mathcal{L}_{\text{geom}} &= \|\mathbf{f}(g \cdot \mathbf{x}) - g \cdot \mathbf{f}(\mathbf{x})\|^2 \\
    \mathcal{L}_{\text{cross}} &= \text{KL}(\mathbf{p}_{\text{mol}}, \mathbf{p}_{\text{prot}})
    \label{eq:ept_loss}
\end{align}

\begin{proposition}[Transfer of Geometric Knowledge]
    If a model $F$ is pretrained on a domain $\mathcal{D}_1$ with $G$-equivariance, then the transfer to a domain $\mathcal{D}_2$ with the same symmetry $G$ preserves the equivariance properties \citep{jiao2024ept}.
\end{proposition}

\begin{proof}
If $F$ is $G$-equivariant, this means its layers are built with equivariant operations (tensor products, group convolutions, equivariant linear layers). Equivariance is an \emph{architectural} property; it does not depend on the training data nor on the values of the learned parameters. When transferring $F$ to $\mathcal{D}_2$ (with or without fine-tuning), the structure of the layers does not change, so $F(g \cdot \mathbf{x}) = g \cdot F(\mathbf{x})$ holds for all $g \in G$ and $\mathbf{x} \in \mathcal{D}_2$. The necessary condition is that both domains admit the same action of the group $G$.
\end{proof}

\subsection{Equivariant Attention on Point Clouds}

3D point clouds present unique challenges for attention due to their unstructured nature and the need to preserve equivariance under rigid transformations.

\subsubsection{Point \gls{Transformer}: Vector Attention}

The Point \gls{Transformer} of \citet{zhao2021point} introduces vector attention specifically designed for point clouds:

\begin{equation}
    \mathbf{y}_i = \sum_{j \in \mathcal{N}(i)} \rho(\gamma(\phi(\mathbf{x}_i) - \psi(\mathbf{x}_j) + \delta))(\alpha(\mathbf{x}_j) + \delta)
    \label{eq:point_transformer_attention}
\end{equation}

where:
\begin{itemize}
    \item $\phi, \psi, \alpha$ are linear transformations
    \item $\delta = \text{\gls{MLP}}(\mathbf{p}_i - \mathbf{p}_j)$ encodes positional information
    \item $\gamma, \rho$ are \glspl{MLP}
\end{itemize}

The function $\delta$ captures the local geometry:

\begin{equation}
    \delta_{ij} = \text{\gls{MLP}}(\mathbf{p}_i - \mathbf{p}_j, \|\mathbf{p}_i - \mathbf{p}_j\|)
    \label{eq:positional_encoding_3d}
\end{equation}

The Point \gls{Transformer} formulation can be understood as an instance of the equivariant attention framework developed in \autoref{subsec:equivariant_attention}. The positional encoding $\delta_{ij}$ depends exclusively on the relative position $\mathbf{p}_i - \mathbf{p}_j$, which is invariant under translations. Likewise, if the \gls{MLP} that computes $\delta$ receives as input only the relative difference and its norm, the architecture is invariant under rotations in the computation of the attention weights, connecting directly with the invariance condition of the compatibility functions of \autoref{subsec:equivariant_attention}. Vector attention (in contrast to standard scalar attention) allows different dimensions of the attention vector to independently modulate the features, providing greater expressive capacity than a single scalar weight $\alpha_{ij}$ per pair of points.

A computational advantage of the Point \gls{Transformer} over architectures such as SE(3)-\glspl{Transformer} lies in avoiding the tensor product of irreducible representations (Equation~\eqref{eq:tensor_product_decomposition_attention}), which has cost $O(\ell_{\max}^3)$ for maximum degree $\ell_{\max}$. By operating directly on differences of positions and features, the Point \gls{Transformer} achieves complexity $O(|\mathcal{N}(i)| \cdot d)$ per point, at the cost of giving up exact equivariance under $SO(3)$ for the output features (only the relative positions are invariant, not the transformed features).
\section{Graph Theory and Spectral Analysis}

The attention mechanisms developed so far operate on data with regular structure (sequences, images, point clouds). However, many problems of interest (social networks, molecules, knowledge graphs) involve data with irregular connectivity. The natural extension of attention to these domains requires combining graph theory with spectral analysis, providing solid mathematical foundations for adaptive weights on irregularly structured data.

\subsection{Graph Foundations}

For the graph theory foundations (definition of a graph, adjacency matrix, Laplacian) needed in attention mechanisms, see \autoref{ch:gdl}, \autoref{sec:grafos_discretizacion}. The spectral analysis concepts are developed in \autoref{ch:cnn}.

\subsubsection{Connection with Message Passing}

Attention mechanisms on graphs generalize the message passing paradigm:

\begin{equation}
    \text{Attention} \equiv \text{MessagePassing with adaptive weights}
    \label{eq:attention_message_passing_equivalence}
\end{equation}

\subsection{Graph Attention Networks}\label{subsec:graph_attention_networks}\label{sec:graph_attention}

\glspl{GAT} represent the natural intersection between the spectral theory of graphs developed in \autoref{sec:analisis_armonico} and adaptive attention mechanisms. While traditional spectral methods require expensive eigendecomposition, \glspl{GAT} achieve similar expressive power through learnable attention coefficients that implicitly capture the structure of the graph and long-range dependencies.

\subsubsection{Mathematical Foundations: Attention over Graphs}\label{subsubsec:fundamentos_atencion_grafos}

\begin{definition}[Labeled Graph for Attention]\label{def:labeled_graph_attention}
A \textit{labeled graph for attention} is a tuple $\mathcal{G} = (V, E, \mathbf{H})$ where:
\begin{itemize}
    \item $V = \{v_1, v_2, \ldots, v_n\}$ is the set of nodes
    \item $E \subseteq V \times V$ is the set of edges
    \item $\mathbf{H} = [\mathbf{h}_1, \mathbf{h}_2, \ldots, \mathbf{h}_n]^T \in \mathbb{R}^{n \times d}$ are the node features
\end{itemize}
The structure of the graph defines the attention constraints: typically, node $v_i$ can attend only to its neighborhood $\mathcal{N}(i) = \{j : (i,j) \in E\}$.
\end{definition}

\begin{definition}[Graph Attention Function]\label{def:graph_attention_function}
A \textit{graph attention function} $f_{\text{GAT}}: \mathcal{G} \to \mathbb{R}^{n \times d'}$ satisfies:
\begin{enumerate}
    \item \textbf{Permutation equivariance}: For any permutation $\pi$ of the nodes:
    \begin{equation}
        f_{\text{GAT}}(\pi(\mathcal{G})) = \pi(f_{\text{GAT}}(\mathcal{G}))
        \label{eq:gat_permutation_equivariance}
    \end{equation}

    \item \textbf{Respect for structure}: It only considers connections defined by $E$:
    \begin{equation}
        \frac{\partial f_{\text{GAT}}(\mathbf{h}_i)}{\partial \mathbf{h}_j} \neq \mathbf{0} \Rightarrow (i,j) \in E \text{ or } i = j
        \label{eq:gat_structural_constraint}
    \end{equation}
\end{enumerate}
\end{definition}

\subsubsection{Classical Graph Attention Networks}\label{subsubsec:gat_clasicos}

The pioneering work of \citet{velickovic2017graph} introduces the first successful formulation of attention specific to graphs.

\begin{definition}[GAT Attention Mechanism]\label{def:gat_mechanism}
For a node $i$ with neighborhood $\mathcal{N}(i)$, the \gls{GAT} mechanism computes:

\textbf{Step 1: Linear projection}
\begin{equation}
    \mathbf{z}_i = \mathbf{W}\mathbf{h}_i, \quad \mathbf{z}_j = \mathbf{W}\mathbf{h}_j
    \label{eq:gat_linear_projection}
\end{equation}
where $\mathbf{W} \in \mathbb{R}^{d' \times d}$ is a shared weight matrix.

\textbf{Step 2: Attention coefficients}
\begin{equation}
    e_{ij} = a^T[\mathbf{z}_i \| \mathbf{z}_j] = a^T[\mathbf{W}\mathbf{h}_i \| \mathbf{W}\mathbf{h}_j]
    \label{eq:gat_attention_coefficients}
\end{equation}
where $a \in \mathbb{R}^{2d'}$ is a learnable attention vector and $\|$ denotes concatenation.

\textbf{Step 3: Masked softmax normalization}
\begin{equation}
    \alpha_{ij} = \frac{\exp(\text{LeakyReLU}(e_{ij}))}{\sum_{k \in \mathcal{N}(i) \cup \{i\}} \exp(\text{LeakyReLU}(e_{ik}))}
    \label{eq:gat_softmax_normalization}
\end{equation}

\textbf{Step 4: Weighted aggregation}
\begin{equation}
    \mathbf{h}_i' = \sigma\left(\sum_{j \in \mathcal{N}(i) \cup \{i\}} \alpha_{ij} \mathbf{z}_j\right)
    \label{eq:gat_weighted_aggregation}
\end{equation}
\end{definition}

\begin{figure}[ht]
\centering
\includegraphics[width=0.75\textwidth]{figures/gat_mechanism.pdf}
\caption{Graph Attention Network (GAT) mechanism. The central node (red) computes attention coefficients $\alpha_{ij}$ with each neighbor (blue). The thickness of the arrows reflects the magnitude of the attention weights: the network learns to assign different importances to each neighbor according to the node features.}
\label{fig:gat_mechanism}
\end{figure}

\begin{proposition}[Expressive Power of GAT]\label{prop:gat_expressive_power}
Standard Graph Attention Networks are bounded in expressive power by the 1-Weisfeiler-Lehman (1-WL) test \citep{xu2019how}. This implies:
\begin{enumerate}
    \item \textbf{Upper bound}: GAT cannot distinguish pairs of non-isomorphic graphs that the 1-WL test does not distinguish.
    \item \textbf{Anisotropic aggregation}: Unlike \gls{GCN} with fixed weights, GAT can learn adaptive aggregation functions through attention coefficients, which allows greater flexibility within the class of functions accessible by message passing architectures.
    \item \textbf{Permutation equivariance}: GAT is permutation-equivariant by construction, since the attention coefficients depend only on the features of the nodes involved and not on an ordering.
\end{enumerate}
\end{proposition}

\begin{proof}
GAT is an instance of the MPNN (message passing neural network) framework with attention-weighted aggregation over the 1-hop neighborhood. \citet{xu2019how} prove that every MPNN architecture is bounded by the 1-WL test in its capacity to distinguish graphs. Permutation equivariance follows from the fact that the attention operation is invariant to the reordering of neighboring nodes: the coefficients $\alpha_{ij}$ depend on $(\mathbf{h}_i, \mathbf{h}_j)$ and the aggregation $\sum_{j \in \mathcal{N}(i)} \alpha_{ij} \mathbf{z}_j$ is symmetric with respect to the order of the neighbors.
\end{proof}

\subsubsection{Multi-Head Graph Attention}\label{subsubsec:multihead_graph_attention}

To improve stability and expressive capacity, \glspl{GAT} employ multi-head attention.

\begin{definition}[Multi-Head GAT]\label{def:multihead_gat}
A Multi-Head GAT with $h$ heads computes:

\begin{equation}
    \mathbf{h}_i' = \sigma\left(\frac{1}{h}\sum_{k=1}^{h} \sum_{j \in \mathcal{N}(i)} \alpha_{ij}^{(k)} \mathbf{W}^{(k)}\mathbf{h}_j\right)
    \label{eq:multihead_gat_average}
\end{equation}

or through concatenation:

\begin{equation}
    \mathbf{h}_i' = \sigma\left(\|_{k=1}^{h} \sum_{j \in \mathcal{N}(i)} \alpha_{ij}^{(k)} \mathbf{W}^{(k)}\mathbf{h}_j\right)
    \label{eq:multihead_gat_concat}
\end{equation}

where each head $k$ has its own parameters $\mathbf{W}^{(k)}$ and $\mathbf{a}^{(k)}$.
\end{definition}

\begin{proposition}[Advantages of Multi-Head GAT]\label{prop:multihead_gat_advantages}
The Multi-Head GAT provides:
\begin{enumerate}
    \item \textbf{Training stability}: Reduces the variance of the gradients
    \item \textbf{Feature specialization}: Different heads capture different relational aspects
    \item \textbf{Greater expressive capacity}: Approximates more complex functions with the same total number of parameters
\end{enumerate}
\end{proposition}

\begin{proof}[Justification]
\textbf{(1)} Let $\hat{\mathbf{h}}_i^{(k)} = \frac{1}{H}\sum_{h=1}^H \mathbf{h}_i^{(k,h)}$ be the averaged output. The variance of the gradient $\text{Var}[\nabla_\theta \hat{\mathbf{h}}_i]$ is reduced by a factor $1/H$ relative to a single head, due to the independence of the attention mechanisms in each head (analogous to the effect of \textit{ensemble methods}).

\textbf{(2)} Empirically, the heads tend to specialize: with $H$ heads and dimension $d/H$ each, every head operates in a different subspace of $\mathbb{R}^d$. The attention weights $\alpha_{ij}^{(h)}$ vary significantly across heads, indicating that they capture different relational patterns.

\textbf{(3)} The space of functions representable by $H$ heads of dimension $d/H$ strictly includes that of a single head of dimension $d$: the concatenation $\|_{h=1}^H \text{Attn}_h$ can represent functions that require attending simultaneously to distinct combinations of neighbors, which a single head cannot achieve.
\end{proof}

\subsubsection{Spectral Analysis of Graph Attention}\label{subsubsec:analisis_espectral_gat}

The connection between GATs and spectral analysis reveals deep insights into how they work.

\begin{definition}[Spectral Interpretation of Attention]\label{def:spectral_interpretation_attention}
Let $\mathbf{A}^{(GAT)} \in \mathbb{R}^{n \times n}$ be the attention matrix with entries $A^{(GAT)}_{ij} = \alpha_{ij}$. This matrix can be interpreted as a \textbf{learnable version of the diffusion operator}:
\begin{equation}
    \mathbf{A}^{(GAT)} \approx f(\mathbf{L}^{-1})
    \label{eq:gat_spectral_approximation}
\end{equation}
where $f$ is a learnable function and $\mathbf{L}$ is the \gls{laplaciano_grafo}.
\end{definition}

\begin{theorem}[GAT-Spectral Connection]\label{thm:gat_spectral_connection}
In the limit of many layers, a \gls{GAT} with stable attention weights converges to performing adaptive spectral filtering:
\begin{equation}
    \lim_{L \to \infty} \text{GAT}^{(L)}(\mathbf{H}) = \mathbf{U} \text{diag}(\hat{g}(\lambda_1), \ldots, \hat{g}(\lambda_n)) \mathbf{U}^T \mathbf{H}
    \label{eq:gat_spectral_limit}
\end{equation}
where $\mathbf{U}$ contains the eigenvectors of the Laplacian and $\hat{g}$ is the implicitly learned filter function.
\end{theorem}

\begin{proof}
\textbf{Step 1}: Stability of the attention weights.
In the limit, the attention weights converge: $\lim_{L \to \infty} \alpha_{ij}^{(L)} = \alpha_{ij}^*$.

\textbf{Step 2}: Limiting transition matrix.
The \gls{GAT} operation becomes: $\mathbf{H}^{(L+1)} = \mathbf{A}^* \mathbf{H}^{(L)}$ where $\mathbf{A}^*$ is the stable attention matrix.

\textbf{Step 3}: Spectral decomposition.
Under the (idealized) hypothesis that $\mathbf{A}^*$ shares eigenvectors with the Laplacian of the graph: $\mathbf{A}^* = \mathbf{U} \mathbf{\Lambda}^* \mathbf{U}^T$.

\textbf{Step 4}: Implicit filtering.
The eigenvalues $\mathbf{\Lambda}^*$ implicitly define the filter function $\hat{g}$.
\end{proof}

\begin{remark}
Step~3 of the previous proof assumes that the learned attention matrix $\mathbf{A}^*$ shares eigenvectors with the Laplacian of the graph $\mathbf{L}$, which is a simplification. In practice, $\mathbf{A}^*$ is generally asymmetric and its eigenvectors do not coincide with those of $\mathbf{L}$. This spectral connection should be understood as an analogy that holds exactly only when the attention coefficients are functions of the graph structure compatible with its spectral decomposition. Moreover, the repeated application of a fixed stochastic matrix $\mathbf{A}^*$ converges to a stationary distribution (the over-smoothing phenomenon), not to the rich spectral filtering described in equation~\eqref{eq:gat_spectral_limit}.
\end{remark}

\subsubsection{Spectral Positional Encoding}\label{subsubsec:codificacion_posicional_espectral}

A fundamental limitation of \glspl{GAT} is their inability to distinguish structurally equivalent nodes. Spectral positional encoding solves this problem.

\begin{definition}[Spectral Positional Encoding]\label{def:spectral_positional_encoding}
Let $\mathbf{L} = \mathbf{U}\mathbf{\Lambda}\mathbf{U}^T$ be the spectral decomposition of the \gls{laplaciano_grafo}. The spectral positional encoding of node $i$ is defined as:
\begin{equation}
    \text{PE}_i = [\mathbf{u}_{i,1}, \mathbf{u}_{i,2}, \ldots, \mathbf{u}_{i,k}] \in \mathbb{R}^k
    \label{eq:spectral_positional_encoding}
\end{equation}
where $\mathbf{u}_{i,j}$ is the $i$-th component of the $j$-th smallest eigenvector (excluding the trivial one).
\end{definition}

\begin{definition}[Spectral Distances for Attention]\label{def:spectral_distances_attention}
To improve attention between distant nodes, spectral metrics are defined:

\textbf{Diffusion Distance}:
\begin{equation}
    d_D^2(i,j) = \sum_{k=2}^{n} e^{-2t\lambda_k} (\mathbf{u}_{i,k} - \mathbf{u}_{j,k})^2
    \label{eq:diffusion_distance}
\end{equation}

\textbf{Green's Function}:
\begin{equation}
    G(i,j) = \sum_{k=2}^{n} \frac{\mathbf{u}_{i,k} \mathbf{u}_{j,k}}{\lambda_k}
    \label{eq:green_function_distance}
\end{equation}

\textbf{Biharmonic Distance}:
\begin{equation}
    d_B^2(i,j) = \sum_{k=2}^{n} \frac{(\mathbf{u}_{i,k} - \mathbf{u}_{j,k})^2}{\lambda_k^2}
    \label{eq:biharmonic_distance}
\end{equation}
\end{definition}

\subsubsection{Spectral Attention Networks (SAN)}\label{subsubsec:spectral_attention_networks}

\citet{kreuzer2021spectral} introduce the first fully connected Transformer for graphs that competes with specialized GNNs.

\begin{definition}[SAN Architecture]\label{def:san_architecture}
A Spectral Attention Network processes graphs through:

\textbf{Step 1: Learnable positional encoding}
\begin{equation}
    \text{LPE}_i = \text{MLP}_{\text{PE}}([\lambda_1, \ldots, \lambda_k, \mathbf{u}_{i,1}, \ldots, \mathbf{u}_{i,k}])
    \label{eq:learned_positional_encoding}
\end{equation}

\textbf{Step 2}: Dual attention (real + virtual)
\begin{align}
    \mathbf{h}_i^{\text{real}} &= \text{Attention}_{\text{sparse}}(\mathbf{h}_i, \{\mathbf{h}_j : j \in \mathcal{N}(i)\}) \\
    \mathbf{h}_i^{\text{virtual}} &= \text{Attention}_{\text{full}}(\mathbf{h}_i, \{\mathbf{h}_j : j \in V\})
    \label{eq:dual_attention_san}
\end{align}

\textbf{Step 3: Adaptive balancing}
\begin{equation}
    \mathbf{h}_i' = \gamma \cdot \mathbf{h}_i^{\text{real}} + (1-\gamma) \cdot \mathbf{h}_i^{\text{virtual}}
    \label{eq:adaptive_balancing_san}
\end{equation}
where $\gamma \in [0,1]$ is a learnable parameter.
\end{definition}

\begin{theorem}[Overcoming Over-Squashing]\label{thm:over_squashing_resolution}
Spectral Attention Networks solve the over-squashing problem by providing direct connections between arbitrarily distant nodes:
\begin{equation}
    \text{Information distance} = \mathcal{O}(1) \quad \forall i,j \in V
    \label{eq:information_distance_constant}
\end{equation}
independently of the geodesic distance in the original graph.
\end{theorem}

\begin{proof}
In a \gls{GNN} with local message passing, the information of node $j$ takes at least $d_G(i,j)$ layers to reach node $i$, where $d_G$ is the geodesic distance. After $L$ layers, the information is compressed into fixed-dimension representations $d$, causing over-squashing when $|\{j: d_G(i,j) \leq L\}|$ grows exponentially \citep{topping2022understanding}. In SAN, global attention allows node $i$ to attend directly to node $j$ in a single layer, with score $\alpha_{ij} = \text{softmax}(\mathbf{q}_i^T \mathbf{k}_j / \sqrt{d})$. The information distance $D_{\text{info}}(i,j) = -\log \alpha_{ij}$ is $O(1)$ (it depends on the similarity of representations, not on the topology of the graph), eliminating the bottleneck.
\end{proof}

\subsubsection{Scalable Extensions: GraphGPS}\label{subsubsec:graphgps}

\begin{definition}[GraphGPS Architecture]\label{def:graphgps_architecture}
GraphGPS \citep{rampasek2022graphgps} combines local message passing with global attention through:

\begin{equation}
    \mathbf{H}^{(\ell+1)} = \text{MLP}^{(\ell)}\left(\mathbf{H}_{\text{MPNN}}^{(\ell+1)} + \mathbf{H}_{\text{Global}}^{(\ell+1)}\right)
    \label{eq:graphgps_architecture}
\end{equation}

where:
\begin{align}
    \mathbf{H}_{\text{MPNN}}^{(\ell+1)} &= \text{MPNN}^{(\ell)}(\mathbf{H}^{(\ell)}) \\
    \mathbf{H}_{\text{Global}}^{(\ell+1)} &= \text{GlobalTransformer}^{(\ell)}(\mathbf{H}^{(\ell)} + \text{PE})
    \label{eq:graphgps_components}
\end{align}
\end{definition}

\begin{proposition}[Linear Scalability of GraphGPS]\label{prop:graphgps_scalability}
When a linear-complexity attention mechanism (Performer \citep{choromanski2021performer} or BigBird) is used for the global component, and excluding the cost of precomputing spectral positional encodings, GraphGPS achieves per-layer complexity $\mathcal{O}(N + E)$ through:
\begin{enumerate}
    \item \textbf{Sub-quadratic attention}: Performer or BigBird with complexity $\mathcal{O}(N)$
    \item \textbf{Local message passing}: Complexity $\mathcal{O}(E)$ for the MPNN component
    \item \textbf{Modular separation}: Decoupling of local and global information
\end{enumerate}
With standard full attention, the global component would have complexity $\mathcal{O}(N^2)$. For the precise breakdown as a function of the model parameters, see \autoref{prop:graphgps_efficient_complexity}.
\end{proposition}

\begin{proof}
The complexity is obtained by summing the contributions of each component. The MPNN component processes each edge once, resulting in $\mathcal{O}(E)$. For the global component, mechanisms like Performer approximate softmax attention through $\text{softmax}(\mathbf{Q}\mathbf{K}^\top) \approx \hat{\phi}(\mathbf{Q})\hat{\phi}(\mathbf{K})^\top$, where $\hat{\phi}: \mathbb{R}^d \to \mathbb{R}^k$ with $k \ll N$. This reduces the computation from $\mathcal{O}(N^2)$ to $\mathcal{O}(Nk) = \mathcal{O}(N)$ for fixed $k$. The modular separation allows summing both contributions independently, obtaining $\mathcal{O}(N + E)$.
\end{proof}

\subsubsection{Limitations and Future Directions}\label{subsubsec:limitaciones_gat}

\paragraph{Current Limitations}

\begin{enumerate}
    \item \textbf{1-WL limitation}: Standard GATs cannot distinguish non-isomorphic graphs that the 1-Weisfeiler-Lehman test cannot distinguish

    \item \textbf{Over-smoothing}: In deep graphs, the node representations converge to similar values:
    \begin{equation}
        \lim_{L \to \infty} \|\mathbf{h}_i^{(L)} - \mathbf{h}_j^{(L)}\| \to 0 \quad \forall i,j
        \label{eq:over_smoothing_limit}
    \end{equation}

    \item \textbf{Sign ambiguity}: The eigenvectors are defined up to multiplication by $\pm 1$, requiring techniques such as SignNet

    \item \textbf{Isospectral graphs}: Graphs with identical spectra but different structures cannot be distinguished spectrally
\end{enumerate}

\paragraph{Future Research}

\begin{definition}[Higher-Order k-GAT]\label{def:higher_order_gat}
Higher-order GATs consider $k$-tuples of nodes:
\begin{equation}
    \alpha_{i,\mathcal{S}} = \text{softmax}\left(a^T\left[\|_{j \in \mathcal{S}} \mathbf{W}\mathbf{h}_j\right]\right)
    \label{eq:higher_order_gat}
\end{equation}
where $\mathcal{S} \subseteq \mathcal{N}^{(k)}(i)$ is a $k$-subset of the $k$-neighborhood.
\end{definition}

\begin{definition}[Continuous Geometric Attention]\label{def:continuous_geometric_attention}
For graphs embedded in continuous spaces:
\begin{equation}
    e_{ij} = a(\mathbf{h}_i, \mathbf{h}_j, \|\mathbf{x}_i - \mathbf{x}_j\|, \angle(\mathbf{x}_i, \mathbf{x}_j))
    \label{eq:continuous_geometric_attention}
\end{equation}
incorporating metric and angular information.
\end{definition}

\subsubsection{Connections with Previous Chapters}

The developments in GATs connect directly with established concepts:

\begin{itemize}
    \item \textbf{Spectral Analysis (\autoref{sec:analisis_armonico})}: GATs implicitly implement adaptive spectral filtering
    \item \textbf{Message Passing (Section~\ref{subsec:message_passing_unificador})}: GATs are a specific instantiation of the unifying framework with attention-based aggregation
    \item \textbf{Equivariance (Section~\ref{subsec:equivarianza_invarianza})}: GATs preserve permutation equivariance while gaining adaptive flexibility
\end{itemize}

\begin{proposition}[Synthesis of GAT in the Unifying Framework]\label{prop:gat_unified_framework}
Graph Attention Networks are naturally expressed in the Geometric Message Passing framework:
\begin{align}
    \mathbf{h}_i^{(k+1)} &= \phi^{(k)}\left(\mathbf{h}_i^{(k)}, \bigoplus_{j \in \mathcal{N}(i)} \alpha_{ij} \cdot \psi^{(k)}(\mathbf{h}_j^{(k)})\right) \\
    \alpha_{ij} &= \text{Attention}(\mathbf{h}_i^{(k)}, \mathbf{h}_j^{(k)}, e_{ij})
    \label{eq:gat_in_unified_framework}
\end{align}
where the attention weights $\alpha_{ij}$ replace the fixed aggregation with learnable adaptive aggregation.
\end{proposition}

\begin{proof}
Let $G = (V, E)$ be a graph and $\mathbf{h}_i^{(k)}$ the representation of node $i$ in layer $k$. In the standard MPNN framework, the update takes the form $\mathbf{h}_i^{(k+1)} = \phi(\mathbf{h}_i^{(k)}, \bigoplus_{j \in \mathcal{N}(i)} \psi(\mathbf{h}_j^{(k)}))$. In a GAT, the message function is $\psi^{(k)}(\mathbf{h}_j) = \mathbf{W}^{(k)} \mathbf{h}_j$ and the aggregation is $\bigoplus = \sum_{j} \alpha_{ij} (\cdot)$, where $\alpha_{ij} = \operatorname{softmax}_j(e_{ij})$ with $e_{ij} = \text{LeakyReLU}(\mathbf{a}^\top [\mathbf{W}\mathbf{h}_i \| \mathbf{W}\mathbf{h}_j])$. Substituting into the general MPNN equation directly yields the expression in the proposition.
\end{proof}

This formulation shows that GATs represent the natural evolution of message passing from rigid aggregation toward flexible aggregation, maintaining equivariance guarantees while gaining adaptive expressiveness.

\subsection{Spectral Analysis on Graphs}

For the spectral analysis foundations of the Laplacian, including eigenvalues, eigenvectors, and spectral decomposition, see \autoref{ch:cnn}, Detailed Spectral Analysis Section. The spectral positional encoding (\autoref{def:spectral_positional_encoding}) and the spectral distances (\autoref{def:spectral_distances_attention}) introduced earlier constitute the fundamental tools for incorporating structural information of the graph into attention mechanisms.

\subsection{Advanced Graph Architectures}\label{subsec:arquitecturas_avanzadas_grafos}

Advanced graph architectures represent the confluence of attention mechanisms with the fundamental principles of \gls{GDL} developed in \autoref{ch:gdl}. This section examines how different architectural designs integrate spectral, structural, and geometric information to overcome the limitations of basic architectures such as GAT. Each architecture addresses specific challenges: Spectral Attention Networks solve the problem of positional ambiguity, GraphGPS combines local and global processing efficiently, Graph \glspl{Transformer} incorporate topological structure into attention, and Structure-Aware \glspl{Transformer} exploit subgraph patterns. We will unify these architectures by showing how they instantiate the Message Passing framework (Section~\ref{subsec:message_passing_unificador}) with different aggregation mechanisms and inductive biases.

\subsubsection{Spectral Attention Networks}\label{subsubsec:san_advanced}

The Spectral Attention Networks of \citet{kreuzer2021spectral} constitute the first fully connected \gls{Transformer} for graphs that competes with specialized \glspl{GNN}. Their fundamental contribution lies in the use of learnable positional encodings based on the spectrum of the graph Laplacian.

\paragraph{Theoretical foundation: the isotropy problem}

Standard \glspl{Transformer} operate on sequences with an implicit natural order, which provides intrinsic positional information. Graphs, however, lack a canonical ordering: they are invariant under node permutations. This invariance, although desirable from the perspective of equivariance (see Section~\ref{subsec:equivarianza_invarianza}), eliminates all structural positional information. The central challenge is: \textit{how can a \gls{Transformer} be endowed with positional information in a domain without a natural order?}

SAN's answer is to use the \textbf{spectrum of the graph Laplacian}. The Spectral Theorem (\autoref{ch:cnn}) guarantees that for the symmetric normalized Laplacian $\tilde{L} = I - D^{-1/2}AD^{-1/2}$ there exists a decomposition:
\begin{equation}
    \tilde{L} = \sum_{k=1}^{n} \lambda_k \mathbf{v}_k \mathbf{v}_k^T
    \label{eq:spectral_decomposition_san}
\end{equation}
where $0 = \lambda_1 \leq \lambda_2 \leq \cdots \leq \lambda_n \leq 2$ are the eigenvalues and $\{\mathbf{v}_k\}_{k=1}^{n}$ the corresponding orthonormal eigenvectors.

\begin{definition}[Learnable Spectral Positional Encoding]\label{def:learned_spectral_pe}
For a graph $G$ with Laplacian $\tilde{L}$, the learnable spectral positional encoding for node $i$ is defined as:
\begin{equation}
    \mathbf{PE}_i = \text{MLP}_{\text{enc}}(\lambda_1, \ldots, \lambda_k, \mathbf{v}_{1,i}, \ldots, \mathbf{v}_{k,i}) \in \mathbb{R}^{d_{\text{model}}}
    \label{eq:san_learned_positional_encoding}
\end{equation}
where:
\begin{itemize}
    \item $k \ll n$ is the number of spectral components used (typically $k \in [8, 20]$)
    \item $\text{MLP}_{\text{enc}}$ is a multilayer network with a nonlinear activation function
    \item $\lambda_1, \ldots, \lambda_k$ are the first $k$ eigenvalues (global information of the graph)
    \item $\mathbf{v}_{k,i}$ is the $i$-th component of the eigenvector $\mathbf{v}_k$ (local positional information of the node)
\end{itemize}
\end{definition}

\textbf{Spectral justification.}\label{rem:justification_spectral_pe} The first eigenvectors of the Laplacian capture fundamental topological structure:
\begin{itemize}
    \item $\mathbf{v}_1 = \frac{1}{\sqrt{n}}\mathbf{1}$ (constant vector, connected component)
    \item $\mathbf{v}_2$ (Fiedler vector) identifies communities and minimum cuts
    \item Subsequent eigenvectors capture geometric patterns of increasing frequency
\end{itemize}
This approach connects with harmonic analysis on graphs (\autoref{ch:cnn}), where the eigenvectors form a Fourier basis for signals over the graph.

\paragraph{SAN architecture: dual attention with adaptive balancing}

The complete SAN architecture implements a \textbf{dual attention} scheme that combines local information (graph neighborhood) with global information (the full graph):

\begin{definition}[SAN Layer]\label{def:san_layer}
A SAN layer processes node features $\mathbf{H}^{(\ell)} \in \mathbb{R}^{n \times d}$ through three stages:

\textbf{Stage 1: Positional encoding injection}
\begin{equation}
    \tilde{\mathbf{H}}^{(\ell)} = \mathbf{H}^{(\ell)} + \mathbf{PE}
    \label{eq:san_pe_injection}
\end{equation}

\textbf{Stage 2: Dual attention}
\begin{align}
    \mathbf{H}_{\text{local}}^{(\ell+1)} &= \text{MultiHeadAttn}_{\text{sparse}}\left(\tilde{\mathbf{H}}^{(\ell)}, \mathcal{E}\right) \label{eq:san_local_attn} \\
    \mathbf{H}_{\text{global}}^{(\ell+1)} &= \text{MultiHeadAttn}_{\text{full}}\left(\tilde{\mathbf{H}}^{(\ell)}\right) \label{eq:san_global_attn}
\end{align}
where $\mathcal{E}$ denotes the edges of the graph and the sparse attention only computes weights for pairs $(i,j) \in \mathcal{E}$.

\textbf{Stage 3: Adaptive combination}
\begin{equation}
    \mathbf{H}^{(\ell+1)} = \text{MLP}\left(\text{LayerNorm}\left(\gamma^{(\ell)} \mathbf{H}_{\text{local}}^{(\ell+1)} + (1-\gamma^{(\ell)}) \mathbf{H}_{\text{global}}^{(\ell+1)}\right)\right)
    \label{eq:san_adaptive_combination}
\end{equation}
where $\gamma^{(\ell)} \in [0,1]$ is a learnable parameter per layer.
\end{definition}

\begin{proposition}[Computational Complexity of SAN]\label{prop:san_complexity}
A SAN layer with $n$ nodes, $m$ edges, $h$ attention heads, and dimension $d$ has complexity:
\begin{equation}
    \mathcal{O}(hd(m + n^2)) = \mathcal{O}(hdm) + \mathcal{O}(hdn^2)
    \label{eq:san_complexity}
\end{equation}
The quadratic term $\mathcal{O}(hdn^2)$ dominates for dense graphs, limiting scalability to graphs with thousands of nodes.
\end{proposition}

\begin{proof}
Local attention operates over the $m$ edges of the graph: for each edge $(i,j)$, each head computes attention scores and messages in $\mathcal{O}(d)$, giving $\mathcal{O}(hdm)$. Global attention computes the matrix $\mathbf{Q}\mathbf{K}^\top \in \mathbb{R}^{n \times n}$ for each head, requiring $\mathcal{O}(dn^2)$ operations per head and $\mathcal{O}(hdn^2)$ in total. The subsequent multiplication with $\mathbf{V}$ is also $\mathcal{O}(hdn^2)$. Summing both contributions gives $\mathcal{O}(hd(m + n^2))$.
\end{proof}

\subsubsection{GraphGPS: Modular and Scalable Design}\label{subsubsec:graphgps_advanced}

GraphGPS (Graph Generalized Positional Signals) by \citet{rampasek2022graphgps} represents an evolution toward modular and composable architectures. Its design philosophy is based on three principles:

\begin{enumerate}
    \item \textbf{Separation of responsibilities}: Specialized components for local vs. global processing
    \item \textbf{Scalability}: Use of efficient attention mechanisms (Performer, BigBird)
    \item \textbf{Flexibility}: Interchangeability of MPNN and attention components
\end{enumerate}

\paragraph{Modular architecture}

\begin{definition}[GraphGPS Layer]\label{def:graphgps_layer}
A GraphGPS layer consists of two parallel paths that are merged:
\begin{equation}
    \mathbf{H}^{(\ell+1)} = \mathbf{H}^{(\ell)} + \text{MLP}^{(\ell)}\left(\text{LayerNorm}\left(\mathbf{H}_{\text{MPNN}}^{(\ell)} + \mathbf{H}_{\text{Global}}^{(\ell)}\right)\right)
    \label{eq:graphgps_main_layer}
\end{equation}

\textbf{Local path (Message Passing):}
\begin{equation}
    \mathbf{H}_{\text{MPNN}}^{(\ell)} = \text{MPNN}^{(\ell)}\left(\mathbf{H}^{(\ell)}, \mathcal{E}\right)
    \label{eq:graphgps_mpnn_path}
\end{equation}
where $\text{MPNN}^{(\ell)}$ can be \gls{GCN}, \gls{GAT}, GIN, or any message passing architecture.

\textbf{Global path (Attention with PE):}
\begin{equation}
    \mathbf{H}_{\text{Global}}^{(\ell)} = \text{Transformer}^{(\ell)}\left(\mathbf{H}^{(\ell)} + \mathbf{PE}\right)
    \label{eq:graphgps_global_path}
\end{equation}
where the \gls{Transformer} can use full attention, Performer, or BigBird for efficiency.
\end{definition}

\textbf{Interpretation as a sum of aggregations.}\label{rem:graphgps_interpretation} GraphGPS implements an additive decomposition of aggregation in the Message Passing framework (Section~\ref{subsec:message_passing_unificador}):
\begin{equation}
    m_i^{(\ell)} = \underbrace{\bigoplus_{j \in \mathcal{N}(i)} \psi_{\text{local}}(\mathbf{h}_j^{(\ell)})}_{\text{Local aggregation (MPNN)}} + \underbrace{\bigoplus_{j \in V} \alpha_{ij}^{(\ell)} \psi_{\text{global}}(\mathbf{h}_j^{(\ell)})}_{\text{Global aggregation (Attention)}}
    \label{eq:graphgps_message_decomposition}
\end{equation}
This decomposition allows capturing both local patterns (topological structure) and long-range dependencies (global interactions).

\paragraph{Scalability mechanisms}

\begin{proposition}[Complexity of GraphGPS with efficient attention]\label{prop:graphgps_efficient_complexity}
Using Performer for global attention, GraphGPS achieves complexity:
\begin{equation}
    \mathcal{O}(hdm + ndk) \quad \text{where } k \ll n
    \label{eq:graphgps_linear_complexity}
\end{equation}
with $k$ the rank of the Performer approximation. For graphs with $m = \mathcal{O}(n)$, the total complexity is $\mathcal{O}(ndk)$, linear in the number of nodes.
\end{proposition}

\begin{proof}
The MPNN component processes the $m$ edges with $h$ heads of dimension $d$, contributing $\mathcal{O}(hdm)$. For the global component, Performer replaces the explicit computation of $\mathbf{Q}\mathbf{K}^\top \in \mathbb{R}^{n \times n}$ with the approximation $\hat{\phi}(\mathbf{Q})(\hat{\phi}(\mathbf{K})^\top \mathbf{V})$, where $\hat{\phi}: \mathbb{R}^d \to \mathbb{R}^k$. The product $\hat{\phi}(\mathbf{K})^\top \mathbf{V} \in \mathbb{R}^{k \times d}$ is computed in $\mathcal{O}(ndk)$ and then $\hat{\phi}(\mathbf{Q}) \cdot (\hat{\phi}(\mathbf{K})^\top \mathbf{V})$ also in $\mathcal{O}(ndk)$. Summing: $\mathcal{O}(hdm + ndk)$. If $m = \mathcal{O}(n)$ and $h, k$ are constant with respect to $n$, this reduces to $\mathcal{O}(ndk)$.
\end{proof}

\begin{example}[Typical GraphGPS configuration]\label{ex:graphgps_configuration}
A standard configuration uses:
\begin{itemize}
    \item \textbf{Local MPNN}: GatedGCN (8 layers)
    \item \textbf{Global attention}: Performer with $k=256$ random features
    \item \textbf{Positional encoding}: Laplacian eigenvectors (10 components) + Random Walk Structural Encoding
    \item \textbf{Model dimension}: $d=512$
    \item \textbf{Attention heads}: $h=8$
\end{itemize}
This configuration processes graphs with $n \sim 10^4$ nodes in reasonable time.
\end{example}

\subsubsection{Graph \gls{Transformer} Networks}\label{subsubsec:graph_transformers}

The Graph \glspl{Transformer} of \citet{dwivedi2021generalization} address a fundamental question: \textit{how can the topological structure of the graph be incorporated directly into the attention mechanism?}

\begin{definition}[Graph \gls{Transformer} with structural bias]\label{def:graph_transformer_bias}
A Graph \gls{Transformer} modifies standard attention by incorporating the adjacency matrix $A \in \{0,1\}^{n \times n}$:
\begin{equation}
    \text{GraphAttn}(\mathbf{Q}, \mathbf{K}, \mathbf{V}, A) = \text{softmax}\left(\frac{\mathbf{Q}\mathbf{K}^T}{\sqrt{d_k}} \odot M(A)\right)\mathbf{V}
    \label{eq:graph_transformer_masked_attention}
\end{equation}
where $M(A)$ is a mask derived from $A$ and $\odot$ is the element-wise (Hadamard) product.
\end{definition}

\paragraph{Variants of structural masks}

\begin{enumerate}
    \item \textbf{Rigid binary mask} (sparsification):
    \begin{equation}
        M(A)_{ij} = \begin{cases}
            0 & \text{if } A_{ij} = 0 \\
            1 & \text{if } A_{ij} = 1
        \end{cases}
        \label{eq:binary_mask}
    \end{equation}
    Reduces attention only to connected neighbors, recovering \gls{GAT}-like behavior.

    \item \textbf{Geodesic distance mask}:
    \begin{equation}
        M(A)_{ij} = \exp\left(-\beta \cdot d_G(i,j)\right)
        \label{eq:geodesic_mask}
    \end{equation}
    where $d_G(i,j)$ is the shortest-path distance between $i$ and $j$, and $\beta > 0$ controls the decay.

    \item \textbf{Spectral bias mask}:
    \begin{equation}
        M(A)_{ij} = 1 + \sum_{k=1}^{K} \theta_k \lambda_k \mathbf{v}_{k,i} \mathbf{v}_{k,j}
        \label{eq:spectral_bias_mask}
    \end{equation}
    where $\{\theta_k\}$ are learnable parameters and $\{\lambda_k, \mathbf{v}_k\}$ the spectrum of the Laplacian.
\end{enumerate}

\begin{theorem}[GAT-GraphTransformer Equivalence]\label{thm:gat_graphtransformer_equivalence}
A Graph \gls{Transformer} with a rigid binary mask and a single attention head is equivalent to a single-layer \gls{GAT} if the projection functions $\mathbf{W}_Q, \mathbf{W}_K$ are parameterized appropriately.
\end{theorem}

\begin{proof}[Sketch]
\gls{GAT} attention is written as:
\begin{equation}
    \alpha_{ij} = \frac{\exp(e_{ij})}{\sum_{k \in \mathcal{N}(i)} \exp(e_{ik})}, \quad e_{ij} = a^T[\mathbf{W}\mathbf{h}_i \| \mathbf{W}\mathbf{h}_j]
\end{equation}
Reparameterizing $a$ and $\mathbf{W}$ as $\mathbf{W}_Q, \mathbf{W}_K$ with an additive bias yields the dot-product form. The restriction to $\mathcal{N}(i)$ is equivalent to the binary mask.
\end{proof}

\subsubsection[Structure-Aware Transformers]{Structure-Aware \glspl{Transformer}}\label{subsubsec:structure_aware_transformers}

The Structure-Aware \glspl{Transformer} (SAT) of \citet{chen2022structureaware} adopt a radically different approach: instead of encoding structure through positions or masks, they extract \textbf{explicit subgraph representations} that are incorporated as additional features.

\begin{definition}[Local subgraph representation]\label{def:subgraph_representation_sat}
For a node $i$ and radius $r \in \mathbb{N}$, the $r$-local subgraph is defined as:
\begin{equation}
    \mathcal{G}_i^{(r)} = (V_i^{(r)}, E_i^{(r)}), \quad V_i^{(r)} = \{j \in V : d_G(i,j) \leq r\}
    \label{eq:local_subgraph}
\end{equation}
where $E_i^{(r)} = \{(u,v) \in E : u, v \in V_i^{(r)}\}$ are the induced edges.

The structural representation of node $i$ is obtained by applying a \gls{GNN} to the local subgraph:
\begin{equation}
    \mathbf{h}_i^{\text{struct}} = \text{GNN}_{\text{struct}}(\mathcal{G}_i^{(r)})
    \label{eq:structural_representation_sat}
\end{equation}
\end{definition}

\paragraph{Attention augmented with structural information}

\begin{definition}[SAT with structural features]\label{def:sat_attention}
SAT extends standard attention by concatenating node features with structural representations:
\begin{align}
    \tilde{\mathbf{h}}_i &= [\mathbf{h}_i \| \mathbf{h}_i^{\text{struct}}] \label{eq:sat_augmented_features} \\
    e_{ij}^{\text{SAT}} &= a\left(\mathbf{W}_Q \tilde{\mathbf{h}}_i, \mathbf{W}_K \tilde{\mathbf{h}}_j\right) \label{eq:sat_attention_score}
\end{align}
This formulation allows attention to depend both on local node features and on the structural patterns of its neighborhood.
\end{definition}

\textbf{Comparison with spectral methods.}\label{rem:sat_vs_spectral} While spectral methods (SAN, GraphGPS) encode global structure through eigenvectors, SAT captures explicit local structural patterns. This difference is fundamental:
\begin{itemize}
    \item \textbf{Advantage of SAT}: Captures local motifs and patterns (triangles, cliques)
    \item \textbf{Advantage of spectral methods}: Scale better and capture global properties
\end{itemize}

\subsubsection{Unifying Framework and Comparative Analysis}\label{subsubsec:unified_framework_advanced}

All advanced architectures are unified under the Generalized Message Passing framework:

\begin{proposition}[Unified representation of advanced architectures]\label{prop:unified_advanced_architectures}
Each advanced architecture instantiates the generalized MPNN framework:
\begin{equation}
    \mathbf{h}_i^{(\ell+1)} = \phi^{(\ell)}\left(\mathbf{h}_i^{(\ell)}, \bigoplus_{j \in V} w_{ij}^{(\ell)} \psi^{(\ell)}(\mathbf{h}_j^{(\ell)}, \mathbf{c}_{ij})\right)
    \label{eq:unified_mpnn_generalized}
\end{equation}
with different choices of weights $w_{ij}^{(\ell)}$ and context $\mathbf{c}_{ij}$:
\begin{itemize}
    \item \textbf{SAN}: $w_{ij} = \gamma \alpha_{ij}^{\text{local}} \mathbb{I}_{(i,j) \in E} + (1-\gamma) \alpha_{ij}^{\text{global}}$, $\mathbf{c}_{ij} = \mathbf{PE}$
    \item \textbf{GraphGPS}: $w_{ij} = \alpha_{ij}^{\text{MPNN}} + \alpha_{ij}^{\text{Transformer}}$, additive decomposition
    \item \textbf{Graph Transformer}: $w_{ij} = \alpha_{ij} \cdot M(A)_{ij}$, structural mask
    \item \textbf{SAT}: $w_{ij} = \alpha_{ij}(\mathbf{h}_i^{\text{struct}}, \mathbf{h}_j^{\text{struct}})$, attention over structural features
\end{itemize}
\end{proposition}

\begin{proof}
We verify each instantiation by identifying the functions $\phi$, $\psi$, $w_{ij}$, and $\mathbf{c}_{ij}$ of the general framework.

\textbf{SAN}: The adaptive combination of \autoref{def:san_layer} with parameter $\gamma$ separates local attention (over edges $E$) and global attention (over all pairs), giving $w_{ij} = \gamma \alpha_{ij}^{\text{local}} \mathbb{I}_{(i,j) \in E} + (1-\gamma) \alpha_{ij}^{\text{global}}$. The context is the spectral positional encoding.

\textbf{GraphGPS}: \autoref{def:graphgps_layer} sums the MPNN and Transformer outputs independently, equivalent to $w_{ij} = \alpha_{ij}^{\text{MPNN}} + \alpha_{ij}^{\text{Transformer}}$ in the unified equation.

\textbf{Graph Transformer}: The structural mask $M(A)_{ij}$ restricts attention to connected pairs or those within a topological radius, multiplying the standard attention weights.

\textbf{SAT}: The attention scores are computed over structural representations $\mathbf{h}_i^{\text{struct}}$ obtained by subgraph extraction, instead of over the direct node representations.
\end{proof}

\paragraph{Comparative table of complexities and expressiveness}

\begin{table}[h]
\centering
\small
\begin{tabular}{|l|c|c|c|c|}
\hline
\textbf{Architecture} & \textbf{Complexity} & \textbf{Expressiveness} & \textbf{Scalability} & \textbf{Information} \\
\hline
GAT & $\mathcal{O}(hdm)$ & 1-WL & High & Local \\
SAN & $\mathcal{O}(hdn^2)$ & $>$1-WL & Low & Global spectral \\
GraphGPS & $\mathcal{O}(hdm + ndk)$ & $>$1-WL & Medium-High & Local + Global \\
Graph Transformer & $\mathcal{O}(hdn^2)$ & 1-WL & Low-Medium & Topological \\
SAT & $\mathcal{O}(hdm + nr^3)$ & $>$1-WL & Medium & Local motifs \\
\hline
\end{tabular}
\caption[Comparison of advanced architectures]{Comparison of advanced architectures. $n$ = nodes, $m$ = edges, $d$ = dimension, $h$ = heads, $k$ = Performer rank, $r$ = subgraph radius.}
\label{tab:advanced_architectures_comparison}
\end{table}

\subsubsection{Fundamental Limitations and Future Directions}\label{subsubsec:limitations_future}

\paragraph{Common limitations}

The analysis of \citet{ying2021transformers} identifies cross-cutting challenges:

\begin{enumerate}
    \item \textbf{Critical dependence on positional encodings}: Without robust PEs, performance degrades significantly, especially in long-range dependency tasks.

    \item \textbf{Expressiveness-scalability trade-off}: More expressive architectures (SAN, SAT) have quadratic or cubic complexity, limiting applications to small/medium graphs.

    \item \textbf{Spectral ambiguity}: Eigenvectors are defined up to sign ($\mathbf{v}_k \sim -\mathbf{v}_k$) and up to permutations in the case of repeated eigenvalues, requiring normalization techniques (SignNet, BasisNet).

    \item \textbf{Isospectral graphs}: Non-isomorphic graphs with identical spectra are indistinguishable for pure spectral methods.
\end{enumerate}

\paragraph{Open research directions}

\begin{itemize}
    \item \textbf{Equivariant attention}: Designing attention mechanisms that preserve additional symmetries (rotations in $SO(3)$, reflections) for 3D molecular data.

    \item \textbf{Topological positional encodings}: Use of persistent homology and Betti numbers as invariant positional features.

    \item \textbf{Hierarchical attention}: Multi-scale mechanisms that process graphs at different levels of granularity (nodes, subgraphs, communities).

    \item \textbf{Connection with extremal graph theory}: Characterization of classes of graphs where different architectures achieve optimal expressiveness.
\end{itemize}

\section{Differential Geometry and Manifolds}\label{sec:differential_geometry_attention}

The extension of attention mechanisms to differentiable manifolds and non-Euclidean spaces opens new possibilities for modeling data with intrinsic geometric structure.

\subsection{Non-Euclidean Spaces}

For the foundations of differentiable manifolds and Riemannian metrics, see \autoref{ch:gdl}, \autoref{subsec:fundamentos_variedades}. In attention mechanisms over manifolds, geodesic distance replaces Euclidean distance in the computation of similarities.

\subsubsection{Hyperbolic Geometry}

Hyperbolic geometry emerges naturally in the context of hierarchical data and tree structures, where volume grows exponentially with the distance from the root.

\begin{definition}[Hyperbolic Space - Hyperboloid Model]\label{def:espacio_hiperbolico_hiperboloide}
The $n$-dimensional hyperbolic space is represented through the hyperboloid model:

\begin{equation}
    \mathbb{H}^n = \{\mathbf{x} \in \mathbb{R}^{n+1} : \langle \mathbf{x}, \mathbf{x} \rangle_{\mathcal{L}} = -1, x_0 > 0\}
    \label{eq:hyperboloid_model}
\end{equation}

where $\langle \cdot, \cdot \rangle_{\mathcal{L}}$ is the Lorentz inner product:

\begin{equation}
    \langle \mathbf{x}, \mathbf{y} \rangle_{\mathcal{L}} = -x_0y_0 + \sum_{i=1}^{n} x_iy_i
    \label{eq:lorentz_inner_product}
\end{equation}
\end{definition}

\begin{definition}[Poincaré Model]\label{def:modelo_poincare}
The Poincaré model represents hyperbolic space as the unit disk:
\begin{equation}
\mathbb{D}^n = \{\mathbf{x} \in \mathbb{R}^n : \|\mathbf{x}\| < 1\}
\end{equation}
with Riemannian metric:
\begin{equation}
g_{ij} = \frac{4}{(1-\|\mathbf{x}\|^2)^2} \delta_{ij}
\end{equation}

\textbf{Mapping between models}: The transformation from the hyperboloid to Poincaré is:
\begin{equation}
\mathbf{x} = \frac{(x_1, \ldots, x_n)}{1 + x_0}
\end{equation}
The proof that this transformation is a Riemannian isometry can be found in \autoref{thm:isometria_hiperboloide_poincare} (Appendix~\ref{ch:teoremas_complementarios}).
\end{definition}

\begin{theorem}[Fundamental Properties of Hyperbolic Space]\label{thm:propiedades_espacio_hiperbolico}
The hyperbolic space $\mathbb{H}^n$ satisfies:
\begin{enumerate}
    \item \textbf{Constant negative curvature}: $K = -1$
    \item \textbf{Exponential growth of volume}: $\text{Vol}(B_r(p)) \sim e^{(n-1)r}$ for large radius $r$
    \item \textbf{Hyperbolic Gauss-Bonnet theorem}: For geodesic triangles with interior angles $\alpha, \beta, \gamma$, the area is $\text{Area}(T) = \pi - (\alpha + \beta + \gamma) > 0$ and $\int_T K\, dA = (\alpha + \beta + \gamma) - \pi < 0$
    \item \textbf{Isometry}: The group of isometries that preserve orientation and the component $x_0 > 0$ is $SO^+(1,n)$ (the connected component of the identity of $SO(1,n)$)
\end{enumerate}
Properties (1)-(3) are direct consequences of the hyperboloid metric; for proofs, see \citep{do1976differential}.
\end{theorem}

\textbf{Implications for machine learning}:
\begin{itemize}
    \item \textbf{Efficient representation of hierarchies}: Trees are embedded naturally with minimal distortion
    \item \textbf{Expressive capacity}: Exponential growth allows representing complex structures
    \item \textbf{Geometric regularization}: Negative curvature induces implicit regularization
\end{itemize}

\begin{definition}[Fundamental Hyperbolic Operations]\label{def:operaciones_hiperbolicas}
In the Poincaré model, the basic operations are:

\textbf{Möbius addition}:
\begin{equation}
\mathbf{x} \oplus \mathbf{y} = \frac{(1+2\langle\mathbf{x},\mathbf{y}\rangle+\|\mathbf{y}\|^2)\mathbf{x} + (1-\|\mathbf{x}\|^2)\mathbf{y}}{1+2\langle\mathbf{x},\mathbf{y}\rangle+\|\mathbf{x}\|^2\|\mathbf{y}\|^2}
\end{equation}

\textbf{Möbius scalar multiplication}:
\begin{equation}
r \otimes \mathbf{x} = \tanh(r \tanh^{-1}(\|\mathbf{x}\|)) \frac{\mathbf{x}}{\|\mathbf{x}\|}
\end{equation}

\textbf{Hyperbolic exponential map}:
\begin{equation}
\exp_{\mathbf{x}}^{\mathbb{H}}(\mathbf{v}) = \mathbf{x} \oplus \left(\tanh\left(\frac{\lambda_{\mathbf{x}} \|\mathbf{v}\|}{2}\right) \frac{\mathbf{v}}{\|\mathbf{v}\|}\right)
\end{equation}
where $\lambda_{\mathbf{x}} = \frac{2}{1-\|\mathbf{x}\|^2}$ is the conformal factor.
\end{definition}

\subsection{Attention in Hyperbolic Spaces}

Hyperbolic attention extends the standard attention mechanism by exploiting the geometric structure of hyperbolic space, particularly useful for data with inherent hierarchical structures.

\subsubsection{Hyperbolic Embeddings for Knowledge Graphs}
\label{subsubsec:embeddings_hiperbolicos_grafos}

Before exploring attention architectures that operate in hyperbolic geometry, it is fundamental to understand why hyperbolic space constitutes the natural geometric framework for hierarchical data. Knowledge graphs (WordNet, biomedical ontologies) and taxonomic relations exhibit an underlying tree structure where the number of nodes grows exponentially with depth. As \autoref{thm:propiedades_espacio_hiperbolico} shows, the volume of geodesic balls in $\mathbb{H}^n$ grows as $e^{(n-1)r}$, which matches the combinatorial expansion of trees. \citet{nickel2017poincare} demonstrated empirically that embeddings in the Poincaré disk $\mathbb{D}^d$ capture lexical hierarchies with significantly less distortion than Euclidean embeddings of the same dimension.

\begin{definition}[Poincaré Embedding]\label{def:embedding_poincare}
Let $\mathcal{E}$ be a finite set of entities with hierarchical relations. A \textit{Poincaré embedding} is a map
\begin{equation}
    \Phi: \mathcal{E} \to \mathbb{D}^d = \{\mathbf{x} \in \mathbb{R}^d : \|\mathbf{x}\| < 1\}
\end{equation}
that assigns to each entity $e \in \mathcal{E}$ a point $\Phi(e) \in \mathbb{D}^d$ in the $d$-dimensional Poincaré disk, so that the hyperbolic distances $d_{\mathbb{D}}(\Phi(e_i), \Phi(e_j))$ reflect the semantic relations between entities: hierarchically related entities lie close together, and the norm $\|\Phi(e)\|$ encodes the position in the hierarchy (values close to $1$ for specific entities, close to $0$ for general entities).
\end{definition}

\begin{theorem}[Embedding of Trees in Hyperbolic Space]\label{thm:embedding_arboles_hiperbolico}
Every finite tree $T$ with $n$ vertices admits an embedding into the hyperbolic plane $\mathbb{H}^2$ with arbitrarily small distortion. More precisely, for every $\varepsilon > 0$ there exists a map $f: V(T) \to \mathbb{H}^2$ such that for every pair of vertices $u, v \in V(T)$:
\begin{equation}
    d_T(u,v) \leq d_{\mathbb{H}}(f(u), f(v)) \leq (1+\varepsilon)\, d_T(u,v)
\end{equation}
where $d_T$ is the distance in the tree \citep{sarkar2011low}.
\end{theorem}

\begin{proof}[Proof idea]
The central argument is based on the correspondence between the exponential growth of hyperbolic space and the combinatorial expansion of trees. A complete $b$-ary tree of depth $h$ has $\Theta(b^h)$ leaves. In $\mathbb{H}^2$, the area of a geodesic disk of radius $r$ is $2\pi(\cosh r - 1) \sim \pi e^r$, which grows exponentially. \citet{sarkar2011low} construct the embedding recursively: they place the root at the origin, and for each internal node with $k$ children, they distribute the children equidistantly on a hyperbolic circle centered on the parent, exploiting the fact that the perimeter of a hyperbolic circle of radius $r$ grows as $2\pi \sinh r \sim \pi e^r$, which provides enough space to separate the branches without distortion. In contrast, the polynomial growth $\Theta(r^d)$ of volume in $\mathbb{R}^d$ prevents this separation in low dimension.
\end{proof}

\begin{proposition}[Contrast with the Euclidean Case]\label{prop:cota_inferior_euclidiano}
Let $T_n$ be a complete binary tree with $n$ vertices. Every embedding $f: V(T_n) \to \mathbb{R}^d$ with distortion bounded by a constant $c \geq 1$ requires dimension $d = \Omega(\log n)$. In contrast, \autoref{thm:embedding_arboles_hiperbolico} shows that $\mathbb{H}^2$ suffices with distortion $(1+\varepsilon)$ for any $\varepsilon > 0$, independently of $n$.
\end{proposition}

\begin{proof}[Proof idea]
A complete binary tree of depth $h = \Theta(\log n)$ has $2^h$ leaves. In an embedding with distortion $c$, the leaves must be at mutual distance at least $\Omega(h)$ (being the distance in the tree between leaves of distinct subtrees), and at distance at most $O(ch)$ from the root. In $\mathbb{R}^d$, a ball of radius $R = O(ch)$ contains at most $O((R/r)^d)$ points with mutual separation $r = \Omega(h)$. To accommodate $2^h$ points one needs $(ch/h)^d \geq 2^h$, that is, $d \geq \Omega(h) = \Omega(\log n)$.
\end{proof}

\begin{definition}[Poincaré Loss Function]\label{def:perdida_poincare}
Given a set of observed hierarchical relations $\mathcal{D} = \{(u, v)\}$ where $u$ is an ancestor of $v$ in the hierarchy, the \textit{Poincaré loss function} \citep{nickel2017poincare} is defined as:
\begin{equation}
    \mathcal{L}(\boldsymbol{\theta}) = \sum_{(u,v) \in \mathcal{D}} \log \frac{e^{-d_{\mathbb{D}}(\Phi(u), \Phi(v))}}{\displaystyle\sum_{v' \in \mathcal{N}(u)} e^{-d_{\mathbb{D}}(\Phi(u), \Phi(v'))}}
    \label{eq:perdida_poincare}
\end{equation}
where $\boldsymbol{\theta} = \{\Phi(e)\}_{e \in \mathcal{E}}$ are the parameters of the embedding and $\mathcal{N}(u) = \{v\} \cup \operatorname{Neg}(u)$ includes the positive example $v$ together with a set of sampled negative examples $\operatorname{Neg}(u)$.
\end{definition}

\begin{remark}[Riemannian Optimization]\label{rem:optimizacion_riemanniana_poincare}
Minimizing $\mathcal{L}$ over the Poincaré disk requires Riemannian gradient descent (RSGD). The Riemannian gradient is obtained by rescaling the Euclidean gradient by the inverse of the metric tensor:
\begin{equation}
    \nabla_{\mathbb{D}} \mathcal{L} = \lambda_{\mathbf{x}}^{-2}\, \nabla_{\mathbb{R}} \mathcal{L}
\end{equation}
where $\lambda_{\mathbf{x}} = \frac{2}{1 - \|\mathbf{x}\|^2}$ is the conformal factor (\autoref{def:operaciones_hiperbolicas}). This rescaling is necessary because the Poincaré metric $g_{\mathbf{x}} = \lambda_{\mathbf{x}}^2 \mathbf{I}$ is not flat: a fixed Euclidean step corresponds to ever larger hyperbolic distances as $\|\mathbf{x}\| \to 1$. \autoref{thm:hvt_convergence} formalizes the convergence of this optimization scheme.
\end{remark}

The geometric justification developed in this section motivates a natural question: if the data possess intrinsically hyperbolic geometry, should the attention mechanism itself not operate in that space? The following subsections answer affirmatively, defining hyperbolic attention networks and transformers over the Poincaré disk.

\begin{figure}[ht]
    \centering
    \includegraphics[width=0.9\textwidth]{figures/hyperbolic_embedding.pdf}
    \caption{Comparison between Euclidean embedding (left) and embedding in the Poincaré disk (right) of a hierarchical tree. In $\mathbb{R}^2$, the leaves of the tree are compressed and the distances are distorted. In the Poincaré disk $\mathbb{D}^2$, the exponential growth of hyperbolic volume allows separating the branches with minimal distortion (\autoref{thm:embedding_arboles_hiperbolico}), whereas in $\mathbb{R}^d$ one needs $d = \Omega(\log n)$ (\autoref{prop:cota_inferior_euclidiano}).}
    \label{fig:hyperbolic_embedding_attention}
\end{figure}

\subsubsection{Hyperbolic Attention Networks}

The HANs of \citet{gulcehre2018hyperbolic} re-express attention in terms of hyperbolic operations, preserving the geometric properties of the space.

\begin{definition}[Hyperbolic Distance]\label{def:distancia_hiperbolica}
In the hyperboloid model, the hyperbolic distance between two points $\mathbf{x}, \mathbf{y} \in \mathbb{H}^n$ is defined as:

\begin{equation}
    d_{\mathbb{H}}(\mathbf{x}, \mathbf{y}) = \text{arccosh}(-\langle \mathbf{x}, \mathbf{y} \rangle_{\mathcal{L}})
    \label{eq:hyperbolic_distance}
\end{equation}

In the Poincaré model, the equivalent distance is:
\begin{equation}
d_{\mathbb{D}}(\mathbf{u}, \mathbf{v}) = \text{arccosh}\left(1 + 2\frac{\|\mathbf{u} - \mathbf{v}\|^2}{(1-\|\mathbf{u}\|^2)(1-\|\mathbf{v}\|^2)}\right)
\end{equation}
\end{definition}

\begin{theorem}[Properties of the Hyperbolic Distance]\label{thm:propiedades_distancia_hiperbolica}
The hyperbolic distance satisfies:
\begin{enumerate}
    \item \textbf{Metric}: $d_{\mathbb{H}}(\mathbf{x}, \mathbf{y}) \geq 0$ with equality if and only if $\mathbf{x} = \mathbf{y}$
    \item \textbf{Symmetry}: $d_{\mathbb{H}}(\mathbf{x}, \mathbf{y}) = d_{\mathbb{H}}(\mathbf{y}, \mathbf{x})$
    \item \textbf{Triangle inequality}: $d_{\mathbb{H}}(\mathbf{x}, \mathbf{z}) \leq d_{\mathbb{H}}(\mathbf{x}, \mathbf{y}) + d_{\mathbb{H}}(\mathbf{y}, \mathbf{z})$
    \item \textbf{Isometric invariance}: $d_{\mathbb{H}}(g\mathbf{x}, g\mathbf{y}) = d_{\mathbb{H}}(\mathbf{x}, \mathbf{y})$ for $g \in SO^+(1,n)$
\end{enumerate}
\end{theorem}

\begin{proof}
Properties (1)-(3) are verified directly from the distance formula. Positivity and equality result from the properties of arccosh. Symmetry is immediate from the symmetry of the Minkowski bilinear form $\langle \mathbf{x}, \mathbf{y} \rangle_M$. The triangle inequality is proven using the fact that arccosh is convex and that the geodesics of the hyperboloid are sections by planes passing through the origin. For (4), the isometries of the hyperboloid preserve the Minkowski form: $\langle g\mathbf{x}, g\mathbf{y} \rangle_M = \langle \mathbf{x}, \mathbf{y} \rangle_M$ for $g \in SO^+(1,n)$, and the distance depends only on $\langle \mathbf{x}, \mathbf{y} \rangle_M$.
\end{proof}

\begin{definition}[Generalized Hyperbolic Attention]\label{def:atencion_hiperbolica}
Hyperbolic attention is defined as:

\begin{equation}
    \text{HypAttn}(\mathbf{q}, \mathbf{K}, \mathbf{V}) = \text{softmax}(f(d_{\mathbb{H}}(\mathbf{q}, \mathbf{K}))) \odot_{\mathbb{H}} \mathbf{V}
    \label{eq:hyperbolic_attention}
\end{equation}

where:
\begin{itemize}
    \item $f: \mathbb{R} \to \mathbb{R}$ is an activation function (typically $f(x) = -\alpha x$ with $\alpha > 0$)
    \item $\odot_{\mathbb{H}}$ is the hyperbolic aggregation (generalization of the weighted sum)
    \item $\mathbf{q} \in \mathbb{H}^d$ is the hyperbolic query
    \item $\mathbf{K} = \{\mathbf{k}_i\}_{i=1}^n$ are the hyperbolic keys
    \item $\mathbf{V} = \{\mathbf{v}_i\}_{i=1}^n$ are the hyperbolic values
\end{itemize}
\end{definition}

\begin{definition}[Weighted Hyperbolic Aggregation]\label{def:agregacion_hiperbolica}
The hyperbolic aggregation of values $\{\mathbf{v}_i\}$ with weights $\{w_i\}$ is defined as:

\textbf{Einstein midpoint} (natural generalization in the Klein model):
\begin{equation}
\bigoplus_{i=1}^n w_i \mathbf{v}_i = \frac{\sum_{i=1}^n w_i \gamma(\mathbf{v}_i) \mathbf{v}_i}{\sum_{i=1}^n w_i \gamma(\mathbf{v}_i)}
\end{equation}
where $\gamma(\mathbf{v}) = \frac{1}{\sqrt{1-\|\mathbf{v}\|^2}}$ is the Lorentz factor and the denominator is a scalar sum (not a vector norm).

\textbf{Fréchet mean} (minimizer of geometric variance):
\begin{equation}
\bigoplus_{i=1}^n w_i \mathbf{v}_i = \arg\min_{\mathbf{x} \in \mathbb{H}^d} \sum_{i=1}^n w_i d_{\mathbb{H}}^2(\mathbf{x}, \mathbf{v}_i)
\end{equation}
\end{definition}

\textbf{Algorithm to compute the Fréchet mean}:

\begin{algorithm}[H]
\caption{Hyperbolic Fréchet Mean}
\begin{algorithmic}[1]
    \State \textbf{Input:} Points $\{\mathbf{v}_i\}_{i=1}^n \subset \mathbb{H}^d$, weights $\{w_i\}_{i=1}^n$
    \State Initialize $\mathbf{m}^{(0)} = \mathbf{v}_1$
    \For{$t = 0, 1, 2, \ldots$ until convergence}
        \State $\mathbf{g}^{(t)} = \sum_{i=1}^n w_i \log_{\mathbf{m}^{(t)}}^{\mathbb{H}}(\mathbf{v}_i)$
        \State $\mathbf{m}^{(t+1)} = \exp_{\mathbf{m}^{(t)}}^{\mathbb{H}}(\eta \mathbf{g}^{(t)})$
    \EndFor
    \State \textbf{Output:} $\mathbf{m}^{(\infty)}$
\end{algorithmic}
\end{algorithm}

\textbf{Theoretical properties of hyperbolic attention}:

\begin{proposition}[Equivariance under Isometries]\label{prop:equivarianza_atencion_hiperbolica}
For any isometry $g \in SO^+(1,n)$ of the hyperbolic space, if the hyperbolic aggregation $\odot_{\mathbb{H}}$ is equivariant under isometries (such as the Fréchet mean):
\begin{equation}
\text{HypAttn}(g\mathbf{q}, g\mathbf{K}, g\mathbf{V}) = g \cdot \text{HypAttn}(\mathbf{q}, \mathbf{K}, \mathbf{V})
\end{equation}
\end{proposition}

\begin{proof}[Proof idea]
The attention weights depend on hyperbolic distances: $\alpha_{ij} \propto \exp(f(d_{\mathbb{H}}(\mathbf{q}, \mathbf{k}_j)))$. Since $g$ is an isometry, $d_{\mathbb{H}}(g\mathbf{q}, g\mathbf{k}_j) = d_{\mathbb{H}}(\mathbf{q}, \mathbf{k}_j)$ (\autoref{thm:propiedades_distancia_hiperbolica}, property 4), so the attention weights are invariant: $\alpha_{ij}(g\mathbf{q}, g\mathbf{K}) = \alpha_{ij}(\mathbf{q}, \mathbf{K})$. For the Fréchet mean, equivariance follows from the fact that isometries preserve the objective function $\sum_i w_i d_{\mathbb{H}}^2(g\mathbf{x}, g\mathbf{v}_i) = \sum_i w_i d_{\mathbb{H}}^2(\mathbf{x}, \mathbf{v}_i)$, so $g$ maps the minimizer to $g$ of the minimizer.
\end{proof}

\begin{corollary}[Similarity Invariance]\label{cor:invarianza_similaridad}
Hyperbolic attention is invariant under transformations that preserve the hierarchical relations in the data.
\end{corollary}

\begin{proof}
The transformations that preserve hierarchical relations are precisely the isometries of hyperbolic space, that is, the group $SO^+(1,n)$: these preserve the hyperbolic distance $d_{\mathbb{H}}$ and, with it, the proximity relations that encode the hierarchy. By \autoref{prop:equivarianza_atencion_hiperbolica}, hyperbolic attention is equivariant under every isometry $g \in SO^+(1,n)$. In particular, the attention weights $\alpha_{ij} \propto \exp(-d_{\mathbb{H}}(\mathbf{q}_i, \mathbf{k}_j))$ are invariant under isometries (since $d_{\mathbb{H}}(g\mathbf{q}_i, g\mathbf{k}_j) = d_{\mathbb{H}}(\mathbf{q}_i, \mathbf{k}_j)$), which implies that the similarities computed by the attention mechanism are invariant under transformations that preserve the hierarchical structure.
\end{proof}

\textbf{Specific computational advantages}:

\begin{itemize}
    \item \textbf{Compact representation}: Complex hierarchies require smaller dimensions than Euclidean embeddings
    \item \textbf{Implicit regularization}: Hyperbolic geometry prevents overfitting on hierarchical data
    \item \textbf{Interpretability}: The hyperbolic distance directly reflects hierarchical relations
    \item \textbf{Scalability}: Efficient operations in the Poincaré model
\end{itemize}

\subsubsection{Hyperbolic Vision \gls{Transformer}}

The Hyperbolic Vision Transformer (HVT) of \citet{hvt2024hyperbolic} represents a fundamental extension of the \gls{ViT} architecture that integrates hyperbolic geometry to efficiently capture hierarchical structures in visual data.

\textbf{Fundamental architecture of HVT}:

The HVT architecture consists of three main components that operate in hyperbolic space:

\begin{definition}[Hyperbolic Projections in HVT]\label{def:hvt_projections}
The query, key, value projections in HVT are reformulated through exponential/logarithmic maps:

\begin{align}
    \mathbf{Q} &= \text{exp}_{\mathbf{0}}^{\mathbb{H}}(\mathbf{W}_Q \text{log}_{\mathbf{0}}^{\mathbb{H}}(\mathbf{X})) \\
    \mathbf{K} &= \text{exp}_{\mathbf{0}}^{\mathbb{H}}(\mathbf{W}_K \text{log}_{\mathbf{0}}^{\mathbb{H}}(\mathbf{X})) \\
    \mathbf{V} &= \text{exp}_{\mathbf{0}}^{\mathbb{H}}(\mathbf{W}_V \text{log}_{\mathbf{0}}^{\mathbb{H}}(\mathbf{X}))
    \label{eq:hyperbolic_projections}
\end{align}

where $\text{exp}_{\mathbf{0}}^{\mathbb{H}}: T_{\mathbf{0}}\mathbb{H}^d \to \mathbb{H}^d$ and $\text{log}_{\mathbf{0}}^{\mathbb{H}}: \mathbb{H}^d \to T_{\mathbf{0}}\mathbb{H}^d$ are the hyperbolic exponential and logarithmic maps at the origin.
\end{definition}

\textbf{Geometric justification.}\label{rem:hvt_geometry} The map $\text{log}_{\mathbf{0}}^{\mathbb{H}}$ projects hyperbolic representations to the Euclidean tangent space at the origin, where standard linear transformations $\mathbf{W}_Q, \mathbf{W}_K, \mathbf{W}_V$ are applied. The map $\text{exp}_{\mathbf{0}}^{\mathbb{H}}$ subsequently projects the result back to hyperbolic space, preserving geometric properties.

\textbf{Complete formulation of the hyperbolic attention mechanism}:

\begin{definition}[Hyperbolic Multi-Head Attention]\label{def:hvt_multihead}
The multi-head attention mechanism in HVT is defined as:

\begin{equation}
\text{HypMultiHead}(\mathbf{X}) = \text{exp}_{\mathbf{0}}^{\mathbb{H}}\left(\mathbf{W}_O \text{Concat}(\text{head}_1, \ldots, \text{head}_h)\right)
\label{eq:hvt_multihead}
\end{equation}

where each head $\text{head}_i$ operates in hyperbolic space:

\begin{equation}
\text{head}_i = \text{softmax}\left(-\alpha d_{\mathbb{H}}(\mathbf{Q}_i, \mathbf{K}_i)\right) \odot_{\mathbb{H}} \mathbf{V}_i
\label{eq:hvt_head}
\end{equation}

with $\alpha > 0$ a scale parameter and $\odot_{\mathbb{H}}$ the hyperbolic aggregation (Fréchet mean).
\end{definition}

\textbf{Training procedures and numerical stability}:

Training HVT requires special considerations to maintain numerical stability in hyperbolic space:

\begin{enumerate}
    \item \textbf{Hyperbolic initialization}: The parameters are initialized through exponential projection of small Gaussian distributions:
    \begin{equation}
    \mathbf{w}^{(0)} = \text{exp}_{\mathbf{0}}^{\mathbb{H}}(\epsilon \mathbf{z}), \quad \mathbf{z} \sim \mathcal{N}(0, \mathbf{I}), \quad \epsilon \ll 1
    \end{equation}

    \item \textbf{Projection to the Poincaré ball}: To maintain numerical stability, all points are projected to $\|\mathbf{x}\| < 1 - \delta$ with $\delta = 10^{-5}$:
    \begin{equation}
    \text{proj}(\mathbf{x}) = \begin{cases}
    \mathbf{x} & \text{if } \|\mathbf{x}\| \leq 1 - \delta \\
    \frac{(1-\delta)\mathbf{x}}{\|\mathbf{x}\|} & \text{otherwise}
    \end{cases}
    \end{equation}

    \item \textbf{Riemannian Optimization}: The RSGD optimizer (Riemannian Stochastic Gradient Descent) is used:
    \begin{equation}
    \mathbf{x}^{(t+1)} = \text{exp}_{\mathbf{x}^{(t)}}^{\mathbb{H}}\left(-\eta \text{grad}^{\mathbb{H}} L(\mathbf{x}^{(t)})\right)
    \end{equation}
    where $\text{grad}^{\mathbb{H}} L = \lambda_{\mathbf{x}}^{-2} \nabla L$ is the Riemannian gradient with $\lambda_{\mathbf{x}} = \frac{2}{1-\|\mathbf{x}\|^2}$.
\end{enumerate}

\begin{theorem}[Convergence of RSGD on Manifolds]\label{thm:hvt_convergence}
Let $L: \mathbb{H}^d \to \mathbb{R}$ be a geodesically convex function with $M$-Lipschitz Riemannian gradient. The RSGD algorithm with learning rate $\eta < 1/M$ satisfies, for the averaged iterate $\bar{\mathbf{x}}^{(T)} = \frac{1}{T}\sum_{t=0}^{T-1} \mathbf{x}^{(t)}$ (average in the tangent space):
\begin{equation}
\mathbb{E}[L(\bar{\mathbf{x}}^{(T)})] - L(\mathbf{x}^*) \leq \frac{d_{\mathbb{H}}^2(\mathbf{x}^{(0)}, \mathbf{x}^*)}{2\eta T} + \frac{\eta M \sigma^2}{2}
\end{equation}
where $\sigma^2$ is the variance of the stochastic gradient and $\mathbf{x}^*$ is a global minimizer. This result is an adaptation of the convergence theorem of \citet{bonnabel2013stochastic} for stochastic gradient descent on Riemannian manifolds.
\end{theorem}

\textbf{Experimental results and ablation}:

The experiments of \citet{hvt2024hyperbolic} demonstrate significant advantages of HVT over standard \gls{ViT}:

\begin{itemize}
    \item \textbf{ImageNet-1K}: HVT-Base achieves 84.2\% top-1 accuracy with 30\% fewer parameters than ViT-Base
    \item \textbf{CIFAR-100}: Improvement of 2.8\% in accuracy with hyperbolic embeddings of dimension $d=256$ vs. Euclidean of $d=512$
    \item \textbf{Hierarchical data}: In biological species classification (taxonomic structure), HVT surpasses \gls{ViT} by 5.3\%
\end{itemize}

\subsubsection{HyPE-GT: Hyperbolic Positional Encodings}

HyPE-GT (Hyperbolic Positional Encodings for Graph Transformers) of \citet{hypegt2023} introduces learnable positional encodings in hyperbolic space that naturally capture hierarchical structures in graphs.

\textbf{Mathematical formulation of HyPE}:

\begin{definition}[Hyperbolic Positional Encodings]\label{def:hype_encoding}
For a graph $\mathcal{G} = (V, E)$, the hyperbolic positional encodings are defined as:

\begin{equation}
    \mathbf{PE}_i^{\mathbb{H}} = \text{HypMLP}(\mathbf{h}_i^{(0)}, \mathbf{s}_i)
    \label{eq:hyperbolic_positional_encoding}
\end{equation}

where:
\begin{itemize}
    \item $\mathbf{h}_i^{(0)} \in \mathbb{R}^{d_0}$ are the initial features of node $i$
    \item $\mathbf{s}_i \in \mathbb{R}^{d_s}$ is the structural vector that encodes topological properties of the node
    \item $\text{HypMLP}: \mathbb{R}^{d_0 + d_s} \to \mathbb{H}^d$ is a multilayer neural network with hyperbolic output
\end{itemize}
\end{definition}

\begin{definition}[Node Structural Vector]\label{def:structural_vector}
The structural vector $\mathbf{s}_i$ captures local and global topological properties:

\begin{equation}
\mathbf{s}_i = [\text{deg}(i), \text{CC}(i), \text{BC}(i), \text{PR}(i), \mathbf{d}_i^{\text{ecc}}]
\end{equation}

where:
\begin{itemize}
    \item $\text{deg}(i)$ is the degree of the node
    \item $\text{CC}(i)$ is the clustering coefficient
    \item $\text{BC}(i)$ is the betweenness centrality
    \item $\text{PR}(i)$ is the PageRank
    \item $\mathbf{d}_i^{\text{ecc}} \in \mathbb{R}^k$ is the vector of eccentricities (distances to representative nodes)
\end{itemize}
\end{definition}

The HyPE encodings are integrated into the \gls{Transformer} architecture through hyperbolic combination with the Möbius operations, and attention is computed over the combined embeddings.

\begin{remark}[Dimensional advantage of hyperbolic embeddings]\label{rem:hypegt_expressiveness}
The hyperbolic positional encodings of HyPE-GT \citep{hypegt2023} benefit from the general property that hyperbolic embeddings require smaller dimensions than Euclidean ones to represent hierarchical structures. By \autoref{thm:embedding_arboles_hiperbolico}, a tree with $n$ vertices can be embedded in $\mathbb{H}^2$ with distortion $(1+\varepsilon)$, whereas Euclidean embeddings with bounded distortion require dimension $d = \Omega(\log n)$ (\autoref{prop:cota_inferior_euclidiano}). This theoretical advantage motivates the use of hyperbolic encodings in HyPE-GT, although \citet{hypegt2023} do not formally establish expressiveness bounds for their architecture, their contribution being mainly empirical.
\end{remark}

\subsection{Attention and Embeddings on the Sphere}\label{subsec:atencion_esfera}

While the previous section explores attention in spaces of negative curvature, the unit sphere $S^{d-1}$ constitutes the most relevant space of constant positive curvature for attention mechanisms. Normalizing embeddings to $S^{d-1}$ induces a direct connection between the dot product (the basis of standard attention) and the spherical geodesic distance.

\begin{proposition}[Dot Product--Cosine Similarity Equivalence on $S^{d-1}$]\label{prop:dot_product_sphere}
Let $\hat{\mathbf{q}}, \hat{\mathbf{k}} \in S^{d-1}$ be normalized embeddings, that is, $\hat{\mathbf{q}} = \mathbf{q}/\|\mathbf{q}\|$ and $\hat{\mathbf{k}} = \mathbf{k}/\|\mathbf{k}\|$. Then:
\begin{equation}
\hat{\mathbf{q}}^T \hat{\mathbf{k}} = \cos\bigl(d_{S^{d-1}}(\hat{\mathbf{q}}, \hat{\mathbf{k}})\bigr)
\label{eq:dot_geodesic_sphere}
\end{equation}
where $d_{S^{d-1}}(\hat{\mathbf{q}}, \hat{\mathbf{k}}) = \arccos(\hat{\mathbf{q}}^T \hat{\mathbf{k}})$ is the geodesic distance on $S^{d-1}$. Consequently, scaled attention over normalized embeddings
\[
\alpha_{ij} = \frac{\exp\bigl(\hat{\mathbf{q}}_i^T \hat{\mathbf{k}}_j / \tau\bigr)}{\sum_l \exp\bigl(\hat{\mathbf{q}}_i^T \hat{\mathbf{k}}_l / \tau\bigr)}
\]
operates exclusively with spherical geodesic distances (cf.\ \autoref{def:dot_product_similarity}).
\end{proposition}

\begin{proof}
By definition of the metric on $S^{d-1}$ (see \autoref{def:esfera_n}), the geodesic distance between $\hat{\mathbf{q}}, \hat{\mathbf{k}} \in S^{d-1}$ is $d_{S^{d-1}}(\hat{\mathbf{q}}, \hat{\mathbf{k}}) = \arccos(\hat{\mathbf{q}}^T \hat{\mathbf{k}})$. Applying the cosine to both sides directly yields identity~\eqref{eq:dot_geodesic_sphere}. Since $\|\hat{\mathbf{q}}\| = \|\hat{\mathbf{k}}\| = 1$, the dot product $\hat{\mathbf{q}}^T \hat{\mathbf{k}}$ does not depend on magnitudes, but only on the geodesic angle, and the attention weights $\alpha_{ij}$ are a monotonically decreasing function of the geodesic distance $d_{S^{d-1}}$.
\end{proof}

\autoref{prop:dot_product_sphere} motivates the use of normalized embeddings in classification tasks: by eliminating the magnitude, the decision boundary depends exclusively on angles.

\begin{definition}[Angular Classifier]\label{def:clasificador_angular}
Let $\hat{\mathbf{f}}(\mathbf{x}) = \mathbf{f}(\mathbf{x})/\|\mathbf{f}(\mathbf{x})\| \in S^{d-1}$ be the normalized embedding of an input $\mathbf{x}$, and let $\hat{\mathbf{w}}_c = \mathbf{w}_c/\|\mathbf{w}_c\| \in S^{d-1}$ be the normalized weights for each class $c \in \{1, \ldots, C\}$. An \textit{angular classifier} \citep{wang2017normface} produces logits of the form:
\begin{equation}
z_c(\mathbf{x}) = s \cdot \hat{\mathbf{f}}(\mathbf{x})^T \hat{\mathbf{w}}_c = s \cos\theta_c
\label{eq:clasificador_angular}
\end{equation}
where $\theta_c = d_{S^{d-1}}(\hat{\mathbf{f}}(\mathbf{x}), \hat{\mathbf{w}}_c) \in [0, \pi]$ is the geodesic distance to the prototype of class $c$, and $s > 0$ is a scale parameter needed to control the concentration of the softmax distribution (without $s$, the logits would be bounded in $[-1, 1]$, producing excessively uniform probability distributions).
\end{definition}

\begin{definition}[Generalized Angular Margin Loss]\label{def:perdida_margen_angular}
Given the angular classifier structure (\autoref{def:clasificador_angular}), an \textit{angular margin loss} modifies the logit of the correct class $y$ through a margin function $g: [0, \pi] \to \mathbb{R}$:
\begin{equation}
\mathcal{L} = -\log \frac{\exp\bigl(s \cdot g(\theta_y)\bigr)}{\exp\bigl(s \cdot g(\theta_y)\bigr) + \sum_{c \neq y} \exp\bigl(s \cos\theta_c\bigr)}
\label{eq:perdida_margen_angular}
\end{equation}
The main instances of this family are:
\begin{enumerate}
    \item \textbf{NormFace} \citep{wang2017normface}: $g(\theta) = \cos\theta$ (no margin, normalized baseline).
    \item \textbf{SphereFace} \citep{liu2017sphereface}: $g(\theta) = \cos(m\theta)$ with $m \in \mathbb{N}$, $m \geq 2$ (multiplicative angular margin).
    \item \textbf{CosFace} \citep{wang2018cosface}: $g(\theta) = \cos\theta - m$ with $m > 0$ (additive margin in cosine space).
    \item \textbf{ArcFace} \citep{deng2019arcface}: $g(\theta) = \cos(\theta + m)$ with $m > 0$ (additive angular margin).
\end{enumerate}
\end{definition}

\begin{proposition}[Geodesic Interpretation of the Angular Margin]\label{prop:interpretacion_geodesica_margen}
In ArcFace, the margin $m > 0$ imposes a geodesic separation of exactly $m$ radians on $S^{d-1}$ between the embedding and the nearest decision boundary.
\end{proposition}

\begin{proof}
Without a margin, the decision boundary between classes $y$ and $c$ lies where $\cos\theta_y = \cos\theta_c$, that is, where the geodesic distances to both prototypes coincide. With ArcFace, the condition for correct classification requires $\cos(\theta_y + m) \geq \cos\theta_c$, which is equivalent to $\theta_y + m \leq \theta_c$ (in the range $[0, \pi]$ where the cosine is monotonically decreasing). At the boundary, $\theta_y + m = \theta_c$, and therefore the embedding must be at geodesic distance at least $m$ radians closer to its class prototype than to any other. Since $\theta_y = d_{S^{d-1}}(\hat{\mathbf{f}}(\mathbf{x}), \hat{\mathbf{w}}_y)$, this separation is intrinsically geodesic.
\end{proof}

\begin{remark}[Contrastive Learning as Attention on $S^{d-1}$]\label{rem:contrastive_sphere}
The InfoNCE loss used in multimodal models constitutes a particular case of angular classification. Given normalized embeddings $\hat{\mathbf{z}}_i^A, \hat{\mathbf{z}}_j^B \in S^{d-1}$ from two modalities, the loss
\[
\mathcal{L}_{\text{InfoNCE}} = -\log \frac{\exp\bigl(\hat{\mathbf{z}}_i^{A\,T} \hat{\mathbf{z}}_i^B / \tau\bigr)}{\sum_{j} \exp\bigl(\hat{\mathbf{z}}_i^{A\,T} \hat{\mathbf{z}}_j^B / \tau\bigr)}
\]
is exactly an angular classifier with $s = 1/\tau$, $g(\theta) = \cos\theta$, and the classes defined dynamically by the batch. Thus, contrastive learning performs angular classification over the sphere, where the temperature $\tau$ controls the concentration of softmax attention over the nearest geodesic neighbors.
\end{remark}

\subsection{Parallel Transport and Connections}

Parallel transport provides a fundamental mechanism for comparing tangent vectors at different points of a Riemannian manifold, preserving the intrinsic geometric structure. This concept is essential for formulating attention mechanisms on non-Euclidean manifolds.

\subsubsection{Mathematical Foundations of Parallel Transport}

We recall that the \textit{Levi-Civita connection} $\nabla$ of a Riemannian manifold $(M, g)$ is the unique affine connection compatible with the metric and torsion-free (\autoref{def:conexion_levi_civita_gdl} of \autoref{ch:gdl}), determined in local coordinates by the Christoffel symbols $\Gamma^k_{ij}$.\label{def:levi_civita_attention}\label{eq:christoffel_symbols_attention}

The \textit{parallel transport} along a curve $\gamma: [0,1] \to M$ is given by the equation $\nabla_{\dot{\gamma}(t)} V(t) = 0$, which in local coordinates translates into $\frac{dV^k}{dt} + \Gamma^k_{ij} \frac{dx^i}{dt} V^j = 0$ (see \autoref{def:transporte_paralelo_formal} of \autoref{ch:gdl} for the general case in principal bundles).\label{def:parallel_transport_attention}\label{eq:parallel_transport}\label{eq:parallel_transport_coordinates}

\begin{theorem}[Existence and Uniqueness of Parallel Transport]\label{thm:existence_parallel_transport}
Given a smooth curve $\gamma: [0,1] \to M$ and an initial vector $v_0 \in T_{\gamma(0)}M$, there exists a unique parallel vector field $V(t)$ along $\gamma$ with $V(0) = v_0$.

Parallel transport defines a linear isometry:
\begin{equation}
\Gamma_\gamma: T_{\gamma(0)}M \to T_{\gamma(1)}M
\end{equation}
that satisfies $\Gamma_\gamma(v_0) = V(1)$.
\end{theorem}

\begin{proof}
Equation~\eqref{eq:parallel_transport_coordinates} is a linear system of first-order ODEs. By the Picard-Lindelöf theorem, given an initial value $V(0) = v_0$, there exists a unique solution $V(t)$ defined for $t \in [0,1]$.

The metric preservation follows from the compatibility of $\nabla$ with $g$:
\begin{align}
\frac{d}{dt}\langle V(t), W(t) \rangle &= \langle \nabla_{\dot{\gamma}} V, W \rangle + \langle V, \nabla_{\dot{\gamma}} W \rangle \\
&= 0 \quad \text{(if $V$ and $W$ are parallel)}
\end{align}

Therefore, $\langle \Gamma_\gamma(v), \Gamma_\gamma(w) \rangle_{\gamma(1)} = \langle v, w \rangle_{\gamma(0)}$ for all $v, w \in T_{\gamma(0)}M$.
\end{proof}

\begin{proposition}[Properties of Parallel Transport]\label{prop:parallel_transport_properties}
Parallel transport $\Gamma_\gamma: T_{\gamma(0)}M \to T_{\gamma(1)}M$ satisfies:

\begin{enumerate}
    \item \textbf{Linearity}: $\Gamma_\gamma(\alpha u + \beta v) = \alpha \Gamma_\gamma(u) + \beta \Gamma_\gamma(v)$
    \item \textbf{Metric preservation}: $g_{\gamma(1)}(\Gamma_\gamma(u), \Gamma_\gamma(v)) = g_{\gamma(0)}(u, v)$
    \item \textbf{Concatenation property}: For paths $\gamma_1, \gamma_2$ with $\gamma_1(1) = \gamma_2(0)$:
    \begin{equation}
    \Gamma_{\gamma_2 * \gamma_1} = \Gamma_{\gamma_2} \circ \Gamma_{\gamma_1}
    \end{equation}
    \item \textbf{Inversion}: $\Gamma_{\gamma^{-1}} = (\Gamma_\gamma)^{-1}$ where $\gamma^{-1}(t) = \gamma(1-t)$
\end{enumerate}
\end{proposition}

\begin{proof}
\textbf{(1) Linearity}: Since parallel transport is defined as the solution of the equation $\nabla_{\dot{\gamma}} V = 0$, which is linear in $V$, the uniqueness of solutions of linear ODEs guarantees the linearity of the operator $\Gamma_\gamma$.

\textbf{(2) Metric preservation}: If $\nabla$ is the Levi-Civita connection, compatible with the metric, then for parallel fields $V, W$ along $\gamma$:
\[
\frac{d}{dt} g(V, W) = g(\nabla_{\dot{\gamma}} V, W) + g(V, \nabla_{\dot{\gamma}} W) = 0
\]
so $g(V,W)$ is constant along $\gamma$.

\textbf{(3) Concatenation}: Let $V$ be a parallel field along $\gamma_1$ with $V(0) = v$. Then $V(1) = \Gamma_{\gamma_1}(v)$. Let $W$ be parallel along $\gamma_2$ with $W(0) = V(1)$. Then $W(1) = \Gamma_{\gamma_2}(\Gamma_{\gamma_1}(v))$. But the concatenation $\gamma_2 * \gamma_1$ with field $V$ followed by $W$ is parallel along the complete path.

\textbf{(4) Inversion}: If $V$ is parallel along $\gamma$, then $V(\gamma^{-1}(t)) = V(\gamma(1-t))$ is parallel along $\gamma^{-1}$, so $\Gamma_{\gamma^{-1}} \circ \Gamma_\gamma = \text{Id}$.
\end{proof}

\begin{figure}[ht]
    \centering
    \includegraphics[width=0.8\textwidth]{figures/parallel_transport_surface.pdf}
    \caption{Parallel transport of a tangent vector along a curve $\gamma$ on a surface. The vector $V(t)$ satisfies $\nabla_{\dot{\gamma}} V = 0$, keeping its norm and relative angle to the curve constant. The dependence of the result on the chosen path (holonomy) reflects the curvature of the surface and motivates attention with parallel transport (\autoref{def:attention_parallel_transport}).}
    \label{fig:parallel_transport_attention}
\end{figure}

\subsubsection{Algorithm to Compute Geodesics and Parallel Transport}

In practice, parallel transport is computed numerically through integration of equation~\eqref{eq:parallel_transport_coordinates} along geodesics.

We recall that a \textit{geodesic} is a curve $\gamma$ that satisfies $\nabla_{\dot{\gamma}} \dot{\gamma} = 0$, or equivalently in local coordinates $\ddot{x}^k + \Gamma^k_{ij} \dot{x}^i \dot{x}^j = 0$ (\autoref{def:geodesica} of \autoref{ch:gdl}).\label{def:geodesic_attention}\label{eq:geodesic_equation}

\begin{algorithm}[H]
\caption{Computation of Parallel Transport on Riemannian Manifolds}
\label{alg:parallel_transport_computation}
\begin{algorithmic}[1]
    \State \textbf{Input:} Initial point $p \in M$, final point $q \in M$, vector $v \in T_p M$
    \State \textbf{Output:} Transported vector $\Gamma_{pq}(v) \in T_q M$
    \State
    \State \textbf{Step 1: Compute geodesic} $\gamma: [0,1] \to M$ with $\gamma(0) = p$, $\gamma(1) = q$
    \State Solve the geodesic equation~\eqref{eq:geodesic_equation} using fourth-order Runge-Kutta
    \State
    \State \textbf{Step 2: Initialize} parallel vector field $V(0) = v$
    \State
    \State \textbf{Step 3: Integrate} parallel transport equation~\eqref{eq:parallel_transport_coordinates}
    \For{$t = 0$ to $1$ (discretization $\Delta t = 1/N$)}
        \State Compute Christoffel symbols $\Gamma^k_{ij}(\gamma(t))$ at the current point
        \State Update $V(t + \Delta t)$ using the Runge-Kutta method:
        \State $\quad V^k(t + \Delta t) = V^k(t) - \Delta t \cdot \Gamma^k_{ij}(\gamma(t)) \dot{\gamma}^i(t) V^j(t)$
    \EndFor
    \State
    \State \textbf{Return} $\Gamma_{pq}(v) = V(1)$
\end{algorithmic}
\end{algorithm}

\subsubsection{Worked Example: Parallel Transport on the Sphere $S^2$}

We develop a complete example of parallel transport on the unit sphere, where the formulas admit closed-form expressions.

\begin{example}[Parallel Transport on $S^2$]\label{ex:parallel_transport_sphere}
Consider the unit sphere $S^2 = \{\mathbf{x} \in \mathbb{R}^3 : \|\mathbf{x}\| = 1\}$ with the metric induced from Euclidean space.

\textbf{Geodesics on $S^2$}: The geodesics are great circles. For two non-antipodal points $p, q \in S^2$, the geodesic is given by:
\begin{equation}
\gamma(t) = \frac{\sin((1-t)\theta)}{\sin\theta} p + \frac{\sin(t\theta)}{\sin\theta} q, \quad t \in [0,1]
\end{equation}
where $\theta = \arccos(\langle p, q \rangle)$ is the angle between $p$ and $q$.

\textbf{Explicit formula for parallel transport}: For a tangent vector $v \in T_p S^2$ (with $\langle v, p \rangle = 0$), the parallel transport to $T_q S^2$ along the geodesic $\gamma$ is:

\begin{equation}
\Gamma_{\gamma_{pq}}(v) = v - \frac{\langle v, q \rangle + \langle v, p \rangle}{1 + \langle p, q \rangle}(p + q)
\label{eq:parallel_transport_sphere_formula}
\end{equation}

Taking advantage of the fact that $v \perp p$ (that is, $\langle v, p \rangle = 0$), this formula simplifies to:
\begin{equation}
\Gamma_{\gamma_{pq}}(v) = v - \frac{\langle v, q \rangle}{1 + \langle p, q \rangle}(p + q)
\end{equation}

\textbf{Verification of properties}:
\begin{enumerate}
    \item \textbf{Tangency}: $\langle \Gamma_{\gamma_{pq}}(v), q \rangle = 0$ (the transported vector is tangent at $q$)
    \item \textbf{Norm preservation}: $\|\Gamma_{\gamma_{pq}}(v)\| = \|v\|$
    \item \textbf{Angle preservation}: For $v, w \in T_p S^2$:
    \begin{equation}
    \langle \Gamma_{\gamma_{pq}}(v), \Gamma_{\gamma_{pq}}(w) \rangle = \langle v, w \rangle
    \end{equation}
\end{enumerate}
\end{example}

\subsubsection{Parallel Transport in Attention Mechanisms}

Parallel transport allows formulating geometrically consistent attention mechanisms on Riemannian manifolds.

\begin{definition}[Attention with Parallel Transport]\label{def:attention_parallel_transport}
Let $M$ be a Riemannian manifold and $\{x_i\}_{i=1}^n \subset M$ a set of points (tokens) on the manifold. For features $\{f_i \in T_{x_i}M\}_{i=1}^n$, attention with parallel transport at point $x_j$ is defined as:

\begin{equation}
\text{Attn}_{\text{PT}}(x_j, \{x_i, f_i\}) = \sum_{i=1}^n \alpha_{ij} \Gamma_{\gamma_{ij}}(f_i)
\label{eq:attention_parallel_transport}
\end{equation}

where:
\begin{itemize}
    \item $\gamma_{ij}$ is the geodesic from $x_i$ to $x_j$
    \item $\Gamma_{\gamma_{ij}}: T_{x_i}M \to T_{x_j}M$ is the parallel transport
    \item $\alpha_{ij} = \text{softmax}_i(-d_M(x_i, x_j)^2)$ are the attention weights based on geodesic distance
\end{itemize}
\end{definition}

\begin{theorem}[Equivariance under Isometries]\label{thm:attention_parallel_transport_equivariance}
Let $\phi: M \to M$ be an isometry of the Riemannian manifold. Attention with parallel transport satisfies:

\begin{equation}
\text{Attn}_{\text{PT}}(\phi(x_j), \{\phi(x_i), d\phi_{x_i}(f_i)\}) = d\phi_{x_j}(\text{Attn}_{\text{PT}}(x_j, \{x_i, f_i\}))
\end{equation}

where $d\phi_x: T_x M \to T_{\phi(x)}M$ is the differential of $\phi$.
\end{theorem}

\begin{proof}
Isometries preserve geodesics and parallel transport:
\begin{enumerate}
    \item $\phi(\gamma_{ij}(t)) = \gamma_{\phi(i)\phi(j)}(t)$ (geodesics map to geodesics)
    \item $d\phi \circ \Gamma_{\gamma_{ij}} = \Gamma_{\phi(\gamma_{ij})} \circ d\phi$ (commutativity with transport)
    \item $d_M(\phi(x_i), \phi(x_j)) = d_M(x_i, x_j)$ (distance preservation)
\end{enumerate}

Applying these properties:
\begin{align}
\text{Attn}_{\text{PT}}(\phi(x_j), \{\phi(x_i), d\phi(f_i)\}) &= \sum_i \alpha_{ij} \Gamma_{\gamma_{\phi(i)\phi(j)}}(d\phi(f_i)) \\
&= \sum_i \alpha_{ij} d\phi(\Gamma_{\gamma_{ij}}(f_i)) \\
&= d\phi\left(\sum_i \alpha_{ij} \Gamma_{\gamma_{ij}}(f_i)\right) \\
&= d\phi(\text{Attn}_{\text{PT}}(x_j, \{x_i, f_i\})) \qquad \square
\end{align}
\end{proof}

\subsubsection{Applications in Geometric Deep Learning}

Parallel transport has found successful applications in attention architectures for non-Euclidean data:

\begin{itemize}
    \item \textbf{Climate models on the sphere}: Attention on meteorological data defined on $S^2$ using parallel transport to aggregate information from different locations while preserving orientations

    \item \textbf{3D mesh processing}: Convolutions on surfaces where parallel transport allows comparing tangent features at different vertices of the mesh

    \item \textbf{Neural networks in shape spaces}: Analysis of anatomical shapes represented as points on Riemannian manifolds (e.g., Stiefel manifold, Grassmannian)

    \item \textbf{Hyperbolic attention with transport}: Combination of hyperbolic geometry (for hierarchies) with parallel transport (for geometrically consistent aggregation)
\end{itemize}

\begin{example}[Practical implementation: Attention on $SO(3)$]\label{ex:attention_so3}
For orientation data represented as rotations $R_i \in SO(3)$, attention with parallel transport allows aggregating information while respecting the Lie geometry of the group.

Given a set of orientations $\{R_i\}$ with features $\{f_i \in T_{R_i}SO(3) \cong \mathfrak{so}(3)\}$:

\begin{enumerate}
    \item \textbf{Compute geodesics}: For $R_i, R_j \in SO(3)$, the geodesic is:
    \begin{equation}
    \gamma(t) = R_i \exp(t \log(R_i^{-1} R_j)), \quad t \in [0,1]
    \end{equation}

    \item \textbf{Parallel transport}: In Lie groups, parallel transport through left translation:
    \begin{equation}
    \Gamma_{\gamma_{ij}}(f_i) = \text{Ad}_{R_i^{-1}R_j}(f_i)
    \end{equation}
    where $\text{Ad}_R: \mathfrak{so}(3) \to \mathfrak{so}(3)$ is the adjoint representation.

    \item \textbf{Aggregation}: Compute weights $\alpha_{ij} \propto \exp(-\|R_i^{-1}R_j\|_F^2)$ and aggregate:
    \begin{equation}
    \text{Attn}(R_j) = \sum_i \alpha_{ij} \text{Ad}_{R_i^{-1}R_j}(f_i)
    \end{equation}
\end{enumerate}

This formulation is invariant under global rotations and preserves the group structure of $SO(3)$.
\end{example}
\section{Geometric Algebra (Clifford)}\label{sec:clifford_attention}

Geometric algebra (GA) provides a unified and extraordinarily powerful mathematical framework for representing and manipulating geometric objects, naturally extending attention mechanisms to richer algebraic structures. This formalism, originally developed by William Kingdon Clifford in the 19th century, has experienced a renaissance in the context of geometric deep learning, offering an elegant and computationally efficient alternative for handling complex symmetries.

\subsection{Mathematical Foundations of Clifford Algebras}

The central idea of Clifford algebra is to extend a vector space with an associative product (the \textit{geometric product}) that simultaneously encodes metric relations (inner product) and orientations (exterior product). This makes it possible to represent scalars, vectors, planes, and volumes as elements of a single algebraic structure, facilitating the construction of equivariant operations that act uniformly on all these objects.

\subsubsection{Definition and Basic Properties}

We recall that the Clifford algebra $\text{Cl}(p,q)$ (\autoref{def:clifford_cnn} of \autoref{ch:cnn}) is the associative algebra generated by orthonormal vectors satisfying $\mathbf{e}_i \mathbf{e}_j + \mathbf{e}_j \mathbf{e}_i = 2\eta_{ij}$, with dimension $2^{p+q}$ as a vector space.\label{def:algebra_clifford}

\begin{theorem}[Bott Periodicity for Clifford Algebras]\label{thm:clasificacion_clifford}
The classification of real Clifford algebras exhibits periodicity of period 8: $\text{Cl}(p+8, q) \cong \text{Cl}(p, q) \otimes M_{16}(\mathbb{R})$ and $\text{Cl}(p, q+8) \cong \text{Cl}(p, q) \otimes M_{16}(\mathbb{R})$. For fixed dimension $n = p + q$, the signature $p - q$ uniquely determines the pair $(p,q)$, and therefore the isomorphism type of the algebra. The structure (whether the algebra is isomorphic to a matrix algebra over $\mathbb{R}$, $\mathbb{C}$ or $\mathbb{H}$, or a direct product of two such) depends only on $p - q \pmod{8}$. For the proof, see \citep{lawson1989spin}.
\end{theorem}

\subsubsection{Multivectors and Grading}

Multivectors $M = \sum_k \langle M \rangle_k$ admit a decomposition by grades, where $\langle M \rangle_0$ is a scalar, $\langle M \rangle_1$ is a vector, $\langle M \rangle_2$ is a bivector, etc.\label{def:multivector}

\subsubsection{Fundamental Geometric Products}

The Clifford algebra admits multiple products that capture different geometric aspects:

The Clifford algebra admits three fundamental products: the geometric product $AB = A \cdot B + A \wedge B$, the inner product (contraction), and the exterior product (wedge), formally defined in \autoref{def:producto_geometrico_cnn} of \autoref{ch:cnn}.\label{def:productos_clifford}

\begin{proposition}[Properties of the Products]\label{prop:propiedades_productos_clifford}
The Clifford products satisfy:
\begin{enumerate}
    \item \textbf{Distributivity}: $(A + B) \wedge C = A \wedge C + B \wedge C$
    \item \textbf{Associativity}: $(A \wedge B) \wedge C = A \wedge (B \wedge C)$
    \item \textbf{Anticommutativity}: $A \wedge B = (-1)^{|A||B|} B \wedge A$
    \item \textbf{Nilpotency}: $A \wedge A = 0$ for vectors $A$
    \item \textbf{Grading}: $|A \wedge B| = |A| + |B|$
\end{enumerate}
\end{proposition}

\begin{proof}
The properties follow directly from the definition of the exterior product as the fully antisymmetric part of the geometric product. Distributivity is a consequence of the bilinearity of the geometric product. Associativity is verified by induction on the grades. Anticommutativity results from the fact that swapping two vectors in an exterior product introduces a factor $(-1)$, and for multivectors of grades $|A|$ and $|B|$ a total of $|A| \cdot |B|$ swaps accumulate. Nilpotency follows from anticommutativity: $A \wedge A = (-1)^{1 \cdot 1} A \wedge A = -A \wedge A$, so $A \wedge A = 0$. Grading is immediate: the exterior product of a $p$-vector with a $q$-vector produces a $(p+q)$-vector by construction.
\end{proof}

\subsubsection{Involutions and Conjugate Operations}

\begin{definition}[Clifford Involutions]\label{def:involuciones_clifford}
The principal involutions are:

\textbf{Reversion} (grade preserving, order reversing):
\begin{equation}
    \widetilde{A} = \sum_k (-1)^{k(k-1)/2} \langle A \rangle_k
\end{equation}

\textbf{Grade involution} (grade reversing):
\begin{equation}
    \overline{A} = \sum_k (-1)^k \langle A \rangle_k
\end{equation}

\textbf{Clifford conjugation} (combination):
\begin{equation}
    A^* = \overline{\widetilde{A}} = \sum_k (-1)^{k(k+1)/2} \langle A \rangle_k
\end{equation}
\end{definition}

\begin{proposition}[Clifford Norm]\label{prop:norma_clifford}
For elements $A, B$ of the Clifford group (Lipschitz group) $\Gamma(V,Q)$, that is, invertible elements expressible as products of non-isotropic vectors, the quadratic norm
\begin{equation}
    N(A) = A \widetilde{A} = \widetilde{A} A
\end{equation}
is a scalar and satisfies $N(AB) = N(A)N(B)$ (multiplicative property). For general multivectors, $A\widetilde{A}$ may not be a scalar and the multiplicative property does not hold in general.
\end{proposition}

\begin{proof}
Let $A = v_1 v_2 \cdots v_k$ be a product of non-isotropic vectors. Then $\widetilde{A} = v_k \cdots v_2 v_1$ (order reversal). We have:
\[
A\widetilde{A} = v_1 v_2 \cdots v_k v_k \cdots v_2 v_1 = v_1 v_2 \cdots v_{k-1} Q(v_k) v_{k-1} \cdots v_2 v_1 = Q(v_k) \cdot v_1 \cdots v_{k-1} v_{k-1} \cdots v_1
\]
Iterating, $N(A) = \prod_{i=1}^k Q(v_i) \in \mathbb{R}$, which is a scalar. The multiplicative property is verified: $N(AB) = AB\widetilde{AB} = AB\widetilde{B}\widetilde{A} = A \cdot N(B) \cdot \widetilde{A} = N(B) \cdot A\widetilde{A} = N(A)N(B)$, where we use that $N(B)$ is a scalar and therefore commutes with all elements.
\end{proof}

\subsection{Geometric Algebra Transformers (GATr): Fundamental Architecture}

\subsubsection{Projective Geometric Algebra}

GATr by \citet{brehmer2023geometric} uses the projective geometric algebra $\text{Cl}(3,0,1)$ (also denoted $\mathbb{G}_{3,0,1}$ or PGA), providing a unified representation of 3D geometric objects.

\begin{definition}[3D Projective Algebra]\label{def:algebra_proyectiva_3d}
The algebra $\text{Cl}(3,0,1)$ is built over a 4-dimensional vector space with basis $\{\mathbf{e}_0, \mathbf{e}_1, \mathbf{e}_2, \mathbf{e}_3\}$ where the signature $(3,0,1)$ indicates 3 positive directions and 1 degenerate:
\begin{align}
    \mathbf{e}_i^2 &= +1 \quad \text{for } i = 1,2,3 \text{ (spatial dimensions)} \\
    \mathbf{e}_0^2 &= 0 \quad \text{(degenerate/null direction)}
\end{align}
The total dimension of the algebra is $2^4 = 16$. The degenerate direction $\mathbf{e}_0$ is essential: it allows representing affine geometry (including translations), not just linear geometry.
\end{definition}

\begin{proposition}[Unified Representation of Geometric Objects]\label{prop:representacion_unificada_gatr}
In $\text{Cl}(3,0,1)$, multivectors decompose into grades $k = 0, 1, 2, 3, 4$ with dimensions $\binom{4}{k} = 1, 4, 6, 4, 1$ respectively (total: 16). The fundamental geometric objects of $\mathbb{R}^3$ admit canonical representations in the ``plane-based'' convention of PGA \citep{brehmer2023geometric}:

\textbf{Scalars} ($k=0$): $\alpha \in \mathbb{R}$

\textbf{Planes} ($k=1$): $\pi = n_1 \mathbf{e}_1 + n_2 \mathbf{e}_2 + n_3 \mathbf{e}_3 + d\, \mathbf{e}_0$

\textbf{Lines} ($k=2$): $\ell \in \langle \mathbf{e}_i \wedge \mathbf{e}_j \rangle$

\textbf{Points} ($k=3$): $P = x_1 (\mathbf{e}_0 \wedge \mathbf{e}_2 \wedge \mathbf{e}_3) + x_2 (\mathbf{e}_0 \wedge \mathbf{e}_3 \wedge \mathbf{e}_1) + x_3 (\mathbf{e}_0 \wedge \mathbf{e}_1 \wedge \mathbf{e}_2) + \mathbf{e}_1 \wedge \mathbf{e}_2 \wedge \mathbf{e}_3$

\textbf{Pseudoscalar} ($k=4$): $I = \mathbf{e}_0 \wedge \mathbf{e}_1 \wedge \mathbf{e}_2 \wedge \mathbf{e}_3$
\end{proposition}

\begin{proof}
The algebra $\text{Cl}(3,0,1)$ has basis $\{\mathbf{e}_0, \mathbf{e}_1, \mathbf{e}_2, \mathbf{e}_3\}$, hence $\dim \text{Cl}(3,0,1) = 2^4 = 16$. The decomposition by grade $k$ follows from the fact that the $k$-blades form a space of dimension $\binom{4}{k}$: $\binom{4}{0} + \binom{4}{1} + \binom{4}{2} + \binom{4}{3} + \binom{4}{4} = 1 + 4 + 6 + 4 + 1 = 16$.

For the canonical representations: a plane $n_1 x_1 + n_2 x_2 + n_3 x_3 + d = 0$ in $\mathbb{R}^3$ is encoded as a 1-vector $\pi = \sum n_i \mathbf{e}_i + d\,\mathbf{e}_0$, where the $\mathbf{e}_0$ component captures the affine offset. Lines, as the intersection of two planes, are represented by the exterior product $\pi_1 \wedge \pi_2$, which is a 2-vector. Points, as the intersection of three planes, are represented as 3-vectors $\pi_1 \wedge \pi_2 \wedge \pi_3$. The explicit formulas are obtained by expanding these exterior products in the canonical basis.
\end{proof}

\subsubsection{16-Dimensional Representation and Geometric Embedding}

\begin{definition}[GATr Representation Space]\label{def:espacio_representacion_gatr}
Each multivector in $\text{Cl}(3,0,1)$ is represented as a vector of 16 components:

\begin{equation}
    M = \alpha_0 + \sum_{i=0}^{3} \alpha_i \mathbf{e}_i + \sum_{0 \leq i < j \leq 3} \alpha_{ij} \mathbf{e}_i \wedge \mathbf{e}_j + \sum_{0 \leq i < j < k \leq 3} \alpha_{ijk} \mathbf{e}_i \wedge \mathbf{e}_j \wedge \mathbf{e}_k + \alpha_{0123} I
    \label{eq:gatr_multivector}
\end{equation}

The canonical 16-dimensional basis decomposes by grades:
\begin{align}
\text{Grade 0 (1):} \quad & \{1\} \\
\text{Grade 1 (4):} \quad & \{\mathbf{e}_0, \mathbf{e}_1, \mathbf{e}_2, \mathbf{e}_3\} \\
\text{Grade 2 (6):} \quad & \{\mathbf{e}_0 \wedge \mathbf{e}_1, \mathbf{e}_0 \wedge \mathbf{e}_2, \mathbf{e}_0 \wedge \mathbf{e}_3, \mathbf{e}_1 \wedge \mathbf{e}_2, \mathbf{e}_1 \wedge \mathbf{e}_3, \mathbf{e}_2 \wedge \mathbf{e}_3\} \\
\text{Grade 3 (4):} \quad & \{\mathbf{e}_0 \wedge \mathbf{e}_1 \wedge \mathbf{e}_2, \mathbf{e}_0 \wedge \mathbf{e}_1 \wedge \mathbf{e}_3, \mathbf{e}_0 \wedge \mathbf{e}_2 \wedge \mathbf{e}_3, \mathbf{e}_1 \wedge \mathbf{e}_2 \wedge \mathbf{e}_3\} \\
\text{Grade 4 (1):} \quad & \{I = \mathbf{e}_0 \wedge \mathbf{e}_1 \wedge \mathbf{e}_2 \wedge \mathbf{e}_3\}
\end{align}
\end{definition}

\begin{algorithm}[H]
\caption{Embedding of Geometric Objects in GATr (PGA)}
\label{alg:embedding_gatr}
\begin{algorithmic}[1]
    \State \textbf{Input:} Geometric object (point, plane, etc.)
    \Function{EmbedPoint}{$\mathbf{x} \in \mathbb{R}^3$}
        \State $P = x_1 (\mathbf{e}_0 \wedge \mathbf{e}_2 \wedge \mathbf{e}_3) + x_2 (\mathbf{e}_0 \wedge \mathbf{e}_3 \wedge \mathbf{e}_1) + x_3 (\mathbf{e}_0 \wedge \mathbf{e}_1 \wedge \mathbf{e}_2) + \mathbf{e}_1 \wedge \mathbf{e}_2 \wedge \mathbf{e}_3$
        \State \Return $P$ \Comment{Trivector, 16 components}
    \EndFunction
    \Function{EmbedPlane}{$\mathbf{n} \in \mathbb{R}^3, d \in \mathbb{R}$}
        \State $\pi = n_1 \mathbf{e}_1 + n_2 \mathbf{e}_2 + n_3 \mathbf{e}_3 + d\, \mathbf{e}_0$
        \State \Return $\pi$ \Comment{Vector (grade 1), 16 components}
    \EndFunction
    \Function{EmbedRotation}{$\mathbf{n} \in \mathbb{R}^3, \theta \in \mathbb{R}$}
        \State $B = \sin(\theta/2) (n_1 \mathbf{e}_2 \wedge \mathbf{e}_3 + n_2 \mathbf{e}_3 \wedge \mathbf{e}_1 + n_3 \mathbf{e}_1 \wedge \mathbf{e}_2)$
        \State \Return $\cos(\theta/2) + B$ \Comment{Unit rotor}
    \EndFunction
\end{algorithmic}
\end{algorithm}

\subsubsection{Equivariant Transformer Architecture}

\begin{definition}[Equivariant Linear Layer in GATr]\label{def:capa_lineal_gatr}
A linear transformation $f: \text{Cl}(3,0,1) \to \text{Cl}(3,0,1)$ is Pin(3,0,1)-equivariant if and only if it can be written as a combination of grade projections and multiplication by $\mathbf{e}_0$:

\begin{equation}
    f(M) = \sum_{k=0}^4 \alpha_k \langle M \rangle_k + \sum_{k=0}^4 \beta_k\, \mathbf{e}_0 \langle M \rangle_k
    \label{eq:linear_equivariant_gatr}
\end{equation}

where $\alpha_k, \beta_k \in \mathbb{R}$ are learnable parameters (9 parameters in total, given that $\beta_0 = 0$ for grade reasons). See Proposition~1 of \citet{brehmer2023geometric}.
\end{definition}

\begin{theorem}[E(3)-Equivariance of Operations in GATr]\label{thm:equivarianza_automatica_gatr}
The geometric product in $\text{Cl}(3,0,1)$ is Pin(3,0,1)-equivariant: if $\rho_u(x) = uxu^{-1}$ denotes the action of the sandwich product for $u \in \text{Pin}(3,0,1)$, then $\rho_u(xy) = \rho_u(x)\rho_u(y)$ for every pair of multivectors $x, y$.
\end{theorem}

\begin{proof}
$\rho_u(xy) = u(xy)u^{-1} = (uxu^{-1})(uyu^{-1}) = \rho_u(x)\rho_u(y)$, where the second equality uses the associativity of the geometric product and the property $u^{-1}u = 1$. See Appendix~A of \citet{brehmer2023geometric} for the verification that Pin(3,0,1) acts by algebra automorphisms.
\end{proof}

\begin{remark}
Not all Clifford operations are automatically Pin-equivariant. In particular, the dual operation $x \mapsto xI^{-1}$ is Spin-equivariant but not Pin-equivariant (see the discussion in \citet{brehmer2023geometric}, Section~3.3). GATr handles this distinction carefully in its architectural design.
\end{remark}

\begin{definition}[Geometric Attention Mechanism]\label{def:atencion_geometrica_gatr}
Attention in GATr uses the Clifford inner product in $\text{Cl}(3,0,1)$ as a similarity measure:

\begin{equation}
    \text{Attn}(Q, K, V) = \text{softmax}\left(\frac{Q \cdot K^T}{\sqrt{16}}\right) V
\end{equation}

where $Q \cdot K^T$ denotes the element-wise application of the Clifford inner product, and $\sqrt{16}$ is the scale factor corresponding to the dimension of the algebra.
\end{definition}

\begin{algorithm}[H]
\caption{GATr Forward Pass}
\label{alg:forward_gatr}
\begin{algorithmic}[1]
    \State \textbf{Input:} Sequence of multivectors $\{M_1, M_2, \ldots, M_n\} \subset \text{Cl}(3,0,1)^{16}$
    \For{each layer $l = 1, 2, \ldots, L$}
        \State $Q^{(l)} = \text{LinearEquivariant}(M^{(l-1)})$
        \State $K^{(l)} = \text{LinearEquivariant}(M^{(l-1)})$
        \State $V^{(l)} = \text{LinearEquivariant}(M^{(l-1)})$
        \State $A^{(l)} = \text{MultiHeadAttention}(Q^{(l)}, K^{(l)}, V^{(l)})$
        \State $M^{(l)} = \text{LayerNorm}(M^{(l-1)} + A^{(l)})$
        \State $M^{(l)} = \text{LayerNorm}(M^{(l)} + \text{FFNEquivariant}(M^{(l)}))$
    \EndFor
    \State \Return $M^{(L)}$
\end{algorithmic}
\end{algorithm}

\subsubsection{Specialized Architectural Components}

\begin{definition}[Grade-Aware Normalization]\label{def:normalizacion_grado}
Layer normalization in GATr is applied separately to each grade:

\begin{equation}
    \text{LayerNorm}(M) = \sum_{k=0}^4 \frac{\langle M \rangle_k - \mu_k}{\sqrt{\sigma_k^2 + \epsilon}} \gamma_k + \beta_k
\end{equation}

where $\mu_k$ and $\sigma_k^2$ are the mean and variance of the grade-$k$ components.
\end{definition}

\begin{definition}[Equivariant Activation Function in GATr]\label{def:activacion_equivariante_gatr}
The equivariant activation in GATr uses a scalar gating mechanism \citep{brehmer2023geometric}: given a multivector $x \in \text{Cl}(3,0,1)$, the activation is defined as:

\begin{equation}
    \text{GatedGELU}(x) = \text{GELU}(\langle x \rangle_0) \cdot x
\end{equation}

where $\langle x \rangle_0$ is the scalar component (grade 0) of the multivector. This function is Pin(3,0,1)-equivariant because the scalar $\langle x \rangle_0$ is invariant under the group action, and the scalar-multivector product preserves the grade structure.
\end{definition}

The Clifford attention framework presented extends naturally to other metric signatures. In particular, the algebra $\text{Cl}(1,3)$ over Minkowski spacetime allows constructing Lorentz-equivariant architectures (L-GATr, \citet{spinner2024lgatr}) for applications in relativistic physics, such as jet classification and the computation of scattering amplitudes.

\subsection{Theoretical Analysis and Fundamental Properties}

\subsubsection{Expressive Power of Geometric Algebra Architectures}

\begin{remark}[Universal Approximation in Clifford Algebras: Open Problem]\label{rem:aproximacion_universal_clifford}
\citet{brehmer2023geometric} conjecture that GATr might be a universal approximator for continuous $E(3)$-equivariant functions $f: (\mathbb{R}^3)^n \to \mathbb{R}^m$ on compact sets. However, the authors themselves explicitly note that ``GATr is not yet shown to be a universal approximator, which is an interesting direction for future work'' (p.~10). Proving this conjecture would require verifying that GATr's operations (equivariant linear layers, geometric product, attention) generate a dense algebra of equivariant functions, which remains an open problem.
\end{remark}

\begin{remark}[Parameter Efficiency]\label{cor:eficiencia_parametrica_clifford}
A practical advantage of GATr is that the representation in $\text{Cl}(3,0,1)$ compactly encodes (in 16 components) geometric objects that, in a direct tensor representation, would require more parameters to guarantee equivariance. However, a rigorous comparison of parameter efficiency between GATr and tensor architectures depends on the class of functions considered and has not been formally established.
\end{remark}

\subsubsection{Computational Complexity Analysis}

\begin{proposition}[Complexity of Clifford Operations]\label{prop:complejidad_clifford}
For multivectors of dimension $2^n$, the fundamental operations have complexity:

\begin{align}
    \text{Geometric product:} &\quad O(4^n) \\
    \text{Wedge product:} &\quad O(4^n) \\
    \text{Involutions:} &\quad O(2^n) \\
    \text{Grade projections:} &\quad O(2^n)
\end{align}

However, for $n \leq 4$ (as in GATr, where $n = 4$ for $\text{Cl}(3,0,1)$), these operations are highly efficient.
\end{proposition}

\begin{proof}
The geometric product of two multivectors of dimension $2^n$ requires multiplying each component of the first multivector with each component of the second, resulting in $2^n \times 2^n = 4^n$ scalar multiplications (plus lookups in the sign table). The wedge product has the same asymptotic complexity, since it is obtained by selecting the components of the geometric product that increase the grade. The involutions (reversion, conjugation) operate component by component, multiplying by $\pm 1$ depending on the grade, which requires $O(2^n)$ operations. Grade projections simply select the corresponding components, also in $O(2^n)$.
\end{proof}

\begin{algorithm}[H]
\caption{Efficient Geometric Product}
\label{alg:producto_geometrico_eficiente}
\begin{algorithmic}[1]
    \State \textbf{Input:} Multivectors $A, B \in \text{Cl}(p,q)$
    \Function{GeometricProduct}{$A, B$}
        \State $C = \text{zeros}(2^{p+q})$
        \For{$i = 0$ to $2^{p+q} - 1$}
            \For{$j = 0$ to $2^{p+q} - 1$}
                \State $(k, \text{sign}) = \text{ProductTable}[i][j]$
                \State $C[k] += \text{sign} \times A[i] \times B[j]$
            \EndFor
        \EndFor
        \State \Return $C$
    \EndFunction
\end{algorithmic}
\end{algorithm}

\subsection{Practical Implementation and Computational Considerations}

\subsubsection{Optimizations for Modern Hardware}

\begin{definition}[Massive Parallelization in GA]\label{def:paralelizacion_ga}
Geometric algebra operations are intrinsically parallelizable. For a batch of $B$ multivectors:

\begin{equation}
    \text{BatchGeometricProduct}(A_{B \times 16}, B_{B \times 16}) = C_{B \times 16}
\end{equation}

can be executed with $B \times 16^2$ parallel operations on a GPU.
\end{definition}

\begin{algorithm}[H]
\caption{GPU-Optimized Implementation of GATr}
\label{alg:gatr_gpu}
\begin{algorithmic}[1]
    \State \textbf{Input:} Batch of multivectors $\mathbf{M} \in \mathbb{R}^{B \times N \times 16}$
    \Function{GATrForwardGPU}{$\mathbf{M}$}
        \State \textbf{// Phase 1: Parallel projections}
        \State $Q, K, V = \text{ParallelLinear}(\mathbf{M})$ \Comment{$O(BN)$ threads}

        \State \textbf{// Phase 2: Attention with Clifford products}
        \State $\text{Scores} = \text{BatchGeometricProduct}(Q, K^T)$ \Comment{Vectorized}
        \State $\text{Attn} = \text{Softmax}(\text{ScalarPart}(\text{Scores}))$

        \State \textbf{// Phase 3: Equivariant aggregation}
        \State $\text{Output} = \text{BatchGeometricProduct}(\text{Attn}, V)$
        \State \Return $\text{Output}$
    \EndFunction
\end{algorithmic}
\end{algorithm}

\subsubsection{Numerical Stability and Practical Considerations}

\begin{proposition}[Numerical Stability of Involutions]\label{prop:estabilidad_numerica_involuciones}
The Clifford involutions are numerically stable under floating-point arithmetic. For rounding errors $\epsilon$:

\begin{equation}
    \|\widetilde{A + \epsilon} - \widetilde{A}\| \leq C\epsilon
\end{equation}

where $C$ is a constant independent of the magnitude of $A$.
\end{proposition}

\begin{proof}
Reversion acts on each homogeneous component of grade $k$ by multiplying it by $(-1)^{k(k-1)/2} \in \{-1, +1\}$. Therefore, $\widetilde{A + \epsilon} - \widetilde{A} = \widetilde{\epsilon}$, and $\|\widetilde{\epsilon}\| = \|\epsilon\|$ since reversion preserves the Euclidean norm (it only changes signs of components). The constant $C = 1$ is independent of $A$. The same argument applies to grade involution and Clifford conjugation, which are equally component-by-component operations with factors $\pm 1$.
\end{proof}

\textbf{Parameter Initialization.}\label{rem:inicializacion_parametros_ga} Parameter initialization in GATr must respect the grade structure. Weights are typically initialized as:

\begin{equation}
    w_k \sim \mathcal{N}\left(0, \frac{1}{\binom{n}{k}}\right)
\end{equation}

where $\binom{n}{k}$ is the number of grade-$k$ elements.

\subsection{Research Directions and Future Developments}

\subsubsection{Extensions to Higher Clifford Algebras}

\begin{definition}[Generalized GATr]\label{def:gatr_generalizado}
A generalized GATr over $\text{Cl}(p,q)$ maintains equivariance under the group $\text{O}(p,q)$ and can represent geometric objects in spaces of arbitrary signature.
\end{definition}

\begin{hypothesis}[Dimensional Scalability]\label{hyp:escalabilidad_dimensional}
There is a fundamental trade-off between the dimensionality of the Clifford algebra ($2^{p+q}$) and the expressive capacity, suggesting optimal dimensions for different classes of problems.
\end{hypothesis}

\subsubsection{Integration with Other GDL Architectures}

The future of architectures based on geometric algebra includes:

\begin{itemize}
    \item \textbf{CNN-GATr Hybrids}: Combination of geometric convolutions with Clifford attention
    \item \textbf{GATr-GNN}: Application of geometric algebra in message-passing networks
    \item \textbf{Generative Models}: Developments in equivariant VAEs and GANs based on GA
    \item \textbf{Quantum Architectures}: Extension to quantum computing leveraging the natural Clifford algebra
\end{itemize}

\subsubsection{Emerging Applications}

The future applications of GATr include:
\begin{itemize}
    \item \textbf{Advanced Robotics}: Control of manipulators with complex symmetries
    \item \textbf{Augmented Reality}: Real-time 3D scene processing
    \item \textbf{Structural Biology}: Prediction of protein structures with crystallographic symmetries
    \item \textbf{Computational Astrophysics}: Simulations of relativistic gravitational fields
\end{itemize}

\section{Algebraic Topology}

The previous sections have shown how attention mechanisms benefit from incorporating geometric structure, from group symmetries to spectral properties of graphs. The next level of abstraction consists of exploiting the \textit{topology} of the data: properties that are preserved under continuous deformations, such as the presence of cycles or cavities. The incorporation of topological concepts into attention mechanisms represents one of the most promising frontiers of \gls{GDL}, allowing the capture of higher-order structural features that go beyond local connections in graphs. This approach, grounded in classical algebraic topology, provides sophisticated mathematical tools for analyzing the global shape of data and its transformations under deep learning operations.

\subsection{Foundations of Algebraic Topology}

\subsubsection{Simplicial Complexes: General Theory}

We recall the definitions of $k$-simplex (\autoref{def:simplicio_gdl}), face of a simplex (\autoref{def:cara_simplicio_gdl}), and simplicial complex (\autoref{def:complejo_simplicial}), introduced in \autoref{ch:gdl}. The fundamental cases are:
\begin{align}
    k=0: &\quad \text{0-simplex (vertex)} \\
    k=1: &\quad \text{1-simplex (edge)} \\
    k=2: &\quad \text{2-simplex (triangle)} \\
    k=3: &\quad \text{3-simplex (tetrahedron)} \\
    k=n: &\quad \text{$n$-simplex (polytope of dimension } n\text{)}
\end{align}

\begin{theorem}[Simplicial Realization Theorem]\label{thm:realizacion_simplicial}
Every finite abstract simplicial complex admits a geometric realization in $\mathbb{R}^n$ for $n$ sufficiently large. For the proof, see \citep{munkres2000topology}.
\end{theorem}

\subsubsection{Orientations and Simplicial Chains}

\begin{definition}[Orientation of a Simplex]\label{def:orientacion_simplicio}
An orientation of a $k$-simplex $\sigma = [v_0, v_1, \ldots, v_k]$ is an equivalence class of orderings of its vertices under even permutations. We denote the positive orientation as $[v_0, v_1, \ldots, v_k]$ and the negative one as $-[v_0, v_1, \ldots, v_k]$.
\end{definition}

\textbf{Number of Orientations.}\label{rem:numero_orientaciones} A $k$-simplex with $k \geq 1$ has exactly two orientations. The 0-simplices have trivial orientation.

\begin{definition}[Space of Simplicial Chains]\label{def:espacio_cadenas}
For a simplicial complex $\mathcal{K}$, the space of $k$-chains is the free $\mathbb{Z}$-module:
\begin{equation}
    C_k(\mathcal{K}) = \bigoplus_{\sigma \in \mathcal{K}_k} \mathbb{Z} \cdot \sigma
\end{equation}
where $\mathcal{K}_k$ denotes the set of oriented $k$-simplices in $\mathcal{K}$.
\end{definition}

A typical element of $C_k(\mathcal{K})$ is written as:
\begin{equation}
    c = \sum_{\sigma \in \mathcal{K}_k} n_{\sigma} \sigma
\end{equation}
where $n_{\sigma} \in \mathbb{Z}$ are the coefficients.

\subsubsection{Boundary Operators}

The boundary operator $\partial_k: C_k(\mathcal{K}) \to C_{k-1}(\mathcal{K})$ (see \autoref{def:cadenas_frontera} of \autoref{ch:cnn}) satisfies $\partial_k[v_0, \ldots, v_k] = \sum_{i=0}^k (-1)^i [v_0, \ldots, \hat{v_i}, \ldots, v_k]$.\label{def:operador_frontera}

\begin{example}[Explicit Boundary Operators]\label{ex:operadores_frontera_explicitos}
\begin{align}
    \partial_0[v] &= 0 \\
    \partial_1[v_0, v_1] &= [v_1] - [v_0] \\
    \partial_2[v_0, v_1, v_2] &= [v_1, v_2] - [v_0, v_2] + [v_0, v_1] \\
    \partial_3[v_0, v_1, v_2, v_3] &= [v_1, v_2, v_3] - [v_0, v_2, v_3] + [v_0, v_1, v_3] - [v_0, v_1, v_2]
\end{align}
\end{example}

\begin{theorem}[Fundamental Property of the Boundary]\label{thm:frontera_frontera_cero}
For all $k \geq 1$:
\begin{equation}
    \partial_{k-1} \circ \partial_k = 0
\end{equation}
That is, the boundary operator applied twice yields zero.
\end{theorem}

\begin{proof}
For a $k$-simplex $[v_0, \ldots, v_k]$:
\begin{align}
    \partial_{k-1} \partial_k[v_0, \ldots, v_k] &= \partial_{k-1} \sum_{i=0}^k (-1)^i [v_0, \ldots, \hat{v_i}, \ldots, v_k] \\
    &= \sum_{i=0}^k (-1)^i \sum_{\substack{j=0 \\ j \neq i}}^k (-1)^{j'} [v_0, \ldots, \hat{v_j}, \ldots, \hat{v_i}, \ldots, v_k]
\end{align}
where $j' = j$ if $j < i$ and $j' = j-1$ if $j > i$. Each $(k-2)$-face appears exactly twice with opposite signs, so the sum is zero.
\end{proof}

\begin{definition}[Chain Complex]\label{def:complejo_cadenas}
The chain complex associated with $\mathcal{K}$ is the sequence:
\begin{equation}
    \cdots \to C_{k+1}(\mathcal{K}) \xrightarrow{\partial_{k+1}} C_k(\mathcal{K}) \xrightarrow{\partial_k} C_{k-1}(\mathcal{K}) \to \cdots \to C_0(\mathcal{K}) \xrightarrow{\partial_0} 0
\end{equation}
\end{definition}

\subsection{Simplicial Homology}

\subsubsection{Cycles, Boundaries, and Homology Groups}

\begin{definition}[Cycles and Boundaries]\label{def:ciclos_fronteras}
For a simplicial complex $\mathcal{K}$:
\begin{align}
    Z_k(\mathcal{K}) &= \ker(\partial_k) \quad \text{(space of } k\text{-cycles)} \\
    B_k(\mathcal{K}) &= \operatorname{im}(\partial_{k+1}) \quad \text{(space of } k\text{-boundaries)}
\end{align}
\end{definition}

By the fundamental property of the boundary, $B_k(\mathcal{K}) \subseteq Z_k(\mathcal{K})$ for all $k$.

The homology groups $H_k(\mathcal{K}) = Z_k/B_k$ and the Betti numbers $\beta_k = \dim H_k$ are defined as in \autoref{def:homologia_betti} of \autoref{ch:gdl}.\label{def:grupos_homologia}\label{def:numeros_betti}

\begin{remark}[Geometric Interpretation of Betti Numbers]\label{rem:interpretacion_betti}
The Betti numbers capture topological features:
\begin{align}
    \beta_0 &= \text{number of connected components} \\
    \beta_1 &= \text{number of independent loops} \\
    \beta_2 &= \text{number of 2-dimensional cavities} \\
    &\vdots \\
    \beta_k &= \text{number of independent } k\text{-cavities}
\end{align}
\end{remark}

\begin{example}[Homology of Basic Structures]\label{ex:homologia_basicas}
\textbf{Sphere $S^2$}: $H_0(S^2) = \mathbb{Z}$, $H_1(S^2) = 0$, $H_2(S^2) = \mathbb{Z}$

\textbf{Torus $T^2$}: $H_0(T^2) = \mathbb{Z}$, $H_1(T^2) = \mathbb{Z} \oplus \mathbb{Z}$, $H_2(T^2) = \mathbb{Z}$

\textbf{Real projective plane $\mathbb{RP}^2$}: $H_0(\mathbb{RP}^2) = \mathbb{Z}$, $H_1(\mathbb{RP}^2) = \mathbb{Z}_2$, $H_2(\mathbb{RP}^2) = 0$
\end{example}

\subsubsection{Persistent Homology}

\begin{definition}[Filtration]\label{def:filtracion}
A filtration of a simplicial complex $\mathcal{K}$ is a nested sequence:
\begin{equation}
    \emptyset = \mathcal{K}^0 \subseteq \mathcal{K}^1 \subseteq \mathcal{K}^2 \subseteq \cdots \subseteq \mathcal{K}^m = \mathcal{K}
\end{equation}
where each $\mathcal{K}^i$ is a simplicial subcomplex of $\mathcal{K}$.
\end{definition}

\begin{definition}[Persistent Homology Modules]\label{def:modulos_homologia_persistente}
For a filtration, the persistent homology modules are:
\begin{equation}
    H_k^{i,j} = \operatorname{im}(H_k(\mathcal{K}^i) \to H_k(\mathcal{K}^j)) \quad \text{for } i \leq j
\end{equation}
\end{definition}

\begin{definition}[Persistence Diagram]\label{def:diagrama_persistencia}
The persistence diagram $\text{Pers}_k(\mathcal{K})$ is the multiset of points $(b,d) \in \mathbb{R}^2$ where:
\begin{itemize}
    \item $b$ is the \textbf{birth time} of a homological feature
    \item $d$ is the \textbf{death time} of the same feature
    \item The \textbf{persistence} is $d - b$
\end{itemize}
\end{definition}

\begin{theorem}[Module Decomposition Theorem]\label{thm:descomposicion_modulos}
Every finitely generated persistence module decomposes as:
\begin{equation}
    M \cong \bigoplus_i \mathbb{k}[b_i, d_i) \oplus \bigoplus_j \mathbb{k}[b_j, \infty)
\end{equation}
where $\mathbb{k}[b,d)$ represents an interval module. This classical result of persistence theory is proved via the classification of finitely generated modules over polynomial rings, see \citep{Hatcher2002}.
\end{theorem}

\begin{algorithm}[H]
\caption{Persistent Homology Algorithm}\label{alg:homologia_persistente}
\textbf{Input}: Filtration $\mathcal{K}^0 \subseteq \mathcal{K}^1 \subseteq \cdots \subseteq \mathcal{K}^m$
\begin{enumerate}
    \item Construct boundary matrices $\partial_k^i$ for each $\mathcal{K}^i$
    \item Apply Smith reduction to diagonalize
    \item Identify (birth, death) pairs of features
    \item Construct persistence diagram
\end{enumerate}
\textbf{Output}: Persistence diagram $\text{Pers}_k(\mathcal{K})$
\end{algorithm}

\subsection{Cohomology and Discrete de Rham Theory}

\subsubsection{Cochain Complexes}

\begin{definition}[Space of Cochains]\label{def:espacio_cocadenas}
The space of $k$-cochains is the dual of the chains:
\begin{equation}
    C^k(\mathcal{K}) = \operatorname{Hom}(C_k(\mathcal{K}), \mathbb{R})
\end{equation}
\end{definition}

\begin{definition}[Coboundary Operator]\label{def:operador_coborde}
The coboundary operator $\delta^k: C^k(\mathcal{K}) \to C^{k+1}(\mathcal{K})$ is defined as:
\begin{equation}
    (\delta^k \omega)(\sigma) = \omega(\partial_{k+1} \sigma)
\end{equation}
for $\omega \in C^k(\mathcal{K})$ and $\sigma \in C_{k+1}(\mathcal{K})$.
\end{definition}

\begin{theorem}[Coboundary Property]\label{thm:coborde_coborde_cero}
$\delta^{k+1} \circ \delta^k = 0$ for all $k$.
\end{theorem}

\begin{proof}
For $\omega \in C^k(\mathcal{K})$ and $\tau \in C_{k+2}(\mathcal{K})$:
\[
(\delta^{k+1} \circ \delta^k)(\omega)(\tau) = (\delta^k \omega)(\partial_{k+2}\tau) = \omega(\partial_{k+1} \circ \partial_{k+2}(\tau)) = \omega(0) = 0
\]
where we use that $\partial_{k+1} \circ \partial_{k+2} = 0$ (fundamental property of the boundary operator).
\end{proof}

\begin{definition}[Cohomology Groups]\label{def:grupos_cohomologia}
The cohomology groups are:
\begin{equation}
    H^k(\mathcal{K}) = \frac{\ker(\delta^k)}{\operatorname{im}(\delta^{k-1})}
\end{equation}
\end{definition}

\begin{theorem}[Poincaré Duality for Complexes]\label{thm:dualidad_poincare_complejos}
For orientable simplicial complexes without boundary of dimension $n$:
\begin{equation}
    H^k(\mathcal{K}) \cong H_{n-k}(\mathcal{K})
\end{equation}
See \citep{Hatcher2002} for the proof.
\end{theorem}

\subsubsection{Discrete Hodge Theory}

\begin{definition}[Discrete Hodge Operators]\label{def:operadores_hodge_discretos}
On a simplicial complex with a discrete metric, we define:
\begin{align}
    \Delta_k &= \partial_{k+1}^* \partial_{k+1} + \partial_k \partial_k^* \quad \text{(Hodge Laplacian)} \\
    \Delta_k^{\text{up}} &= \partial_{k+1}^* \partial_{k+1} \quad \text{(up Laplacian)} \\
    \Delta_k^{\text{down}} &= \partial_k \partial_k^* \quad \text{(down Laplacian)}
\end{align}
where $\partial_k^*$ is the adjoint of $\partial_k$.
\end{definition}

\begin{theorem}[Discrete Hodge Decomposition]\label{thm:descomposicion_hodge_discreta}
The space of $k$-chains decomposes orthogonally as:
\begin{equation}
    C_k(\mathcal{K}) = B_k(\mathcal{K}) \oplus \mathcal{H}_k(\mathcal{K}) \oplus B_k^*(\mathcal{K})
\end{equation}
where $\mathcal{H}_k(\mathcal{K}) = \ker(\Delta_k)$ is the space of harmonic forms (\autoref{def:formas_harmonicas_discretas}) and $B_k^*(\mathcal{K}) = \operatorname{im}(\partial_k^*)$ are the co-boundaries. The proof is found in \autoref{thm:hodge_decomposition_gdl} of \autoref{ch:gdl}.
\end{theorem}

\begin{definition}[Discrete Harmonic Forms]\label{def:formas_harmonicas_discretas}
A $k$-chain $\omega \in C_k(\mathcal{K})$ is harmonic if:
\begin{equation}
    \Delta_k \omega = 0
\end{equation}
The space of harmonic forms is $\mathcal{H}_k(\mathcal{K}) = \ker(\Delta_k)$.
\end{definition}

\begin{theorem}[Harmonic-Homology Isomorphism]\label{thm:isomorfismo_harmonico_homologia}
For finite simplicial complexes:
\begin{equation}
    \mathcal{H}_k(\mathcal{K}) \cong H_k(\mathcal{K})
\end{equation}
\end{theorem}

\begin{proof}
By the Hodge decomposition (\autoref{thm:descomposicion_hodge_discreta}), $C_k = B_k \oplus \mathcal{H}_k \oplus B_k^*$. The cycles are $Z_k = \ker(\partial_k) = B_k \oplus \mathcal{H}_k$ (since $\partial_k$ vanishes on $B_k$ and $\mathcal{H}_k$, but acts injectively on $B_k^*$). Then $H_k = Z_k / B_k = (B_k \oplus \mathcal{H}_k) / B_k \cong \mathcal{H}_k$, where the isomorphism sends each class $[\omega]$ to its harmonic component.
\end{proof}

\subsection{Simplicial Attention Networks: Complete Formulation}

\subsubsection{Modeling Data as Simplicial Complexes}

Structured data can be naturally represented as simplicial complexes, where higher-order interactions are explicitly encoded (see \autoref{subsec:tdl_grafos} of \autoref{ch:gdl} for the foundations of the Hodge Laplacian and simplicial message passing that motivate this section).

\begin{definition}[Simplicial Dataset]\label{def:dataset_simplicial}
A simplicial dataset $\mathcal{D} = (\mathcal{K}, \mathbf{X})$ consists of:
\begin{itemize}
    \item A simplicial complex $\mathcal{K}$ that encodes the topological structure
    \item Simplex features $\mathbf{X} = \{\mathbf{x}_{\sigma}\}_{\sigma \in \mathcal{K}}$ where $\mathbf{x}_{\sigma} \in \mathbb{R}^{d_{\dim(\sigma)}}$
\end{itemize}
\end{definition}

\begin{example}[Applications of Simplicial Datasets]\label{ex:aplicaciones_datasets_simpliciales}
\leavevmode
\begin{itemize}
    \item \textbf{Social networks}: Vertices = people, edges = friendships, triangles = groups of friends
    \item \textbf{Proteins}: Atoms = vertices, bonds = edges, rings = cycles
    \item \textbf{Brain}: Neurons = vertices, synapses = edges, circuits = higher simplices
    \item \textbf{Citation graphs}: Papers = vertices, citations = edges, co-citations = triangles
\end{itemize}
\end{example}

\subsubsection{Architecture of Simplicial Attention Networks}

\begin{figure}[ht]
\centering
\includegraphics[width=0.8\textwidth]{figures/simplicial_messages.pdf}
\caption{Three types of messages in Simplicial Attention Networks. A target edge (red) receives: \emph{lower} messages from its boundary vertices (green), \emph{upper} messages from triangles of which it is a face (orange), and \emph{horizontal} messages from adjacent edges (purple).}
\label{fig:simplicial_messages}
\end{figure}

\begin{definition}[Simplicial Message]\label{def:mensaje_simplicial}
For a $k$-simplex $\sigma$, we define different types of messages:

\textbf{Lower Message}: From incident $(k-1)$-simplices:
\begin{equation}
    \mathbf{m}_{\sigma}^{\text{lower}} = \sum_{\tau \in \partial_k(\sigma)} \alpha_{\tau \to \sigma} \mathbf{W}^{\text{lower}} \mathbf{h}_{\tau}
\end{equation}

\textbf{Upper Message}: From incident $(k+1)$-simplices:
\begin{equation}
    \mathbf{m}_{\sigma}^{\text{upper}} = \sum_{\rho: \sigma \in \partial_{k+1}(\rho)} \alpha_{\rho \to \sigma} \mathbf{W}^{\text{upper}} \mathbf{h}_{\rho}
\end{equation}

\textbf{Horizontal Message}: From adjacent $k$-simplices:
\begin{equation}
    \mathbf{m}_{\sigma}^{\text{horizontal}} = \sum_{\tau \in \mathcal{N}_k(\sigma)} \alpha_{\tau \to \sigma} \mathbf{W}^{\text{horizontal}} \mathbf{h}_{\tau}
\end{equation}
\end{definition}

\begin{definition}[Simplicial Attention Weights]\label{def:pesos_atencion_simplicial}
The attention weights incorporate orientation information:
\begin{align}
    \alpha_{\tau \to \sigma} &= \frac{\exp(\text{LeakyReLU}(\mathbf{a}^T [\mathbf{h}_{\tau} \| \mathbf{h}_{\sigma} \| \mathbf{e}_{\text{orient}}(\tau,\sigma)]))}{\sum_{\rho} \exp(\text{LeakyReLU}(\mathbf{a}^T [\mathbf{h}_{\rho} \| \mathbf{h}_{\sigma} \| \mathbf{e}_{\text{orient}}(\rho,\sigma)]))}
\end{align}
where $\mathbf{e}_{\text{orient}}(\tau,\sigma)$ encodes the orientation relation between simplices.
\end{definition}

\begin{definition}[Complete Simplicial Update]\label{def:actualizacion_simplicial}
The update of a $k$-simplex integrates all message types:
\begin{align}
    \mathbf{h}_{\sigma}^{(l+1)} = \sigma\left(\mathbf{W}^{\text{self}} \mathbf{h}_{\sigma}^{(l)} + \mathbf{m}_{\sigma}^{\text{lower}} + \mathbf{m}_{\sigma}^{\text{upper}} + \mathbf{m}_{\sigma}^{\text{horizontal}}\right)
\end{align}
\end{definition}

\subsubsection{Preservation of Topological Structure}

\begin{theorem}[Simplicial Equivariance]\label{thm:equivarianza_simplicial}
Simplicial Attention Networks are equivariant under automorphisms of simplicial complexes. If $\phi: \mathcal{K} \to \mathcal{K}$ is a simplicial automorphism, then:
\begin{equation}
    \text{SAN}(\phi(\mathbf{X})) = \phi(\text{SAN}(\mathbf{X}))
\end{equation}
\end{theorem}

\begin{proof}
A simplicial automorphism $\phi$ induces a permutation $P_k$ on the set of $k$-simplices that preserves the incidence relations: $\mathbf{B}_k P_{k-1} = P_k \mathbf{B}_k$. The simplicial update (\autoref{def:actualizacion_simplicial}) depends on the boundary matrices through the Laplacians $\mathbf{L}_k^{\text{up}}$ and $\mathbf{L}_k^{\text{down}}$. Since $P_k \mathbf{L}_k P_k^T = \mathbf{L}_k$ (the Laplacians commute with simplicial automorphisms), and the attention weights $\alpha_{\sigma\tau}$ depend on the representations through functions symmetric with respect to $\phi$, it follows that $P_k \cdot \text{SAN}(\mathbf{X}) = \text{SAN}(P_k \cdot \mathbf{X})$ for each dimension $k$.
\end{proof}

\begin{remark}[Chain Structure in SANs]\label{rem:respeto_estructura_cadenas}
For the linear component of a SAN (without nonlinear activation functions), the chain map property
\begin{equation}
    \text{SAN}_{\text{lin}}(\partial_k \mathbf{c}) = \partial_k \text{SAN}_{\text{lin}}(\mathbf{c})
\end{equation}
holds if the linear layer commutes with the boundary operator $\partial_k$. However, this property is generally false for the full network with nonlinear activation functions: element-wise nonlinearities violate the chain map condition. In practice, SAN architectures are designed to approximately preserve the topological structure (processing signals of each simplicial grade separately and respecting adjacencies), but exact commutativity with $\partial_k$ is not guaranteed.
\end{remark}

\begin{remark}[Approximate Homology Preservation]\label{cor:preservacion_homologia}
If a network were an exact chain map (that is, if it commuted with $\partial_k$), then it would induce a well-defined map on homology: cycles would map to cycles and boundaries to boundaries. Real SANs, by incorporating nonlinearities, only preserve the homological structure approximately. The ``Hodge-aware'' design of these architectures (which processes the components of the Hodge decomposition separately) seeks to maintain this property as far as possible.
\end{remark}

\subsection{Hodge Laplacian Attention}

\subsubsection{Motivation from Differential Geometry}

On differentiable manifolds, the Hodge Laplacian $\Delta = d \delta + \delta d$ (where $d$ is the exterior derivative and $\delta$ its adjoint) plays a central role in the theory of differential forms. Its discrete analog provides a powerful tool for the analysis of signals on simplicial complexes.

\begin{definition}[Discrete Hodge Laplacian]\label{def:hodge_laplacian_discreto}
For a simplicial complex $\mathcal{K}$ with an inner product, the Hodge Laplacian of order $k$ is:
\begin{align}
    \mathbf{L}_k &= \mathbf{B}_{k+1} \mathbf{B}_{k+1}^T + \mathbf{B}_k^T \mathbf{B}_k \\
    &= \mathbf{L}_k^{\text{up}} + \mathbf{L}_k^{\text{down}}
\end{align}
where $\mathbf{B}_k$ is the boundary matrix of the complex.
\end{definition}

\begin{theorem}[Spectral Decomposition of the Hodge Laplacian]\label{thm:descomposicion_hodge_laplacian}
The Hodge Laplacian admits the decomposition:
\begin{equation}
    \mathbf{L}_k = \mathbf{U}_k \boldsymbol{\Lambda}_k \mathbf{U}_k^T
\end{equation}
where $\boldsymbol{\Lambda}_k = \text{diag}(\lambda_1, \lambda_2, \ldots, \lambda_{n_k})$ with $0 = \lambda_1 \leq \lambda_2 \leq \cdots \leq \lambda_{n_k}$.
\end{theorem}

\begin{proof}
The Hodge Laplacian $\mathbf{L}_k = \mathbf{B}_{k+1}\mathbf{B}_{k+1}^T + \mathbf{B}_k^T\mathbf{B}_k$ is symmetric and positive semidefinite: for any vector $\mathbf{x}$, $\mathbf{x}^T \mathbf{L}_k \mathbf{x} = \|\mathbf{B}_{k+1}^T \mathbf{x}\|^2 + \|\mathbf{B}_k \mathbf{x}\|^2 \geq 0$. By the spectral theorem for real symmetric matrices, $\mathbf{L}_k$ admits an orthonormal basis of eigenvectors with nonnegative real eigenvalues. The multiplicity of $\lambda = 0$ is $\dim(\ker(\mathbf{L}_k)) = \beta_k$, the $k$-th Betti number.
\end{proof}

\textbf{Interpretation of the Eigenvalues.}\label{rem:interpretacion_autovalores_hodge}
\begin{itemize}
    \item $\lambda_1 = 0$ with multiplicity $\beta_k$ (Betti number)
    \item The eigenvectors corresponding to $\lambda_1$ are the harmonic forms
    \item Larger eigenvalues correspond to signals with greater ``curvature''
\end{itemize}

\subsubsection{Hodge Laplacian Attention Networks}

\begin{definition}[Hodge Spectral Filter]\label{def:filtro_espectral_hodge}
A spectral filter on the simplicial complex is defined as:
\begin{equation}
    \mathbf{g}_{\theta}(\mathbf{L}_k) = \sum_{i=0}^{K-1} \theta_i \mathbf{T}_i(\tilde{\mathbf{L}}_k)
\end{equation}
where $\mathbf{T}_i$ are Chebyshev polynomials and $\tilde{\mathbf{L}}_k = \frac{2\mathbf{L}_k}{\lambda_{\max}} - \mathbf{I}$.
\end{definition}

\begin{definition}[Hodge Laplacian Attention]\label{def:hodge_laplacian_attention}
Hodge Laplacian-based attention incorporates spectral information:
\begin{align}
    \mathbf{Q}_k &= \mathbf{g}_{\theta}^Q(\mathbf{L}_k) \mathbf{H}_k \mathbf{W}^Q \\
    \mathbf{K}_k &= \mathbf{g}_{\theta}^K(\mathbf{L}_k) \mathbf{H}_k \mathbf{W}^K \\
    \mathbf{V}_k &= \mathbf{g}_{\theta}^V(\mathbf{L}_k) \mathbf{H}_k \mathbf{W}^V
\end{align}

The attention mechanism is defined as:
\begin{equation}
    \text{HodgeAttention}(\mathbf{H}_k) = \text{softmax}\left(\frac{\mathbf{Q}_k \mathbf{K}_k^T}{\sqrt{d_k}}\right) \mathbf{V}_k
\end{equation}
\end{definition}

\begin{theorem}[Frequency Interpretation of Hodge Attention]\label{thm:interpretacion_frecuencial_hodge}
Hodge Laplacian Attention acts as a filter in the spectral domain of the simplicial complex:
\begin{equation}
    \widehat{\text{HodgeAttention}}(\lambda) = g_{\theta}^Q(\lambda) g_{\theta}^K(\lambda) \hat{\mathbf{h}}(\lambda)
\end{equation}
where $\hat{\mathbf{h}}(\lambda)$ is the spectral transform of the features.
\end{theorem}

\begin{proof}
Projecting onto the eigenvector basis $\{\mathbf{u}_i\}$ of the Hodge Laplacian (\autoref{thm:descomposicion_hodge_laplacian}), the spectral filters act diagonally: $g_\theta(\mathbf{L}_k) = \mathbf{U}_k g_\theta(\boldsymbol{\Lambda}_k) \mathbf{U}_k^T$. Then $\mathbf{Q}_k = g_\theta^Q(\mathbf{L}_k) \mathbf{H}_k \mathbf{W}^Q$ has spectral transform $\hat{\mathbf{Q}}_k(\lambda_i) = g_\theta^Q(\lambda_i) \hat{\mathbf{H}}_k(\lambda_i) \mathbf{W}^Q$, and analogously for $\mathbf{K}_k$. The attention score in the spectral domain is $\hat{\mathbf{Q}}_k^T \hat{\mathbf{K}}_k(\lambda_i) \propto g_\theta^Q(\lambda_i) g_\theta^K(\lambda_i)$, which modulates the original signal $\hat{\mathbf{h}}(\lambda_i)$. The full mechanism acts as a band-pass filter whose transfer function is the product of the Q and K filters.
\end{proof}

\paragraph{Topological Spectral Attention}

The eigenvalues of the Hodge Laplacian (\autoref{thm:descomposicion_hodge_laplacian}) provide spectral features that capture topological properties. Topological spectral attention incorporates these features via a learnable spectral bias:

\begin{equation}
\operatorname{SpectralAttn}_k = \operatorname{softmax}\left(\frac{\mathbf{Q}_k \mathbf{K}_k^T}{\sqrt{d_k}} + \mathbf{W}_{\lambda} \operatorname{diag}(\boldsymbol{\lambda}^{(k)})\right) \mathbf{V}_k
\label{eq:spectral_topological_attention}
\end{equation}

\paragraph{Filtering in Hodge Components}

The Hodge decomposition allows filtering information by separating the signals into exact, harmonic, and co-exact components:

\begin{align}
\mathbf{H}_{\text{exact}} &= \mathbf{P}_{\operatorname{Im}(\delta_{k-1})} \mathbf{H} \label{eq:exact_component} \\
\mathbf{H}_{\text{harmonic}} &= \mathbf{P}_{\operatorname{Ker}(\mathcal{L}_k)} \mathbf{H} \\
\mathbf{H}_{\text{coexact}} &= \mathbf{P}_{\operatorname{Im}(\delta_k^*)} \mathbf{H}
\end{align}
where $\mathbf{P}$ denotes orthogonal projections onto each subspace.

\subsection{Applications in Network and Complex Graph Analysis}

\subsubsection{Topological Data Analysis (TDA)}

\begin{definition}[TDA Pipeline for Deep Learning]\label{def:pipeline_tda_dl}
The integrated TDA-Deep Learning pipeline consists of:
\begin{enumerate}
    \item \textbf{Complex construction}: From data, construct $\mathcal{K}$ (Vietoris-Rips, Alpha, etc.)
    \item \textbf{Persistence computation}: Obtain diagrams $\text{Pers}_k(\mathcal{K})$
    \item \textbf{Vectorization}: Convert diagrams to vector representations
    \item \textbf{Learning}: Use as features for attention models
\end{enumerate}
\end{definition}

\begin{definition}[Vector Representations of Persistence]\label{def:representaciones_persistencia}
Common methods include:

\textbf{Persistence Images} \citep{adams2017persistence}:
\begin{equation}
    \rho_{\sigma,u}(x,y) = \sum_{(b,d) \in \text{Dgm}} \omega(b,d) \cdot \frac{1}{\sigma\sqrt{2\pi}} \exp\left(-\frac{(x-b)^2 + (y-(d-b))^2}{2\sigma^2}\right)
\end{equation}

\textbf{Persistence Landscapes}:
\begin{equation}
    \lambda_k(t) = k\text{-th largest value of } \min(t-b, d-t)
\end{equation}
for points $(b,d)$ in the diagram.
\end{definition}

\begin{algorithm}
\caption{Simplicial Attention with TDA}\label{alg:sat_tda}
\textbf{Input}: Dataset $\mathcal{X}$, filtration parameters $\epsilon_1 < \epsilon_2 < \cdots < \epsilon_m$

\begin{enumerate}
    \item \textbf{Build filtration}:
    \begin{align}
        K_i &= \text{VietorisRips}(\mathcal{X}, \epsilon_i) \quad \forall i
    \end{align}

    \item \textbf{Compute persistent homology}:
    \begin{align}
        \text{Dgm}_k &= \text{PersistentHomology}(K_0 \subseteq K_1 \subseteq \cdots \subseteq K_m)
    \end{align}

    \item \textbf{Generate topological features}:
    \begin{align}
        \mathbf{f}_{\text{topo}} &= \text{PersistenceImage}(\text{Dgm}_k)
    \end{align}

    \item \textbf{Apply Simplicial Attention}:
    \begin{align}
        \mathbf{h}^{(l+1)} &= \text{SAN}([\mathbf{h}^{(l)} \| \mathbf{f}_{\text{topo}}])
    \end{align}
\end{enumerate}

\textbf{Output}: Topologically enriched representations
\end{algorithm}

\subsubsection{Specific Applications}

\begin{application}[Brain Network Analysis]\label{app:redes_cerebrales}
In connectomics, brain data are modeled as simplicial complexes where:
\begin{itemize}
    \item \textbf{Vertices}: Brain regions (ROIs)
    \item \textbf{Edges}: Pairwise functional connections
    \item \textbf{Triangles}: Order-3 functional interactions
    \item \textbf{k-simplices}: Higher-order activation patterns
\end{itemize}

Hodge Laplacian Attention allows detecting anomalous patterns in:
\begin{align}
    \mathbf{A}_{\text{brain}} &= \text{HodgeAttention}(\mathbf{H}_{\text{roi}}, \mathbf{L}_{\text{connectome}})
\end{align}
\end{application}

\begin{application}[Social Network Analysis]\label{app:redes_sociales}
To detect communities and hierarchical structures:
\begin{enumerate}
    \item Build the clique complex of the social graph
    \item Apply a filtration by connection strength
    \item Use persistence to identify stable communities
    \item Apply Simplicial Attention for user classification
\end{enumerate}
\end{application}

\begin{theorem}[Expressive Power of SANs vs GATs]\label{thm:poder_expresivo_sans}
For graphs with rich topological structure, Simplicial Attention Networks are strictly more expressive than traditional Graph Attention Networks:
\begin{equation}
    \exists G: \quad \text{SAN}(G) \not\in \text{span}\{\text{GAT}(G'): G' \text{ is a graph}\}
\end{equation}
\end{theorem}

\begin{proof}
We construct a separating example. Consider two graphs $G_1$ and $G_2$ with the same 1-skeleton graph (same nodes and edges) but different higher-order simplicial complexes: $G_1$ contains a 2-simplex (filled triangle) $[v_1, v_2, v_3]$ and $G_2$ does not (it only has the three edges). Any \gls{GAT} operates exclusively on nodes and edges, so it produces the same output for $G_1$ and $G_2$ (since the 1-skeleton is identical). However, a SAN accesses the 2-simplices: the signal on the filled triangle of $G_1$ generates ``upper''-type messages (\autoref{def:actualizacion_simplicial}) that do not exist in $G_2$. In particular, $\beta_1(G_1) \neq \beta_1(G_2)$ (the filled triangle eliminates an independent cycle), and this topological difference is detectable by the harmonic components of the order-1 Hodge Laplacian.
\end{proof}

\begin{corollary}[Computational Advantage]\label{cor:ventaja_computacional_sans}
For problems with higher-order interactions, SANs require fewer parameters than equivalent multilayer \glspl{GAT}:
\begin{equation}
    |\theta_{\text{SAN}}| = O(d \cdot k_{\max}) \ll |\theta_{\text{GAT-equiv}}| = O(d \cdot L^2)
\end{equation}
where $k_{\max}$ is the maximum dimension of simplices and $L$ is the depth of the \gls{GAT}.
\end{corollary}

\begin{proof}
A SAN directly processes $k$-simplices for each dimension $k = 0, 1, \ldots, k_{\max}$, with $O(d)$ parameters per simplicial dimension (attention and feature transformation matrices). This gives a total of $O(d \cdot k_{\max})$ parameters.

A standard \gls{GAT} only operates on edges (1-simplices), so capturing order-$k$ interactions requires propagating information across $k$ hops. To capture interactions up to order $k_{\max}$, at least $L \geq k_{\max}$ layers are needed, each with $O(d \cdot L)$ parameters (the intermediate dimensions must grow with $L$ to avoid information loss). The total is $O(d \cdot L^2)$. Since SANs directly access higher-order interactions, they avoid the quadratic accumulation of parameters.
\end{proof}

\subsection{Connections with Diffusion Models and Differential Equations}

\subsubsection{Diffusion on Simplicial Complexes}

\begin{definition}[Simplicial Diffusion Equation]\label{def:ecuacion_difusion_simplicial}
The diffusion equation on a simplicial complex is written as:
\begin{equation}
    \frac{\partial \mathbf{f}_k}{\partial t} = -\mathbf{L}_k \mathbf{f}_k
\end{equation}
where $\mathbf{f}_k(t) \in \mathbb{R}^{n_k}$ represents signals on $k$-simplices.
\end{definition}

\begin{theorem}[Solution of the Diffusion Equation]\label{thm:solucion_difusion_simplicial}
The solution of the diffusion equation is:
\begin{equation}
    \mathbf{f}_k(t) = \exp(-t\mathbf{L}_k) \mathbf{f}_k(0) = \sum_{i=1}^{n_k} \exp(-t\lambda_i) \langle \mathbf{u}_i, \mathbf{f}_k(0) \rangle \mathbf{u}_i
\end{equation}
\end{theorem}

\begin{proof}
The equation $\dot{\mathbf{f}}_k = -\mathbf{L}_k \mathbf{f}_k$ is a linear system of ODEs with solution $\mathbf{f}_k(t) = \exp(-t\mathbf{L}_k) \mathbf{f}_k(0)$, where the matrix exponential is defined by its series $\exp(A) = \sum_{n=0}^\infty A^n/n!$. Using the spectral decomposition $\mathbf{L}_k = \mathbf{U}_k \boldsymbol{\Lambda}_k \mathbf{U}_k^T$ (\autoref{thm:descomposicion_hodge_laplacian}): $\exp(-t\mathbf{L}_k) = \mathbf{U}_k \exp(-t\boldsymbol{\Lambda}_k) \mathbf{U}_k^T = \mathbf{U}_k \operatorname{diag}(e^{-t\lambda_1}, \ldots, e^{-t\lambda_{n_k}}) \mathbf{U}_k^T$. Expanding $\mathbf{f}_k(0) = \sum_i \langle \mathbf{u}_i, \mathbf{f}_k(0)\rangle \mathbf{u}_i$, the result is obtained. The harmonic components ($\lambda_i = 0$) are stationary, while the high-frequency ones decay exponentially.
\end{proof}

\begin{definition}[Neural Diffusion on Complexes]\label{def:neural_diffusion_complejos}
A neural diffusion layer on simplicial complexes:
\begin{align}
    \mathbf{H}_k^{(l+1)} &= \sigma\left(\mathbf{W}^{(l)} \exp(-\alpha^{(l)} \mathbf{L}_k) \mathbf{H}_k^{(l)}\right)
\end{align}
where $\alpha^{(l)}$ controls the diffusion scale.
\end{definition}

This framework unifies concepts from differential geometry, algebraic topology, and deep learning, providing sophisticated mathematical tools for the analysis of data with complex topological structure.

The mathematical tools developed in the preceding sections (equivariant, spectral, hyperbolic, Clifford, and topological attention) find direct applications in molecular science, computational physics, and structural biology, where the preservation of symmetries such as $E(3)$, $SE(3)$, or the Lorentz group is essential for predicting physical properties. For an overview of these application domains, see the introduction (\autoref{ch:intro}).

\subsection[Transformers on CW Complexes]{\glspl{Transformer} on CW Complexes}

Cell complexes represent a natural generalization of graphs, allowing the capture of higher-order relations through cells of different dimensions \citep{bodnar2021cell}.

\paragraph{Multi-Dimensional Attention}

In a CW complex, attention must operate over multiple types of cells:

\begin{align}
\text{CWAttention}(\mathbf{X}) &= \bigoplus_{k=0}^{d} \text{Attention}^{(k)}(\mathbf{X}_k) \label{eq:cw_attention} \\
\text{Attention}^{(k)}(\mathbf{X}_k) &= \text{softmax}\left(\frac{\mathbf{Q}_k \mathbf{K}_k^T}{\sqrt{d_k}} + \mathbf{A}^{(k)}\right) \mathbf{V}_k
\end{align}

where $\mathbf{X}_k$ represents the features of $k$-cells and $\mathbf{A}^{(k)}$ is the adjacency matrix between $k$-cells.

\paragraph{Preservation of Incidence Relations}

The incidence relations between cells of different dimensions are preserved through:

\begin{equation}
\mathbf{M}^{(k,k-1)} = \text{IncidenceMatrix}(\text{Cell}_k, \text{Cell}_{k-1})
\label{eq:incidence_preservation}
\end{equation}

Inter-dimensional attention incorporates these relations:

\begin{equation}
\text{InterAttn}_{k \to k-1} = \text{softmax}(\mathbf{Q}_{k-1} (\mathbf{M}^{(k,k-1)} \mathbf{K}_k)^T) \mathbf{M}^{(k,k-1)} \mathbf{V}_k
\label{eq:inter_dimensional_attention}
\end{equation}

\subsection{Persistent Homology Attention}

Persistent homology provides a multiscale description of topological features, capturing how these evolve across spatial scales.

\paragraph{Persistence Diagrams}

For a filtration $X_0 \subseteq X_1 \subseteq \cdots \subseteq X_n$, the persistence diagrams record the birth and death time of topological features:

\begin{equation}
\text{Dgm}_k = \{(b_i, d_i) : H_k(X_{b_i}) \to H_k(X_{d_i}) \text{ is an isomorphism}\}
\label{eq:persistence_diagram}
\end{equation}

\paragraph{Integration into Attention}

The persistence diagrams are integrated into the attention mechanism through:

\begin{equation}
\alpha_{ij} = \text{softmax}\left(\frac{\mathbf{q}_i^T \mathbf{k}_j}{\sqrt{d}} + w \cdot \text{pers}(i,j)\right)
\label{eq:persistent_attention_weights}
\end{equation}

where $\text{pers}(i,j)$ encodes the topological persistence between points $i$ and $j$.

\paragraph{Persistence Functions}

The persistence function can take various forms:

\begin{align}
\text{pers}_{\text{lifetime}}(i,j) &= \sum_{(b,d) \in \text{Dgm}} (d-b) \cdot \mathbb{I}[\text{connected}(i,j,b,d)] \label{eq:lifetime_persistence} \\
\text{pers}_{\text{birth}}(i,j) &= \min\{b : (b,d) \in \text{Dgm}, \text{connected}(i,j,b,d)\} \\
\text{pers}_{\text{stability}}(i,j) &= \max_{(b,d) \in \text{Dgm}} (d-b) \cdot \mathbb{I}[\text{connected}(i,j,b,d)]
\end{align}

\paragraph{Advantages of Persistence}

The incorporation of persistent homology offers:

\begin{enumerate}
    \item \textbf{Robustness to noise}: Persistent features are more stable
    \item \textbf{Multiscale detection}: Captures patterns at multiple resolutions
    \item \textbf{Topological invariance}: Resistance to continuous deformations
    \item \textbf{Interpretability}: Direct connection with known geometric properties
\end{enumerate}

\subsection{Applications of Topological Algebra}

This advanced formulation finds applications in:

\begin{itemize}
    \item \textbf{Analysis of flows in networks}: Separation of irrotational and solenoidal components using Hodge
    \item \textbf{Electromagnetic processing}: Electric and magnetic fields as differential forms
    \item \textbf{Computational topology}: Analysis of cavities and holes in complex data
    \item \textbf{Time series analysis}: Capture of recurrent topological patterns with persistence
    \item \textbf{Anomaly detection}: Identification of unexpected topological changes
    \item \textbf{Dynamical systems}: Characterization of attractors and orbits through topology
\end{itemize}


\subsection{Outstanding Theoretical and Practical Challenges}

Despite significant advances, geometric attention mechanisms face fundamental challenges that require ongoing research.

\subsubsection{Limitations of Expressivity vs Efficiency}

\paragraph{The Complexity Dilemma}

The trade-off between expressivity and computational efficiency remains one of the central challenges. While traditional attention achieves $O(n^2)$ complexity, geometric approaches frequently increase the complexity due to:

\begin{enumerate}
    \item \textbf{Operations on Lie groups}: The cost of projections and lifting to high-dimensional spaces
    \item \textbf{Spectral analysis}: Eigendecomposition of Laplacians on large graphs
    \item \textbf{Parallel transport}: Geodesic computations on complex manifolds
    \item \textbf{Tensor products}: Dimensional expansion in irreducible representations
\end{enumerate}

\begin{definition}[Geometric Complexity Index]
For a geometric attention architecture with symmetry $G$, we define:
\begin{equation}
\text{ICG}(G, n) = \frac{\text{Cost}(G\text{-equivariance})}{\text{Cost}(\text{standard attention})} \cdot \frac{|\text{Orbit}(G)|}{n}
\label{eq:geometric_complexity_index}
\end{equation}
where $|\text{Orbit}(G)|$ measures the size of the orbits under the group action.
\end{definition}

\paragraph{Approximations and Trade-offs}

Recent developments in sub-quadratic architectures suggest that it is possible to preserve essential geometric properties with linear complexity:

\begin{remark}[Asymptotic Preservation of Equivariance]
Under certain regularity conditions, it is conjectured that linear-complexity models can approximate equivariant architectures with bounded error:
\begin{equation}
\|f_{\text{linear}}(g \cdot x) - g \cdot f_{\text{quadratic}}(x)\| \leq \epsilon_{\text{precision}}
\label{eq:linear_equivariance_approximation}
\end{equation}
The rigorous formalization of this bound and the precise conditions under which it holds remain an open problem.
\end{remark}

\subsubsection{Scalability to Massive Data}

\paragraph{Memory Limitations}

Modern foundation models require the processing of:
\begin{itemize}
    \item \textbf{Protein structures}: $>10^6$ atoms per macromolecular complex
    \item \textbf{Molecular simulations}: $10^9$ particles in temporally extended dynamics
    \item \textbf{Social network graphs}: $10^{12}$ connections in global platforms
    \item \textbf{LiDAR point clouds}: $10^8$ points per city scan
\end{itemize}

The required GPU memory scales as:

\begin{equation}
\text{Memory}_{\text{GDL}} = O(n^2 \cdot d_{\text{model}} \cdot |\text{Rep}(G)| \cdot \text{Precision})
\label{eq:gdl_memory_scaling}
\end{equation}

where $|\text{Rep}(G)|$ represents the dimension of the group representations.

\paragraph{Distribution Strategies}

Emerging solutions include:

\begin{enumerate}
    \item \textbf{Distributed Ring Attention}: Ring parallelization with communication overlapping
    \item \textbf{Geometric gradient checkpointing}: Selective recompute of expensive equivariant operations
    \item \textbf{Symmetry-aware quantization}: Preservation of equivariance under reduced precision
    \item \textbf{Structured sparsification}: Pruning that respects the group structure
\end{enumerate}

\subsubsection{Interpretability of Geometric Attention}

\paragraph{Visualization Complexity}

The interpretability of geometric attention mechanisms presents unique challenges:

\begin{itemize}
    \item \textbf{High-dimensional spaces}: Visualization of representations in $SO(3)$ or $SE(3)$
    \item \textbf{Multiple scales}: Attention operating simultaneously at various geometric resolutions
    \item \textbf{Non-linear interactions}: Tensor products and complex group operations
    \item \textbf{Temporal dependencies}: Evolution of attention patterns in dynamic sequences
\end{itemize}

\begin{definition}[Geometric Saliency]
For a $G$-equivariant attention mechanism, the saliency with respect to the group is defined as:
\begin{equation}
\text{Sal}_G(x_i, x_j) = \sum_{g \in G} w(g) \cdot \alpha_{ij}(g \cdot x_i, g \cdot x_j)
\label{eq:geometric_saliency}
\end{equation}
where $w(g)$ are weights associated with group elements.
\end{definition}

\paragraph{Analysis Tools}

The development of specialized tools includes:

\begin{enumerate}
    \item \textbf{Manifold projection}: Visualization of attention in interpretable 2D/3D spaces
    \item \textbf{Spectral analysis of patterns}: Decomposition of attention maps into frequency components
    \item \textbf{Symmetry traceability}: Tracking the preservation of equivariance across layers
    \item \textbf{Directed perturbation}: Sensitivity analysis with respect to specific geometric transformations
\end{enumerate}

\subsubsection{Theoretical Generalization}

\paragraph{Universal Approximation Limits}

The universal approximation results for geometric attention require extension:

\begin{remark}[Group-Conditioned Universal Approximation]
Let $\mathcal{F}_G$ be the space of $G$-equivariant functions on a geometric domain $\mathcal{M}$. It is conjectured that a geometric attention network can uniformly approximate any $f \in \mathcal{F}_G$ under the following conditions:
\begin{enumerate}
    \item The action of $G$ on $\mathcal{M}$ is piecewise transitive
    \item The representations used generate all the necessary irreducible representations
    \item The attention mechanism respects the causal structure of the domain
\end{enumerate}
This result extends the universal approximation conjecture for equivariant functions (see \autoref{rem:aproximacion_universal_clifford}) and remains an open problem.
\end{remark}

\paragraph{Generalization Bounds}

The generalization capacity is limited by:

\begin{equation}
\text{Error}_{\text{gen}} \leq \text{Error}_{\text{emp}} + \sqrt{\frac{\text{Dim}(\text{Rep}(G)) \log N}{N}}
\label{eq:geometric_generalization_bound}
\end{equation}

where $\text{Dim}(\text{Rep}(G))$ is the effective dimension of the representation space.

\section{Connections with CNNs and Future Directions}\label{sec:cnn_attention_connections}

This chapter has developed geometric attention mechanisms from multiple mathematical perspectives: group theory, spectral analysis, differential geometry, geometric algebra, and algebraic topology. To complete the unifying theoretical framework of Geometric Deep Learning, this final section establishes the fundamental connections with the Geometric Convolutional Neural Networks developed in \autoref{ch:cnn}, demonstrating the conceptual convergence of both paradigms under shared mathematical principles.

\subsection{Unifying Theoretical Framework: Geometric Message Passing}\label{subsec:marco_teorico_unificador}

Resuming the synthesis of \autoref{ch:gdl} (Section~\ref{subsec:message_passing_unificador}) and connecting with the deep analysis of geometric CNNs, we establish the following unifying framework that subsumes both convolutions and attention under a common mathematical abstraction.

\begin{definition}[Generalized Geometric Message Passing]\label{def:message_passing_geometrico_generalizado}
Let $\mathcal{D}$ be a geometric domain (Euclidean grid, graph, manifold, homogeneous space), $G$ a group of symmetries, and $\mathcal{F}(\mathcal{D})$ the space of signals over $\mathcal{D}$. A Geometric Message Passing operator is a family of transformations:
\begin{equation}
    \mathcal{T}^{(G)}: \mathcal{F}(\mathcal{D}) \to \mathcal{F}(\mathcal{D})
    \label{eq:message_passing_geometrico_generalizado}
\end{equation}
that satisfies:
\begin{enumerate}
    \item \textbf{G-Equivariance}: $\mathcal{T}^{(G)}(g \cdot f) = g \cdot \mathcal{T}^{(G)}(f)$ for all $g \in G$
    \item \textbf{Geometric Locality}: Support determined by the geometry of $\mathcal{D}$
    \item \textbf{Compositionality}: $\mathcal{T}^{(G)}_2 \circ \mathcal{T}^{(G)}_1$ preserves geometric properties
\end{enumerate}
\end{definition}

\begin{theorem}[Fundamental CNN-Attention Correspondence]\label{thm:correspondencia_cnn_attention}
There exists a bijective correspondence between geometric convolutions and geometric attention mechanisms through Geometric Message Passing:
\begin{align}
    \text{\textbf{Geometric CNN}} &\leftrightarrow \text{Message Passing with fixed weights} \\
    \text{\textbf{Geometric Attention}} &\leftrightarrow \text{Message Passing with adaptive weights} \\
    \text{\textbf{Hybrid}} &\leftrightarrow \text{Message Passing with partial structure}
\end{align}
This correspondence preserves the properties of equivariance, expressive power, and computational complexity.
\end{theorem}

\begin{proof}[Proof sketch]
The proof is constructed by establishing formal equivalence at three levels:

\textbf{Algebraic Level}: Both operations are expressed as generalized convolutions in the group algebra $\mathbb{C}[G]$. For CNNs, the kernel is fixed: $k_g$; for attention, it is dynamic: $\alpha_{ij}(\mathbf{h}_i, \mathbf{h}_j)$.

\textbf{Spectral Level}: In the Fourier basis of the geometric domain, both operations correspond to multiplication by transfer functions. CNNs use fixed spectral filters $\hat{k}(\lambda)$; attention uses adaptive filters $\hat{\alpha}(\lambda; \mathbf{h})$.

\textbf{Geometric Level}: The neighborhood structure determines the support of integration. CNNs integrate over fixed neighborhoods; attention integrates with weights that depend on the content and respect the underlying geometry.
\end{proof}

\subsection{Table of Theoretical Equivalences}\label{subsec:tabla_equivalencias}

The following systematic analysis establishes the direct correspondences between fundamental concepts developed in both chapters:

\begin{table}[H]
\centering
\caption{Fundamental Mathematical Equivalences: CNNs $\leftrightarrow$ Attention}
\label{tab:cnn_attention_equivalencias_fundamentales}
\begin{tabular}{|p{4cm}|p{4cm}|p{6cm}|}
\hline
\textbf{Geometric CNN} & \textbf{Geometric Attention} & \textbf{Unifying Principle} \\
\hline\hline
\textit{Group Convolution} \\
(\autoref{sec:group_convolutions}) & \textit{G-Equivariant Attention} \\
(Section~\ref{subsec:equivariant_attention}) & Operations that commute with the group action $G$ \\
\hline
\textit{Spectral Filters} \\
$\hat{f}(\lambda_i) = g(\lambda_i) \hat{x}(\lambda_i)$ & \textit{Spectral Attention} \\
$A_{ij} = \text{softmax}(f(\lambda_i, \lambda_j))$ & Multiplication in the Fourier domain of the Laplacian \\
\hline
\textit{Gauge Convolution} \\
(\autoref{sec:gauge_cnns}) & \textit{Gauge-Equivariant Attention} & Respect for local coordinate transformations \\
\hline
\textit{Steerable CNNs} \\
$f \mapsto \rho(g) f$ & \textit{Steerable Attention} \\
$A \mapsto \rho(g) A \rho(g)^{-1}$ & Coherent action on vector fields \\
\hline
\textit{Spherical CNNs} \\
Spherical harmonics $Y_\ell^m$ & \textit{Spherical Attention} \\
Inner products in $L^2(S^2)$ & Harmonic analysis on the sphere \\
\hline
\textit{Manifold CNNs} \\
Parallel transport & \textit{Riemannian Attention} \\
Riemannian connections & Intrinsic geometry of manifolds \\
\hline
\textit{ChebNet} \\
Chebyshev polynomials & \textit{Polynomial Attention} \\
Polynomial approximation & Efficient spectral localization \\
\hline
\end{tabular}
\end{table}

\begin{remark}[Hierarchy of Architectural Generalization]\label{rem:jerarquia_generalizacion}
The architectures of Geometric Deep Learning form a natural hierarchy of generalization:
\begin{equation}
\text{ANN} \subset \text{CNN} \subset \text{G-CNN} \subset \text{Manifold-CNN} \stackrel{\text{dual}}{\leftrightarrow} \text{Attention} \subset \text{G-Attention} \subset \text{Manifold-Attention}
\end{equation}
where each inclusion preserves fundamental geometric properties while increasing expressivity.
\end{remark}

\subsection{Deep Comparative Analysis: Rigidity vs. Adaptability}\label{subsec:rigidez_adaptabilidad}

One of the most fundamental distinctions between convolutional and attention paradigms lies in the balance between imposed structure (rigidity) and learnable flexibility (adaptability). This tension can be formalized mathematically:

\begin{definition}[Rigidity-Adaptability Spectrum]\label{def:espectro_rigidez_adaptabilidad}
For a neural architecture $\mathcal{A}$ operating over a domain $\mathcal{D}$, we define:
\begin{align}
    \text{Rigidity}(\mathcal{A}) &= \frac{\dim(\text{Imposed symmetries})}{\dim(\text{Parameter space})} \\
    \text{Adaptability}(\mathcal{A}) &= \frac{\dim(\text{Content-dependent weights})}{\dim(\text{Parameter space})}
\end{align}
with the fundamental constraint: $\text{Rigidity}(\mathcal{A}) + \text{Adaptability}(\mathcal{A}) \leq 1$.
\end{definition}

\begin{remark}[Fundamental Rigidity-Adaptability Trade-off]\label{rem:tradeoff_rigidez_adaptabilidad}
For Geometric Deep Learning architectures there is a fundamental trade-off:
\begin{equation}
    \text{Theoretical Guarantees} \propto \text{Rigidity}(\mathcal{A})
\end{equation}
\begin{equation}
    \text{Expressive Capacity} \propto \text{Adaptability}(\mathcal{A})
\end{equation}
The optimal architectures maximize $\text{Rigidity} \times \text{Adaptability}$ subject to the constraints of the problem.
\end{remark}

\begin{example}[Analysis of the Spectrum]\label{ex:analisis_espectro}
Consider the progression:
\begin{align}
    \text{Classical CNN}: &\quad \text{Rigidity} = 0.9, \text{Adaptability} = 0.1 \\
    \text{Group CNN}: &\quad \text{Rigidity} = 0.7, \text{Adaptability} = 0.3 \\
    \text{Graph Attention}: &\quad \text{Rigidity} = 0.3, \text{Adaptability} = 0.7 \\
    \text{Full Transformer}: &\quad \text{Rigidity} = 0.1, \text{Adaptability} = 0.9
\end{align}
This explains why CNNs provide better generalization guarantees with less data, while Transformers require massive data but achieve greater expressivity.
\end{example}

\subsection{Convergence of Paradigms: Mathematical Analysis}\label{subsec:convergencia_paradigmas}

Recent developments (2023-2024) suggest a conceptual convergence between convolutional and attention paradigms. This convergence can be formalized mathematically:

\begin{conjecture}[CNN-Attention Convergence]\label{conj:convergencia_cnn_attention}
For sequences of architectures $\{\mathcal{C}_n\}$ (CNNs with increasing complexity) and $\{\mathcal{A}_n\}$ (Attention with increasing structure), there exists a common limit:
\begin{equation}
    \lim_{n \to \infty} d(\mathcal{C}_n, \mathcal{A}_n) = 0
\end{equation}
where $d$ is a metric on the space of functionals over $\mathcal{F}(\mathcal{D})$, and the convergence occurs in the topology of compact operators.
\end{conjecture}

\textbf{Heuristic evidence.} The conjecture is supported by three observations:

\textbf{1. Spectral approximation}: Both families of operators can be approximated arbitrarily well by finite-dimensional spectral operators.

\textbf{2. Density in function spaces}: The set of geometric convolutions is dense in the space of equivariant operators, and attention mechanisms with geometric constraints converge to the same limit set.

\textbf{3. Adaptive regularization}: Modern regularization techniques (geometric dropout, adaptive normalization) cause both families to converge toward the same optimal bias-variance regime.

\subsubsection{Convergence Architectures}

The architectures that exemplify this convergence include:

\begin{definition}[ConvNext-Attention Hybrid]\label{def:convnext_attention}
A hybrid architecture that combines:
\begin{itemize}
    \item \textbf{Convolutional kernel}: To process local features with geometric guarantees
    \item \textbf{Attention heads}: To capture adaptive global dependencies
    \item \textbf{Dynamic switching}: A learnable mechanism that decides when to use each component
\end{itemize}
\end{definition}

\begin{definition}[Geometric Retentive Networks]\label{def:retentive_networks_geometricas}
Extension of the retentive mechanism to geometric domains:
\begin{equation}
    \text{Retention}(Q, K, V) = Q \odot \sum_{j} \text{decay}(d_\mathcal{M}(i,j)) \cdot K_j \odot V_j
\end{equation}
where $d_\mathcal{M}$ is the geodesic distance in the geometric domain $\mathcal{M}$.
\end{definition}

\subsection{Future Theoretical Developments}\label{subsec:desarrollos_teoricos_futuros}

\subsubsection{Unified Expressivity Limits}

\begin{conjecture}[Universal Limits of Geometric Approximation]\label{conj:limites_universales}
For every continuous function $f: \mathcal{M} \to \mathbb{R}$ over a compact manifold $\mathcal{M}$, and every $\epsilon > 0$, there exist CNN and Attention architectures of complexities $C_{CNN}(\epsilon)$ and $C_{Att}(\epsilon)$ such that:
\begin{equation}
    \frac{C_{Att}(\epsilon)}{C_{CNN}(\epsilon)} \to 1 \quad \text{as } \epsilon \to 0
\end{equation}
This conjecture suggests that the difference in architectural complexity vanishes in the limit of perfect approximation.
\end{conjecture}

\subsubsection{Geometric Information Theory}

\begin{definition}[Geometric Fisher Entropy]\label{def:entropia_geometrica_fisher}
For a distribution $p$ over a Riemannian manifold $(\mathcal{M}, g)$, the geometric entropy is defined as:
\begin{equation}
    H_g(p) = -\int_\mathcal{M} p(x) \log p(x) \sqrt{\det g(x)} \, dx
\end{equation}
This quantity captures both the statistical uncertainty and the underlying geometric complexity.
\end{definition}

\begin{conjecture}[Principle of Minimal Geometric Information]\label{conj:principio_informacion_minima}
The optimal GDL architectures minimize the geometric information necessary to represent the target function:
\begin{equation}
    \mathcal{A}^* = \arg\min_\mathcal{A} \left[ \mathcal{L}(\mathcal{A}) + \lambda H_g(\text{representations}_\mathcal{A}) \right]
\end{equation}
where $\mathcal{L}$ is the empirical loss and $\lambda$ controls the information-precision trade-off.
\end{conjecture}

\subsubsection{Geometric Meta-Learning}

\begin{definition}[Geometric Meta-Learner]\label{def:meta_learner_geometrico}
A system that automatically learns the relevant symmetries of a domain:
\begin{equation}
    \text{Meta-GDL}: \mathcal{D} \mapsto G_{\text{optimal}} \mapsto \mathcal{A}_{G_{\text{optimal}}}
\end{equation}
where $\mathcal{D}$ is the data, $G_{\text{optimal}}$ is the discovered group of symmetries, and $\mathcal{A}_{G_{\text{optimal}}}$ is the resulting architecture.
\end{definition}

\subsection{Final Reflection: Toward a Unified Theory}\label{subsec:reflexion_final}

The analyses presented in this chapter and in \autoref{ch:cnn} converge toward a fundamental conclusion: the distinction between "convolutional" and "attention" is increasingly less relevant from a deep theoretical perspective. Both paradigms represent different points on a continuum of geometric operators that process structured information.

The future of Geometric Deep Learning does not lie in the competition between paradigms, but in their intelligent synthesis. The most powerful architectures will combine the stability and efficiency of geometric convolutions with the flexibility and expressivity of adaptive attention, guided by deep mathematical principles rather than ad hoc heuristics.

\begin{remark}[Principle of Architectural Synthesis]\label{rem:principio_sintesis}
For every geometric problem with known structure $\mathcal{S}$ and required adaptability $\mathcal{A}$, there exists an optimal hybrid architecture:
\begin{equation}
    \mathcal{H}^* = \arg\min_{\mathcal{H}} \left[ \mathcal{L}(\mathcal{H}) + \lambda_1 d(\mathcal{H}, \mathcal{S}) + \lambda_2 (1 - \text{Adaptability}(\mathcal{H})) \right]
\end{equation}
where $d(\mathcal{H}, \mathcal{S})$ measures the distance to the known structure and the parameters $\lambda_i$ balance the constraints.
\end{remark}

This synthesis represents more than a technical convergence: it symbolizes the maturation of the field toward unifying principles that transcend historical divisions. As we will demonstrate in the next experimental chapter, this theoretical unification translates into measurable practical advantages and opens new possibilities for addressing the great scientific and technological challenges of the 21st century.

Geometry, as has been the case throughout the history of mathematics and physics, once again provides the unifying language needed to understand and exploit the apparent complexity of modern machine learning systems. In this geometric synthesis lies not only the understanding of the present, but also the conceptual tools to build the future of deep learning.
